\documentclass[11pt,oneside,openany]{book}
\usepackage[T1]{fontenc}
\usepackage{lmodern}
\usepackage{amsmath,amssymb,amsthm,mathtools}
\usepackage[letterpaper,margin=1.12in]{geometry}
\usepackage{microtype,enumitem,booktabs,xcolor}
\usepackage[unicode,colorlinks=true,linkcolor=blue!55!black,citecolor=teal!58!black,urlcolor=blue!55!black]{hyperref}
\usepackage{bookmark,cleveref}
\newtheorem{theorem}{Theorem}[chapter]
\newtheorem{lemma}[theorem]{Lemma}
\newtheorem{proposition}[theorem]{Proposition}
\newtheorem{corollary}[theorem]{Corollary}
\theoremstyle{definition}
\newtheorem{definition}[theorem]{Definition}
\newtheorem{example}[theorem]{Example}
\theoremstyle{remark}
\newtheorem{remark}[theorem]{Remark}
\allowdisplaybreaks[1]
\setlength{\emergencystretch}{2em}
\setcounter{tocdepth}{1}
\hypersetup{pdftitle={Rings, Modules, Representations, and Homological Algebra},
 pdfauthor={MathGenome},pdfsubject={P1-03 Unified Mathematical Notes}}

\begin{document}
\frontmatter
\hypersetup{pageanchor=false}
\begin{titlepage}
\centering
\vspace*{0.17\textheight}
{\large\scshape MathGenome\par}
\vspace{2em}
{\Huge\bfseries Rings, Modules, Representations, and Homological Algebra\par}
\vspace{1.8em}
{\Large Unified mathematical notes\par}
\vfill
{\large P1-03\par}
\vspace{0.5em}
{7 November 2026\par}
\end{titlepage}
\hypersetup{pageanchor=true}

\chapter*{Scope and conventions}
\addcontentsline{toc}{chapter}{Scope and conventions}

The first three chapters establish a common ring and module language. One branch then treats commutative rings, from localization and prime ideals to dimension and completion. A second develops ordinary finite-group representations, including induction and its double-coset decomposition. The third constructs the homological tools for modules, including coefficient theorems. The last chapter uses the two algebraic branches together in a local flatness criterion.

Rings are associative and unital and homomorphisms preserve the identity. The zero ring is allowed unless a field, a domain, a proper maximal ideal, or a nonzero ring is specified. Modules are unital and are left modules by default. A right module is a left module over the opposite ring. Over a possibly noncommutative ring \(R\), the tensor symbol \(A_R\otimes_R {}_RB\) always places the right module first. A second scalar action on a bimodule is written explicitly when it matters. Rings in Chapters 4--7 and 12 are commutative. In Chapters 8--9 the base is a field unless a coefficient ring is specified, and the ordinary complex character formulas use the Hermitian inner product linear in the first variable. Homological differentials lower degree, cohomological differentials raise it, and ordinary resolutions start in degree zero. The tensor total differential on degrees \(p,q\) is \(d_C\otimes 1+(-1)^p 1\otimes d_D\).

We use elementary group, ring, and field algebra, linear algebra, the axiom of choice in maximality and basis arguments, and the usual set and category language for universal properties. In particular, existence of maximal ideals in nonzero unital rings, elementary properties of polynomial rings over fields and UFDs, algebraic and separable field extensions, and basic linear algebra of tensor products over a field can be read in \cite{dummit,d07,d05}. Specific external results used in later proofs are identified where applied. The accompanying result map connects every writing-checklist identifier to its statement, proof context, and precise textbook location.

{\small\tableofcontents}
\mainmatter
\chapter[Rings, modules, and exact sequences]{Rings, modules, and exact sequences}
\label{p103:ch01}

\section{Quotients, scalars, and maps}

\begin{proposition}[The quotient of a ring]\label{p103:ALG-03-C01-R01}
Let \(I\) be a two-sided ideal of a unital ring \(R\). The quotient
additive group \(R/I\) has a unique ring structure for which the
projection \(\pi\) is a unital homomorphism. For a unital ring map
\(f:R\to S\), there is a unique unital map \(\bar f:R/I\to S\) with
\(\bar f\pi=f\) if and only if \(I\subseteq\ker f\).
\end{proposition}
\begin{proof}
If \(r-r',s-s'\in I\), then
\(rs-r's'=(r-r')s+r'(s-s')\in I\). This proves well-definedness of
multiplication, and the ring laws pass to classes by surjectivity of
\(\pi\). Necessity follows by applying \(f\) to \(I\). Sufficiency and
uniqueness follow from the formula \(\bar f(r+I)=f(r)\).
The case \(I=R\) gives the zero ring. Factorization then forces
\(S=0\) because \(f\) is unital.
\end{proof}

\begin{proposition}[Submodules and exactness]\label{p103:ALG-03-C01-R02}
For any \(R\)-linear \(f:M\to N\), the map \(M/\ker f\to
\operatorname{im}f\) induced by \(f\) is an isomorphism.
The inclusion of \(\ker f\) has the kernel universal property, and
\(N/\operatorname{im}f\) has the cokernel universal property.
A sequence \(L\xrightarrow{u}M\xrightarrow{v}N\) is exact at \(M\)
when \(\operatorname{im}u=\ker v\). Thus a zero term immediately to
the left of a map tests injectivity, and a zero term immediately to
its right tests surjectivity.
\end{proposition}
\begin{proof}
The action on cosets of a submodule \(U\) is \(r(m+U)=rm+U\).
A map annihilating \(U\) descends uniquely by evaluation on
representatives, whereas a map whose composite with \(f\) vanishes
has image in \(\ker f\) and factors uniquely through that submodule.
The induced map to the image is onto by its definition and is
injective by the defining relation of the quotient by \(\ker f\).
All assertions hold on either side, using \(R^{\mathrm{op}}\) to
replace a right action by a left action.
\end{proof}

\begin{definition}[Algebras and restriction of scalars]\label{p103:ALG-03-C01-R03}
For a commutative ring \(k\), a unital \(k\)-algebra is a ring \(A\)
with a unital map \(k\to Z(A)\). Its \(k\)-action is multiplication
by the image. A \(k\)-algebra map is a unital ring map compatible
with these structure maps. Any ring map \(u:R\to S\), without a
centrality requirement, restricts a left \(S\)-module to a left
\(R\)-module by \(rx=u(r)x\).
\end{definition}
The formulas satisfy the scalar laws by the ring-map laws. In the
algebra case, centrality makes multiplication \(k\)-bilinear.
Conversely, a compatible unital multiplication on a \(k\)-module
has \(a1_A\) central for every \(a\in k\), since
\((a1_A)x=a(1_Ax)=a(x1_A)=x(a1_A)\). This recovers the structure
map and the same convention for morphisms. See \cite[Sections 7.2,
7.3, and 10.1--10.2]{dummit} for the underlying ring and module
constructions.

\begin{proposition}[Free modules and split sequences]\label{p103:ALG-03-C01-S01}
The free left module on a set \(X\) is \(R^{(X)}\), the functions
of finite support, with its standard basis. A function from \(X\)
into a module extends uniquely to an \(R\)-linear map.
For a short exact sequence \(0\to A\xrightarrow{i} B
\xrightarrow{p} C\to0\), the following conditions are equivalent.
\(p\) has an \(R\)-linear section, \(i\) has an \(R\)-linear
retraction, and the sequence is isomorphic over \(A,C\) to the
direct-sum sequence. A map of short exact sequences with both end
maps isomorphisms has its middle map an isomorphism.
\end{proposition}
\begin{proof}
The unique extension is the sum of the finitely many nonzero
coefficient multiples. Arbitrary products, in contrast, impose no
finite-support condition. A section \(s\) gives the inverse maps
\((a,c)\mapsto i(a)+s(c)\) and
\(b\mapsto (i^{-1}(b-sp(b)),p(b))\). A retraction \(r\) splits
\(B=i(A)\oplus\ker r\), and \(p\) induces a bijection from \(\ker r\)
to \(C\), giving a section. The standard sum supplies both maps.
For the short five lemma, a middle element mapping to zero has
zero image in the right term by injectivity, and then vanishes by
injectivity on the left. For surjectivity, lift the right image of
a target element using the right isomorphism and the top quotient.
The remaining difference lies in the lower left image, so is lifted
using surjectivity on the left.
\end{proof}

\section{Diagram lemmas}

\begin{theorem}[Snake lemma and the four and five lemmas]
\label{p103:ALG-10-C01-R04}
In a commuting diagram of two sequences \(A\xrightarrow{i}B
\xrightarrow{p}C\) and \(A'\xrightarrow{i'}B'\xrightarrow{p'}C'\),
assume exactness in the middle of each, \(p\) onto, and \(i'\)
injective. Write \(a,b,c\) for the three vertical maps. There is a
natural exact sequence
\[
 \ker a\longrightarrow\ker b\longrightarrow\ker c
 \xrightarrow{\delta}\operatorname{coker}a
 \longrightarrow\operatorname{coker}b
 \longrightarrow\operatorname{coker}c.
\]
If \(i\) is injective the first arrow is injective. If \(p'\) is
onto the last arrow is onto.

In a commuting five-term diagram exact at the three interior
positions in each row, with vertical maps \(h_1,\ldots,h_5\), the
map \(h_3\) is injective if \(h_2,h_4\) are injective and \(h_1\)
is onto. It is onto if \(h_2,h_4\) are onto and \(h_5\) is
injective. In particular it is an isomorphism if the four outer
vertical maps are isomorphisms.
\end{theorem}
\begin{proof}
For \(z\in\ker c\), take \(y\in B\) with \(p(y)=z\). Then
\(p'b(y)=0\), so there is a unique \(x'\in A'\) with
\(i'(x')=b(y)\). Set \(\delta z=[x']\) modulo \(a(A)\).
Changing \(y\) adds \(i(x)\), hence adds \(a(x)\) to \(x'\).
This also proves additivity, linearity, and naturality for maps of
such diagrams. The induced arrows before and after \(\delta\)
are the restrictions and quotients of the horizontal maps.
An element of \(\ker b\) with zero \(p\)-image is \(i(x)\), and
\(i'a(x)=0\) implies \(x\in\ker a\). An element \(z\) in
\(\ker\delta\) admits a lift \(y\) with \(b(y)=i'a(x)\).
Replace \(y\) by \(y-i(x)\). If \(i'(x')=b(y)\) then
\(p(y)\in\ker c\) and \([x']=\delta p(y)\), proving exactness
at \(\operatorname{coker}a\). Finally, if \(p'(y')=c(z)\),
lift \(z\) to \(y\). Then \(y'-b(y)\in i'(A')\). These are all
the interior positions. The endpoint assertions follow from the
indicated injectivity or surjectivity of the horizontal maps.

For the five-term assertion let \(d_j,d_j'\) denote arrows from
position \(j\) to \(j+1\). If \(h_3(x)=0\), injectivity of \(h_4\)
gives \(d_3x=0\), so \(x=d_2y\). Then \(h_2(y)=d_1'y'\) by
exactness. Lift \(y'\) under \(h_1\). Injectivity of \(h_2\) gives
\(y=d_1(y_0)\), hence \(x=0\). For a target \(x'\) at position
3, lift \(d_3'x'\) to \(z\) at position 4 using \(h_4\).
Then \(h_5(d_4z)=0\), hence \(z=d_3 x\). The difference
\(x'-h_3x=d_2'y'\) is lifted using \(h_2\) and added to \(x\).
\end{proof}
These arguments apply on both sides of an arbitrary ring. We will
use the connecting-map construction for complexes in
Chapter~\ref{p103:ch10}. Compare \cite[Lemma 1.3.2]{weibel}
and \cite[Proposition 2.72]{rotman}.

\section{Balanced tensors and exact functors}

\begin{theorem}[Balanced tensor and the tensor--Hom adjunction]
\label{p103:ALG-03-C02-S01}
For a right module \(A_R\) and a left module \({}_RB\), the
abelian group \(A\otimes_R B\) represents biadditive pairings
with \((ar,b)\) and \((a,rb)\) identified. With bimodules,
the unused scalar actions pass to the tensor. Associativity,
unit and direct-sum isomorphisms are natural.

If \(B\) is an \(R\)-\(S\) bimodule and \(C\) a right \(S\)-module,
there is a natural bijection of abelian groups
\[
 \operatorname{Hom}_S(A\otimes_R B,C)
 \simeq \operatorname{Hom}_R(A,\operatorname{Hom}_S(B,C)),\qquad
 (g\cdot r)(b)=g(rb).
\]
For a ring map \(R\to S\), extension of left scalars
\(S\otimes_R-\) is left adjoint to restriction.
\end{theorem}
\begin{proof}
Quotient the free abelian group on pairs by the two additivity
relations and \([ar,b]-[a,rb]\). Their images give the asserted
universality. If \(A\) also has a commuting left action, or \(B\)
a commuting right action, acting on the relevant variable preserves
each relation. The two groupings of a triple tensor represent
triadditive maps balanced at both adjoining rings, so the maps
\((a\otimes b)\otimes c\leftrightarrow a\otimes(b\otimes c)\)
are inverse and respect the outer actions. The unit map
\(A\otimes_R R\to A\) has inverse \(a\mapsto a\otimes1\).
The other side is symmetric. For direct sums, the comparison and
its inverse are induced by the coordinate inclusions and by
expanding each finite-support vector into its coordinates.
These maps are natural because they commute on simple tensors.

Send \(f\) to \(a\mapsto(b\mapsto f(a\otimes b))\).
Balance gives right \(R\)-linearity of this map, and conversely
evaluation of such a map at a pair is balanced. The displayed
right action is associative since applying \(r\) and then \(r'\)
evaluates at \(r(r'b)\). These constructions are inverse and
commute with maps in all variables. For extension of left scalars,
an \(S\)-linear map out of \(S\otimes_R M\) is determined by
\(m\mapsto f(1\otimes m)\), and the inverse sends an \(R\)-linear
map \(u\) to \(s\otimes m\mapsto s\,u(m)\). The unit is
\(m\mapsto1\otimes m\) and the counit is multiplication on the
restricted module. The two triangle laws follow on generators.
\end{proof}

\begin{proposition}[Exactness of Hom and tensor]
\label{p103:ALG-03-C02-R01}
Both \(\operatorname{Hom}_R(P,-)\) and \(\operatorname{Hom}_R(-,Q)\)
are left exact, with arrows reversed for the latter.
The balanced tensor is right exact in each variable. With outer
bimodule actions these are statements in the corresponding module
categories as well.
\end{proposition}
\begin{proof}
For covariant Hom, a map whose composite with the quotient map
vanishes factors uniquely through its kernel. For contravariant
Hom, a map vanishing on the included submodule factors uniquely
through the quotient. These factorizations and uniqueness give the
two left-exact sequences. For tensor, a balanced map on \(A\times B\)
that vanishes on \(U\times B\) for a submodule \(U\subseteq A\)
descends uniquely to \(A/U\times B\) and remains balanced.
Thus \((A/U)\otimes_R B\) is the cokernel of
\(U\otimes_R B\to A\otimes_R B\). Replacing \(U\) by the image of
an arbitrary map gives right exactness. Descent through \(B/V\)
is the same argument. The factorizations preserve commuting
scalar actions by their formulas.
\end{proof}

For \(0\to\mathbb Z\xrightarrow{2}\mathbb Z\to\mathbb Z/2\to0\),
tensoring with \(\mathbb Z/2\) makes the first arrow zero on a
nonzero group. In the covariant Hom sequence for
\(P=\mathbb Z/2\), the identity map to \(\mathbb Z/2\) cannot be
lifted to \(\mathbb Z\). In the contravariant Hom sequence with
\(Q=\mathbb Z\), restriction along multiplication by 2 on
\(\mathbb Z\) is multiplication by 2 on Hom, hence is not onto.
The general balanced and diagrammatic arguments are elaborations
of \cite[Sections 10.4--10.5]{dummit} and
\cite[Theorems 2.75--2.76]{rotman}.
\chapter[Projective, injective, flat, and finite modules]
{Projective, injective, flat, and finite modules}
\label{p103:ch02}

\section{Lifting, extension, and presentations}

\begin{theorem}[Projective modules]\label{p103:ALG-03-C02-R02}
For a left module \(P\), the following are equivalent.
\begin{enumerate}[label=(\roman*)]
\item Maps out of \(P\) lift along all surjections.
\item Every short exact sequence ending in \(P\) splits.
\item \(P\) is a direct summand of a free module.
\item \(\operatorname{Hom}_R(P,-)\) is exact.
\end{enumerate}
The corresponding assertions hold for right modules.
\end{theorem}
\begin{proof}
Lifting the identity gives (ii). A free module on the underlying
set of \(P\) maps onto it, giving (iii) from (ii). Maps from a free
module lift by choosing preimages of basis elements, and a direct
summand inherits this property by extending a map by zero on its
complement. Finally, \(\operatorname{Hom}_R(P,-)\) is left exact,
and surjectivity on Hom for every quotient map is precisely (i).
\end{proof}

\begin{theorem}[Injective modules and Baer's criterion]
\label{p103:ALG-03-C02-R03}
A left module \(E\) is injective if every map to \(E\) extends along
every inclusion of left modules. It suffices to test only the
inclusions of left ideals \(I\hookrightarrow R\).
Equivalently, \(\operatorname{Hom}_R(-,E)\) is exact, or every short
exact sequence starting in \(E\) splits. The right version uses
right ideals.
\end{theorem}
\begin{proof}
Assume the ideal test. Extend \(f:U\to E\) on a submodule
\(U\subseteq M\) to a maximal partial extension \(g:V\to E\),
using the union of every chain of extensions. If \(x\notin V\),
then \(I=\{r\in R:rx\in V\}\) is a left ideal and \(r\mapsto
g(rx)\) is left linear. An extension to \(R\) has form
\(r\mapsto re\). Now \(g(v)+re\) defines a linear value on \(v+rx\).
The difference between two choices has \(r-r'\in I\), where the
chosen value agrees with \(g\). This contradicts maximality.
The reverse ideal test is immediate. Surjectivity of restriction
on Hom is exactly the extension condition, and Hom is left exact.
Injectivity splits any inclusion of \(E\) by extending its identity.
Conversely, given \(f:U\to E\) and \(U\subseteq M\), form
\((E\oplus M)/\{(f(u),-u):u\in U\}\). The induced inclusion of
\(E\) is injective. A retraction of it, applied to the image of
\(M\), extends \(f\).
\end{proof}
See \cite[Sections 2A--2B and Theorem 3.7]{lam}.

\begin{theorem}[Finite presentations, Schanuel, and dual bases]
\label{p103:ALG-03-C02-S02}
A module \(M\) is finitely presented if and only if it has a
surjection from a finite-rank free module with finitely generated
kernel. Then every such surjection has finitely generated kernel.
For projective presentations \(0\to K_i\to P_i\to M\to0\) for
\(i=1,2\), there is an isomorphism \(K_1\oplus P_2\simeq
K_2\oplus P_1\). A left module is projective if and only if it
admits elements \(p_j\) and linear maps \(\phi_j:P\to R\) such
that for each \(x\) only finitely many \(\phi_j(x)\) are nonzero
and \(x=\sum_j\phi_j(x)p_j\). For a finitely generated projective
module the index set can be taken finite. On the right, coefficients
in this last formula are on the right.
\end{theorem}
\begin{proof}
The fiber product \(P_1\times_M P_2\) projects onto \(P_1\) with
kernel \(K_2\) and onto \(P_2\) with kernel \(K_1\). Each projection
splits by projectivity, proving the comparison. A finite presentation
is precisely the cokernel of a map between finite free modules.
The kernel condition for any second finite free cover follows
from the comparison. Its kernel is a direct summand of a finitely
generated module. For a split embedding \(P\hookrightarrow
R^{(J)}\) take its coordinates as \(\phi_j\) and send the standard
basis under the retraction to \(p_j\). Conversely these formulas
exhibit a split embedding. The supports of finitely many generators
have finite union, which suffices for the whole embedding.
\end{proof}
The finite-presentation and fiber-product arguments follow
\cite[Sections 4D and 5A]{lam}.

\begin{theorem}[Enough injectives and injective envelopes]
\label{p103:ALG-03-C02-S03}
Every module embeds in an injective module and has an injective
envelope. An extension \(M\subseteq E\) is essential if every
nonzero submodule of \(E\) meets \(M\) nontrivially. For an
extension of \(M\), being injective and essential is equivalent to
being a maximal essential extension (no proper essential extension
over \(M\)), and equivalent to being an injective extension
minimal among injective extensions inside it. Any two injective
envelopes are isomorphic by an isomorphism restricting to the
identity on \(M\).
\end{theorem}
\begin{proof}
The group \(\mathbb Q/\mathbb Z\) is divisible, hence injective over
\(\mathbb Z\) by the ideal test. Any nonzero element of an abelian
group can first be detected by a nonzero character on its cyclic
subgroup, then by an extension of that character. For a left
module use \(J=\operatorname{Hom}_{\mathbb Z}(R,\mathbb Q/\mathbb Z)\)
with \((rf)(s)=f(sr)\). Maps from \(N\) to \(J\) correspond to
group homomorphisms \(N\to\mathbb Q/\mathbb Z\) by evaluation at
1, the inverse sending \(u\) to \(n\mapsto(s\mapsto u(sn))\).
Thus \(J\) is injective and the product of these maps for all such
\(u\) embeds \(N\) in a product of copies of \(J\). A product of
injectives is injective by coordinatewise extension.

In an injective ambient module choose an essential extension \(E\)
of \(M\) maximal by inclusion, using unions of chains. If \(E\)
has an essential extension \(E'\), the inclusion of \(E\) into
the ambient injective extends to \(E'\), and the kernel is zero
by essentiality. Essentiality is transitive, so the image would
enlarge \(E\) as an essential extension of \(M\). Hence \(E\)
has no proper essential extension. For any module \(E\) with
this property, embed it into an injective \(I\) and choose a
submodule \(K\subset I\) maximal with \(K\cap E=0\). The image of
\(E\) in \(I/K\) is essential: a disjoint nonzero submodule would
have preimage strictly larger than \(K\) and still disjoint from
\(E\). It follows that \(I=E\oplus K\), so \(E\) is injective.
Conversely an essential extension of an injective module cannot
be proper, since the inclusion splits. Any injective extension
of \(M\) contains an isomorphic copy of an envelope over \(M\).
Extend the inclusion of \(M\) along the envelope and use
essentiality for injectivity of the resulting map. Thus minimal
injective extensions are exactly the envelopes. The same argument
in another envelope gives a split inclusion with zero complement,
proving uniqueness up to isomorphism over \(M\).
\end{proof}
For the right-handed statements use the opposite ring.
Compare \cite[Sections 3C--3D]{lam}.

\section{Flatness and its finite tests}

For a left module \(M\), a finite relation criterion will use a
finite list \(x_i\in M\), and finitely many rows
\(\sum_i r_{ji}x_i=0\). A factorization of the relations is an
expression \(x_i=\sum_\ell s_{i\ell}y_\ell\) with finitely many
\(y_\ell\in M\) and \(\sum_i r_{ji}s_{i\ell}=0\) for all \(j,\ell\).
The order of multiplication is part of the assertion.

\begin{theorem}[Flatness, finite relations, and Govorov--Lazard]
\label{p103:ALG-03-C02-R04}
For a left module \(M\), the following are equivalent.
\begin{enumerate}[label=(\roman*)]
\item \(-\otimes_R M\) preserves injections (equivalently it is exact).
\item \(I\otimes_R M\to R\otimes_R M\) is injective for each
finitely generated right ideal \(I\).
\item Every finite system of relations admits the above factorization.
\item \(M\) is a filtered colimit of finite-rank free left modules.
\end{enumerate}
Projective modules are flat, and finitely presented flat modules
are projective. The transition maps in (iv) may be arbitrary.
All statements have right-module versions with the sides reversed.
\end{theorem}
\begin{proof}
Right exactness was proved in Chapter~\ref{p103:ch01}, so (i)
implies (ii). Under (ii), for one relation take the right ideal
generated by its \(r_i\) and the surjection \(R^n_R\to I\)
sending the \(i\)-th basis element to \(r_i\). The relation says
that \(\sum_i e_i\otimes x_i\) maps to zero in \(R\otimes M\),
hence in \(I\otimes M\). By right exactness this tensor comes from
the kernel tensored with \(M\). Write that lift as a finite sum of
simple tensors, giving the desired factorization. For several
rows first factor the first relation, then substitute the
factorization in each remaining row. Repeating gives a product of
coefficient matrices. Rows already killed remain zero under
further right multiplication.

For (iii) to (iv) take the category of all maps \(u:R^n\to M\)
from the standard finite free left modules, with arrows linear
maps commuting with \(u\). This is small. Two objects map to
their direct sum with the sum of their specified maps to \(M\).
For parallel arrows \(\alpha,\beta\) to \(u:R^n\to M\), the
vectors \((\alpha-\beta)(e_j)\) give finitely many relations on
\(x_i=u(e_i)\). Factor them. The map \(R^n\to R^\ell\) sending
\(e_i\) to \(\sum_\ell s_{i\ell}e_\ell\) lies over \(M\) and
equalizes the arrows. Thus the category is filtered. Its colimit
maps onto \(M\) by the objects of rank one. If a representative
\(v\in R^n\) maps to zero, the two maps from the rank-one object
with zero map to \(M\), sending \(1\) respectively to \(v,0\),
can be equalized. Thus the representative already vanishes at
some later object, proving injectivity. This is the filtered
direct-limit form of the theorem.

Conversely finite free modules are flat because their tensor
functors are finite direct sums. Tensor commutes with filtered
colimits by its universal property. Such colimits of modules
preserve injections. If an element maps to zero in the target
limit, represent it at one object and use a transition at which
its target representative becomes zero, then apply injectivity
there. Hence (iv) gives (i). Direct summands of flat modules are
flat, and a free module of any rank is flat by the direct-sum
formula, proving projective implies flat.

Finally choose a finite free cover with basis images \(x_i\) and
a finite set of kernel generators giving relations indexed by \(j\).
Factor them using (iii), and lift each \(y_\ell\) to \(z_\ell\) in
the free cover. Sending \(x_i\) to \(\sum_\ell s_{i\ell}z_\ell\)
descends to the presented module, because all the chosen kernel
generators map to zero. The composite with the cover fixes the
generators, so this is a section and the module is projective.
\end{proof}
See \cite[Theorems 4.12, 4.24, 4.30, and 4.34]{lam}. The finite
presentation in the last claim is essential to this argument.
Finite generation alone does not ensure finitely many relations.

\begin{proposition}[Faithfully flat modules]\label{p103:ALG-03-C02-S04}
For a flat right module \(P\), the following are equivalent.
\(P\otimes_R-\) detects nonzero left modules, reflects exactness,
and has \(P\otimes_R R/\mathfrak m\ne0\) for every maximal left
ideal \(\mathfrak m\). In this case the exact tensor functor is
faithful. A left-module version follows by reversal of sides.
\end{proposition}
\begin{proof}
Reflection of exactness tests a module by the zero sequence
through it. Detection implies reflection since an exact functor
commutes with the homology at a position of a complex, and detects
whether a composite map has zero image. For the maximal ideal test,
a nonzero left module has a nonzero cyclic submodule. Its simple
quotient is \(R/\mathfrak m\) for some maximal left ideal.
Right exactness maps the tensor of the cyclic submodule onto
the nonzero tested tensor, and flatness injects this tensor into
the tensor of the original module. The test is necessary by
detection. If a map is nonzero, its image is nonzero and preserved
by the exact functor, which therefore sends the map to a nonzero map.
\end{proof}
See \cite[Section 4I]{lam}.

\section{Chain conditions and finite length}

\begin{proposition}[Noetherian and Artinian modules]
\label{p103:ALG-03-C03-R01}
For a module, ascending chains of submodules stabilize if and only
if every nonempty family has a maximal element, or equivalently
every submodule is finitely generated. This is the Noetherian
condition. For the Artinian condition, descending chains stabilize
if and only if every nonempty family has a minimal element.
Each property is inherited by submodules and quotients and is
preserved by extensions.
\end{proposition}
\begin{proof}
A family with no maximal element yields an infinite strictly
ascending chain, and the set of members of a strictly ascending
chain has no maximal element. To prove the finite-generation
criterion, take a maximal finitely generated submodule of any
given submodule, adding one generator if it is proper. In the
opposite direction, the union of an ascending chain is a submodule
and its finite generating set belongs to a single stage.
The minimal-element argument is the descending analogue.
Submodules inherit chains, and chains in quotients lift. In
\(0\to A\to B\to C\to0\), once the intersections with \(A\)
and images in \(C\) of a chain in \(B\) have both stabilized,
two nested terms are equal. Lift the image of an element of
the larger term in the smaller and use equality of intersections
on the difference. This treats both directions.
\end{proof}

\begin{theorem}[Jordan--H\"older and length]\label{p103:ALG-03-C03-S01}
A module has a finite composition series if and only if it is both
Noetherian and Artinian. The multiset of its simple composition
factors is independent of the series. Its length, denoted
\(\ell_R(M)\), is additive on short exact sequences of such
modules, with \(\ell_R(0)=0\).
\end{theorem}
\begin{proof}
Finite composition series imply both chain conditions by the
extension property. Conversely choose a maximal proper submodule
of any nonzero term and descend, using the descending condition
for termination. For a fixed series, intersect it with a submodule
and project it to the quotient. At each simple factor precisely
one of the induced factors is simple and the other is zero.
In \(M_i/M_{i-1}\) the image of the intersection is a submodule
and the projected factor is its quotient. In particular, a proper
submodule has an induced series with strictly fewer factors
because the quotient is nonzero. Induct on the length of a
fixed series to compare two series. If their top proper terms
\(A,B\) are equal, apply induction. If distinct, maximality gives
\(A+B=M\). With \(D=A\cap B\), the two quotients \(A/D\) and
\(B/D\) identify with \(M/B\) and \(M/A\), respectively.
Use induction inside each proper term to compare its original
series with a series through \(D\). Both lists for \(M\) now
consist of the factors of \(D\) and the same two simple factors.
The zero case starts the induction. For a short exact sequence,
concatenate a series in the submodule with preimages of a series
in the quotient, giving additivity by uniqueness.
\end{proof}
These lattice arguments apply to either side and supplement
\cite[Notes to the Reader]{lam} and \cite[Section 18.1]{dummit}.
\chapter[Radicals, semisimplicity, and Morita equivalence]
{Radicals, semisimplicity, and Morita equivalence}
\label{p103:ch03}

\section{Simple modules and the radical}

\begin{theorem}[Simple and semisimple modules]\label{p103:ALG-03-C03-R03}
A simple module is a nonzero module with no proper nonzero
submodule. A nonzero map between simple modules is an isomorphism.
The endomorphisms of a simple module form a division ring.
For a module \(M\), the following are equivalent.
\(M\) is a direct sum of simple modules, \(M\) is the sum of its
simple submodules, and every submodule of \(M\) is a direct
summand. Such modules are semisimple, and their submodules and
quotients are semisimple. These statements include the zero
module and arbitrary index sets.
\end{theorem}
\begin{proof}
For the first assertion take the kernel and image. A direct
sum is a sum. If \(M\) is a sum of simples and \(U\subseteq M\),
choose a maximal family of simple submodules whose sum \(V\)
is direct and satisfies \(U\cap V=0\). Unions of chains preserve
these conditions because relations involve only finitely many
summands. If \(U+V\ne M\), some simple submodule is not contained
in \(U+V\), and its intersection with \(U+V\) is zero. Adding
it contradicts maximality. Thus \(U\) has a complement, and
with \(U=0\) this argument also gives a direct-sum decomposition.

Conversely the complement property descends to submodules by
restricting retractions. A nonzero cyclic submodule has a maximal
proper submodule, since a union of proper submodules in a chain
cannot contain its generator. The complement in this cyclic
module is simple. Consequently every nonzero complement contains
a simple submodule. If the sum of all simple submodules were proper,
its nonzero complement would contain a simple not in that sum.
This proves the remaining implication. Quotients of sums of
simples are sums of their images, each simple or zero.
A submodule is a summand, hence isomorphic to a quotient.
\end{proof}

\begin{theorem}[Jacobson radical and Nakayama's lemma]
\label{p103:ALG-03-C03-R02}
For a unital ring \(R\), the intersection \(J(R)\) of its maximal
left ideals equals the intersection of the annihilators of simple
left modules. It is a two-sided ideal and agrees with the
corresponding right radical. A criterion is
\[
 x\in J(R)\quad\Longleftrightarrow\quad
             1-rx\text{ is a unit for every }r\in R.
\]
For a two-sided ideal \(I\subseteq J(R)\) and a finitely generated
left module \(M\), \(IM=M\) implies \(M=0\). Equivalently,
\(M=N+IM\) implies \(M=N\), so generators of \(M/IM\) lift to
generators of \(M\). The right version uses \(MI\).
\end{theorem}
\begin{proof}
If \(S\) is simple, every nonzero \(s\in S\) gives a surjection
\(R\to S,\ r\mapsto rs\) with maximal kernel. Hence an element of
the intersection kills every simple module. Conversely test on
\(1+L\) in \(R/L\) for each maximal left ideal \(L\).
Annihilators of whole modules are two-sided.

If \(x\in J(R)\) and \(R(1-rx)\) is proper, a maximal left
ideal containing it would contain \(rx\) and therefore 1.
Thus \(1-rx\) has a left inverse \(y\), and
\(y=1+yrx\) belongs to \(1+J(R)\). The same argument gives a
left inverse of \(y\). The latter must be \(1-rx\), so \(y\)
is two-sided. If \(x\) is missing from a maximal left ideal \(L\),
then \(L+Rx=R\), so \(1-rx\) belongs to \(L\) for some \(r\)
and cannot be a unit. Finally, \(1-ab\) invertible implies
\(1-ba\) invertible by the inverse \(1+b(1-ab)^{-1}a\).
This transfers the criterion to the opposite ring and proves
the equality of radicals. For the zero ring the assertions use
the empty intersection convention.

For Nakayama, take a generating set of minimal finite size.
If its last vector is \(x_n\), \(IM=M\) expresses
\(x_n=\sum_i a_i x_i\) with \(a_i\in I\).
Since \(1-a_n\) is a unit, \(x_n\) can be discarded, a
contradiction unless the size is zero. Apply this to \(M/N\)
for the equivalent form.
\end{proof}

\section{Artinian structure and matrix rings}

\begin{theorem}[Hopkins--Levitzki and the Artinian radical]
\label{p103:ALG-03-C03-S02}
If \(R\) is left Artinian, \(J=J(R)\) is nilpotent, \(R/J\)
is semisimple as a module over itself, and the left regular
module has finite length. In particular \(R\) is left Noetherian.
For a right Artinian ring, the corresponding right-module
assertions follow by taking the opposite ring.
If a ring has nilpotent \(J\) and semisimple \(R/J\), any
Noetherian left module over it has finite length.
A commutative Artinian ring is a finite product of local
Artinian rings.
\end{theorem}
\begin{proof}
Work first with nonzero \(R\). The powers of \(J\) stabilize by
the descending chain condition. If a stable \(J^n\) is nonzero
with \(n\ge1\), choose a minimal left ideal \(L\) subject to
\(J^nL\ne0\). Since \(J^{2n}L=J^nL\), minimality gives
\(J^nL=L\). Choose \(x\in L\) with \(J^nx\ne0\).
The left ideal \(Rx\subseteq L\) has \(J^n(Rx)\ne0\), so
\(L=Rx\). Nakayama applied to the cyclic module \(L\) and
\(J^n\) now contradicts \(L\ne0\).

In \(R/J\) the maximal left ideals correspond to those of \(R\),
so their intersection is zero. A minimal finite intersection of
these ideals exists by the descending chain condition, and
intersecting with any further maximal ideal leaves it fixed.
Thus it is zero. This embeds the left regular module of \(R/J\)
in a finite sum of simple modules, so it is semisimple.
Over a ring with semisimple left regular module all left modules
are semisimple. Each is a quotient of a direct sum of copies of
the regular module. The factors \(J^i/J^{i+1}\) are Artinian
semisimple \(R/J\)-modules, so their simple direct-sum
decompositions are finite (otherwise omitting successively
distinct summands gives a strictly descending chain).
The finite filtration by powers of \(J\) therefore gives finite
length. The same filtration \(J^iM\) works for a Noetherian
module over a ring with the two indicated properties.
Its semisimple factors are Noetherian, which likewise forces
finitely many summands, this time by ascending chains.

In the commutative case, \(R/J\) is a finite product of fields
by the finite intersection just constructed and the Chinese
remainder theorem. Its finitely many factors give all the maximal
ideals \(\mathfrak m_i\) of \(R\). For \(J^N=0\), the powers
\(\mathfrak m_i^N\) are pairwise comaximal. Expand
\(1=(a+b)^{2N-1}\) for \(a\in\mathfrak m_i,b\in\mathfrak m_j\)
with \(a+b=1\). Their product is \(J^N\), so their intersection
is zero, and the Chinese remainder theorem gives the product
of \(R/\mathfrak m_i^N\). Each quotient is local, since any
maximal ideal containing \(\mathfrak m_i^N\) contains
\(\mathfrak m_i\). The zero ring gives the empty product.
\end{proof}
For comparisons, see \cite[Notes to the Reader and Section 3K]{lam},
\cite[Section 18.2]{dummit}, and the commutative version in
\cite[Theorem 2.14 and Corollaries 2.16--2.19]{eisenbud}.

\begin{theorem}[Wedderburn--Artin]\label{p103:ALG-03-C03-R04}
For a nonzero ring \(R\), the following are equivalent.
The left regular module is semisimple, \(R\) is left Artinian
with \(J(R)=0\), and \(R\) is isomorphic to a finite product
\(\prod_i M_{n_i}(D_i)\) with \(n_i>0\) and \(D_i\) division
rings. The factors are unique up to permutation and
isomorphism, and each \(n_i\) and \(D_i\) is determined by its
factor. A simple Artinian ring has one factor.
\end{theorem}
\begin{proof}
The Artinian implication to semisimplicity is the preceding
theorem. A semisimple regular module is cyclic, so a direct sum
decomposition into simples is finite. The generator has finite
support. Group it as \(\bigoplus_i S_i^{n_i}\) with the \(S_i\)
pairwise nonisomorphic. Maps between distinct \(S_i\)'s vanish,
and within each block endomorphisms are matrices over the division
ring \(\operatorname{End}_R(S_i)\).
Right multiplication identifies \(\operatorname{End}_R({}_R R)\)
with \(R^{\mathrm{op}}\). Taking opposites, with transpose
identifying the opposite matrix ring with matrices over the
opposite division ring, gives the product in the statement.

Conversely, for \(M_n(D)\) the left regular module is the direct
sum of its \(n\) column modules. Each is simple: a nonzero column
can be sent to any given column by a matrix, using one nonzero
coordinate and its inverse. The corresponding left regular
module of a finite product is then semisimple and has finite
length. The radical is zero because its action on every simple
factor vanishes and the action on the regular module is faithful.
Each matrix factor is simple as a ring, since a nonzero entry
in an element of a two-sided ideal yields the elementary
matrices by multiplication on each side. Its center is a field,
as is seen by commuting with all elementary matrices and then
with matrices having arbitrary entries on the diagonal. Its
identity is therefore a primitive central idempotent of the
product, and these idempotents recover the individual factors.
Within one factor all simple modules are isomorphic to a column.
A simple module is a quotient of the regular module and receives
a nonzero map from one summand. Thus the regular composition
length recovers \(n\) and the opposite of the endomorphism
division ring of a simple module recovers \(D\).
The zero ring again corresponds to the empty product.
\end{proof}

\section{Progenerators and Morita equivalence}

\begin{theorem}[Morita equivalence from a progenerator]
\label{p103:ALG-03-C02-S05}
Work with right modules. Let \(P_R\) be a finitely generated
projective generator, so every right \(R\)-module is a quotient of
a sum of copies of \(P\). Put \(S=\operatorname{End}_R(P_R)\) with
composition as multiplication, and \(Q=\operatorname{Hom}_R(P,R)\).
Then \(P\) is an \(S\)-\(R\) bimodule and \(Q\) an \(R\)-\(S\)
bimodule. The maps
\[
 \begin{aligned}
 P\otimes_R Q&\longrightarrow S,&
 p\otimes q&\longmapsto (x\mapsto p\,q(x)),\\
 Q\otimes_S P&\longrightarrow R,&
 q\otimes p&\longmapsto q(p)
 \end{aligned}
\]
are bimodule isomorphisms. Tensoring on the right with \(Q\)
and \(P\) gives inverse equivalences of the right module categories
of \(R\) and \(S\). Conversely, an equivalence of such categories
recovers a progenerator by transporting a regular module.
In particular \(R\) and \(M_n(R)\) have equivalent module
categories for \(n>0\). Taking opposite rings gives the left
form, with the opposite endomorphism ring when using a left
progenerator.
\end{theorem}
\begin{proof}
The action \(sp=s(p)\) commutes with the right \(R\)-action,
and \((rqs)(x)=r\,q(sx)\) defines the two actions on \(Q\).
The two displayed maps respect balance by these formulas.
Write \(\theta(p,q)(x)=p\,q(x)\). A finite right dual basis
\(\sum_i p_iq_i(x)=x\) gives the inverse
\(s\mapsto\sum_i s(p_i)\otimes q_i\) to the first map.
On \(x\) the composite is \(s(x)\), and on \(p\otimes q\) the
other composite uses \(q=\sum_i q(p_i)q_i\).
Since a sum of copies of \(P\) surjects onto \(R_R\),
there are finite families \(u_j\in P,v_j\in Q\) with
\(\sum_j v_j(u_j)=1\). The inverse to the second map sends
\(r\mapsto\sum_j v_j\otimes u_jr\).
On a simple tensor the composite is the identity because
\[
 \sum_j v_j\otimes u_jq(p)
 =\sum_j v_j\theta(u_j,q)\otimes p=q\otimes p.
\]
Associativity and the tensor unit from Chapter~\ref{p103:ch01}
now give natural isomorphisms for both composites of the two
functors. Their compatibility with the two evaluations follows
from \(\theta(p,q)p'=p\,q(p')\) and
\(q\theta(p,q')=q(p)q'\). On tensors these are the two triangle
identities.

If an equivalence from right \(S\)-modules to right \(R\)-modules
takes \(S_S\) to \(P\), it preserves quotients, sums, and lifting
against surjections, hence \(P\) is a projective generator.
It also preserves the property that maps out of \(S_S\) into
arbitrary direct sums factor through a finite subsum, since
colimits and zero morphisms are preserved by equivalences.
Apply this property to a section of a free cover of \(P\).
The section factors through a finite free subsum, which then
surjects onto \(P\). In particular \(P\) is finitely generated.
The equivalence identifies the endomorphism rings
\(\operatorname{End}_R(P)\simeq\operatorname{End}_S(S_S)\simeq S\).
Addition in a module category is recovered from the diagonal,
codiagonal and biproduct maps. Finally \(R^n_R\) is a
progenerator with the asserted matrix endomorphism ring.
\end{proof}
The dual-basis argument and the two evaluations are the
Morita I--II construction of \cite[Sections 17B and 18B--18D]{lam}.
\chapter{Localization, spectra, and integrality}
\label{p103:ch04}

The rings in Chapters~\ref{p103:ch04}--\ref{p103:ch07} are commutative and have an identity. Ring maps preserve the identity. The zero ring is allowed, and prime and maximal ideals are proper. In particular its spectrum is empty. A localization may be the zero ring, even when the original ring is nonzero. The treatment follows \cite[Chapters 2 and 4 and Section 13.1]{eisenbud} and \cite[Sections 15.4--15.5]{dummit}.

\section{Fractions, modules, and local detection}

\begin{theorem}[The ring of fractions]\label{p103:ALG-05-C01-R01}
Let \(S\) be a multiplicatively closed subset of \(R\) containing \(1\). There is a ring \(S^{-1}R\) and a map \(j:R\to S^{-1}R\) sending \(S\) to units. Every ring map \(g:R\to B\) with \(g(S)\subset B^\times\) has a unique factorization through \(j\). The fractions satisfy
\[
 \frac{a}{s}=\frac{b}{t}\quad\Longleftrightarrow\quad
 \text{there is \(u\in S\) such that \(u(ta-sb)=0\)}.
\]
If \(0\in S\), the localization is the zero ring.
\end{theorem}
\begin{proof}
On \(R\times S\) impose the displayed relation. For transitivity, if \(u(ta-sb)=0\) and \(v(wb-tc)=0\), then \(uvt(wa-sc)=0\). Reflexivity and symmetry are immediate. Define
\[
 \frac{a}{s}+\frac{b}{t}=\frac{ta+sb}{st},
 \qquad
 \frac{a}{s}\frac{b}{t}=\frac{ab}{st}.
\]
To check independence of representatives, the differences obtained after taking the new common denominator are sums of multiples of the original two differences \(s'a-sa'\) and \(t'b-tb'\). For multiplication the difference is
\(s't'ab-sta'b'=t'b(s'a-sa')+sa'(t'b-tb')\).
The product of the two original witnesses kills these differences. The ring laws follow by writing the expressions over a common denominator. The class \(1/s\) inverts \(s/1\). The proposed factorization sends \(a/s\) to \(g(a)g(s)^{-1}\). Multiplying the equivalence equation by the inverses of the denominator images and of \(g(u)\) makes this well defined. Every factorization has this value. When \(0\in S\), zero witnesses equality of every two fractions. Conversely \(1/1=0/1\) forces an element of \(S\) to be zero.
\end{proof}

\begin{theorem}[Module localization and local detection]\label{p103:ALG-05-C01-S01}
For any \(R\)-module \(M\), localization gives \(S^{-1}M\cong S^{-1}R\otimes_R M\). It preserves exact sequences and finite intersections of submodules. A map of modules is injective, surjective or bijective exactly when the corresponding property holds at every maximal ideal of \(R\). Exactness of sequences can likewise be detected at those localizations.

If \(M\) is finitely presented, \(N\) is arbitrary, and \(B\) is a flat \(R\)-algebra, then the natural map
\[
 B\otimes_R\operatorname{Hom}_R(M,N)
 \longrightarrow\operatorname{Hom}_B(B\otimes_RM,B\otimes_RN)
\]
is an isomorphism. This applies to \(B=S^{-1}R\).
\end{theorem}
\begin{proof}
Use pairs \((m,s)\) and the same witness relation as for ring fractions, with \(ta-sb\) replaced by \(tm-sn\). In particular \(m/s=0\) exactly when some member of \(S\) kills \(m\). The inverse tensor maps are
\((a/s)\otimes m\mapsto am/s\) and \(m/s\mapsto (1/s)\otimes m\).
The first map respects the balanced relation. For the second, from \(u(tm-sn)=0\) obtain
\((1/(ust))\otimes u(tm-sn)=0\), which proves that the two representatives have the same image. Every pure tensor is recovered by these two maps. See the tensor construction in Theorem~\ref{p103:ALG-03-C02-S01}.

Suppose \(M'\xrightarrow{f}M\xrightarrow{g}M''\) is exact at \(M\). If \(g(m)/s=0\), choose \(u\in S\) with \(g(um)=0\) and then \(m'\) with \(f(m')=um\). The element \(m/s\) is the image of \(m'/(us)\). The other inclusion follows from \(gf=0\), and applying this argument to sequences with a zero endpoint deals with injectivity and surjectivity. A finite intersection is the kernel of \(M\to\bigoplus_i M/N_i\), so localization preserves it.

If \(m\ne0\) in \(M\), its annihilator is a proper ideal and is contained in a maximal ideal \(\mathfrak m\). No element of \(R\setminus\mathfrak m\) kills \(m\), so \(m/1\ne0\) in \(M_{\mathfrak m}\). Thus a module whose maximal localizations all vanish is zero, with the same conclusion for the zero ring since its unital modules are zero. Apply this to the kernel, cokernel and, for a complex, \(\ker(g)/\operatorname{im}(f)\). For a pair of maps whose localizations compose to zero, the image of their composite also vanishes locally and hence globally.

Take a finite presentation \(R^a\to R^b\to M\to0\). Then \(\operatorname{Hom}_R(M,N)\) is the kernel of \(N^b\to N^a\), since a map on generators descends precisely when it kills the relations. Flat base change preserves that kernel and both finite powers. The tensored presentation computes the kernel over \(B\) as \(\operatorname{Hom}_B(B\otimes_R M,B\otimes_R N)\). On the generators this is the stated natural map. Localization is flat because its tensor functor is the exact functor just proved.
\end{proof}

\section{The prime spectrum}

For an ideal \(I\), put \(V(I)=\{\mathfrak p:I\subseteq\mathfrak p\}\) in \(\operatorname{Spec}R\), and put \(D(f)=\{\mathfrak p:f\notin\mathfrak p\}\). Intersections of closed sets correspond to sums of ideals, and \(V(I)\cup V(J)=V(IJ)\) by primality. These identities give the Zariski topology.

\begin{theorem}[The radical and principal open sets]\label{p103:ALG-05-C01-S02}
For every ideal \(I\) of \(R\), its radical \(\sqrt I=\{a:a^n\in I\text{ for some }n\ge1\}\) is the intersection of the prime ideals containing \(I\). In particular the nilradical is \(\sqrt{(0)}\). The sets \(D(f)\) form a basis and satisfy \(D(fg)=D(f)\cap D(g)\). The localization map at \(f\) identifies \(\operatorname{Spec}(R_f)\) with \(D(f)\) as a topological space.
\end{theorem}
\begin{proof}
Any prime containing \(I\) contains its radical. If \(f\notin\sqrt I\), the powers of \(f\) avoid \(I\). By the chain union argument in Zorn's lemma there is an ideal \(J\supseteq I\) maximal among ideals avoiding those powers. If \(a,b\notin J\), both \(J+(a)\) and \(J+(b)\) meet the powers. Multiplying their two witnesses shows that \(J+(ab)\) meets the powers. Thus \(ab\notin J\), which proves primality. An empty intersection here is the unit ideal. The formula for the complement of \(V(I)\) as \(\bigcup_{f\in I}D(f)\) and primality prove the assertions about the basis. The principal localization identification is the specialization of the spectrum correspondence below.
\end{proof}

\begin{theorem}[Maps, quotient spectra, and localization spectra]\label{p103:ALG-05-C01-R02}
A ring map \(R\to A\) induces a continuous map \(\operatorname{Spec}A\to\operatorname{Spec}R\) by contraction, functorially in the reverse direction. The quotient \(R\to R/I\) identifies its spectrum with \(V(I)\). The map \(R\to S^{-1}R\) identifies the localization spectrum homeomorphically with the subspace of primes disjoint from \(S\).
\end{theorem}
\begin{proof}
The inverse image of \(D(f)\) under contraction is \(D(g(f))\). Quotient ideals correspond to ideals containing \(I\), with primality preserved in both directions, and the basic opens of the quotient are the restrictions of the basic opens of \(R\).

For a prime \(\mathfrak p\) disjoint from \(S\), the image of \(S\) in \(R/\mathfrak p\) does not contain zero. The domain \(S^{-1}(R/\mathfrak p)\) is therefore nonzero, and the kernel of \(S^{-1}R\to S^{-1}(R/\mathfrak p)\) is \(S^{-1}\mathfrak p\). A fraction is in that kernel just when \(u a\in\mathfrak p\) for some \(u\in S\), equivalently when its numerator \(a\) is in \(\mathfrak p\). This proves contraction back to \(\mathfrak p\). Conversely a prime of the localization contracts to a prime disjoint from \(S\), and a fraction belongs to the original prime exactly when its numerator over \(1\) does, since its denominator is a unit. Finally, \(D(a/s)=D(a/1)\) in the localization, so its basic opens are precisely the restrictions of the basic opens on the image subspace. If \(S\) is generated by \(f\) the image is \(D(f)\). A general such image need not be open.
\end{proof}

\section{Finite generation, integral elements, and prime chains}

An element of an \(R\)-algebra is \emph{integral} over \(R\) if it satisfies a monic polynomial with coefficients in \(R\). A finite algebra means finite as a module, and a finite-type algebra means finitely generated as an algebra.

\begin{theorem}[Finite integral generation and localization]\label{p103:ALG-05-C01-S03}
An element \(a\in A\) is integral if and only if \(R[a]\) is a finite \(R\)-module. An \(R\)-algebra \(A\) is finite if and only if it is integral and of finite type. If \(C\) is the integral closure of the image of \(R\) in \(A\), the image of \(S^{-1}C\) in \(S^{-1}A\) is the integral closure of the image of \(S^{-1}R\).
\end{theorem}
\begin{proof}
A monic relation expresses all sufficiently high powers in terms of lower powers. Conversely, for any finite \(R\)-algebra \(B\), take module generators \(b_1,\ldots,b_m\) and express multiplication by \(a\in B\) by a matrix \(U\) of coefficients in \(R\): \((aI-U)\mathbf b=0\). Multiplying by the adjugate shows that the monic polynomial \(\det(TI-U)\) evaluated at \(a\) kills each \(b_i\). It kills \(1\in B\), so gives the required equation. This argument also applies to \(B=R[a]\), and the zero algebra is covered by the monic equation \(T=0\) for its only element. A finite set of integral generators yields a module spanned by their products with each exponent below the degree of the corresponding relation. A finite module generating set is also an algebra generating set, which gives the other direction.

Monic equations survive localization. Conversely suppose \(a/s\) is integral over \(S^{-1}R\). Then \(a/1\) is integral there, by multiplying a monic equation by a suitable power of \(s\). Write such an equation for \(a/1\), of degree \(n\), with all coefficients having a common denominator \(t\in S\). Substituting \(b=ta\) gives a monic polynomial \(P(T)\in R[T]\) with \(P(b)/1=0\) in \(S^{-1}A\). Indeed the coefficient of \(b^i\) in the equation is \(t^{n-i}\) times the old coefficient, which has no denominator for \(i<n\). Some \(u\in S\) kills \(P(b)\) in \(A\). If \(n=0\), then \(P=1\), so \(u\cdot 1_A=0\) and \(ub=0\) is integral. If \(n\geq1\), the monic polynomial with coefficients \(u^{n-i}P_i\) for \(i<n\) vanishes at \(ub\), since its value is \(u^nP(b)=0\). In either case \(a/s=(ub)/(uts)\), without cancellation of a possible zero divisor in \(A\). Since localization preserves the injection \(C\subseteq A\), the localization of \(C\) is the indicated subring.
\end{proof}

\begin{theorem}[Integral elements, transitivity, and lying over]\label{p103:ALG-05-C01-R03}
The integral elements form an \(R\)-subalgebra. Integrality is transitive in a tower of rings. If \(g:R\to A\) is integral, its map on spectra has image \(V(\ker g)\).
\end{theorem}
\begin{proof}
If \(a,b\) are integral then \(R[a,b]\) is finite by the preceding theorem. All its elements are integral, in particular \(a+b\), \(-a\), and \(ab\). Scalar images themselves are integral. For transitivity, take the finitely many intermediate coefficients of a monic equation for an element \(c\). They generate a finite algebra over the bottom ring. Adjoining \(c\) to it is finite over that algebra and hence finite over the bottom ring, by multiplying the two sets of module generators. Thus \(c\) is integral over the bottom ring.

Every contraction contains the kernel. Quotienting the base by the kernel reduces the converse to an inclusion \(R\subseteq A\). Given a prime \(\mathfrak p\) of \(R\), localize both rings at \(R\setminus\mathfrak p\). The inclusion remains injective, so the target is nonzero. The extension is still integral by localizing monic equations. Take a maximal ideal in the target. Its residue field is integral over the image of the localized base modulo its contraction. If a field \(L\) is integral over a subdomain \(D\), then \(D\) is a field: a monic equation for \(a^{-1}\), \(0\ne a\in D\), multiplied by \(a^{n-1}\), expresses \(a^{-1}\) in \(D\). Thus the contraction is maximal in the local base and is its unique maximal ideal. Contracting again to \(A\) and \(R\) gives a prime lying over \(\mathfrak p\).
\end{proof}

\begin{theorem}[Incomparability and dimension under integral maps]\label{p103:ALG-05-C01-S04}
For an integral map \(R\to A\), a prime \(\mathfrak q\) of \(A\) is maximal if and only if its contraction is maximal. If \(\mathfrak q\subseteq\mathfrak q'\) and both primes have the same contraction, they are equal. The Krull dimension of \(A\) equals that of \(R/\ker(R\to A)\).
\end{theorem}
\begin{proof}
The residue inclusion \(R/\mathfrak p\subseteq A/\mathfrak q\) is integral. A domain integral over a field is a field: a nonzero element has an algebraic equation of least degree, whose constant term is nonzero, and the equation solves for the inverse. The converse follows from the inverse argument in Theorem~\ref{p103:ALG-05-C01-R03}. For incomparability, quotient by the smaller prime and invert the nonzero elements of \(R/\mathfrak p\). The resulting domain is integral over a field, hence is a field. A nonzero prime with zero contraction would localize to a nonzero proper prime in this field, by the localization correspondence and injectivity of localization of the domain. This is impossible.

Any strict chain upstairs therefore contracts strictly in \(V(\ker g)\). Conversely a finite chain in \(V(\ker g)\) lifts by first using lying over and then, successively, applying lying over to \(R/\mathfrak p\subseteq A/\mathfrak q\) for the already chosen prime \(\mathfrak q\) over \(\mathfrak p\). This yields a lift containing \(\mathfrak q\). Take the supremum of finite chain lengths, allowing infinity. When either spectrum is empty so is the other by lying over, and both have the empty-spectrum dimension.
\end{proof}

\begin{theorem}[Going up and going down]\label{p103:ALG-05-C01-R04}
Every integral extension has going up. Given \(\mathfrak p\subseteq\mathfrak p'\) and \(\mathfrak q\) lying over \(\mathfrak p\), there is \(\mathfrak q'\supseteq\mathfrak q\) over \(\mathfrak p'\). If \(R\subseteq A\) is an integral extension of domains and \(R\) is integrally closed in its fraction field, it has going down. Given \(\mathfrak q'\) over \(\mathfrak p'\), a prime \(\mathfrak q\subseteq\mathfrak q'\) lies over \(\mathfrak p\). A flat ring map also has going down, independently of these integrality hypotheses.
\end{theorem}
\begin{proof}
The quotient argument in Theorem~\ref{p103:ALG-05-C01-S04} gives going up, including for a noninjective map when the given prime exists.

For going down with integrally closed base, first suppose \(F=\operatorname{Frac}A\) is finite over \(K=\operatorname{Frac}R\). Let \(L/K\) be a finite normal extension containing \(F\), and let \(T\) be the integral closure of \(R\) in \(L\). Then \(A\subseteq T\) is integral, since an equation over \(R\) is also an equation over \(A\). Lift \(\mathfrak q'\) to a prime \(Q_2'\) of \(T\) by lying over. Independently lift \(\mathfrak p\subseteq\mathfrak p'\) to a chain \(Q_1\subseteq Q_2\) in \(T\). The automorphisms of \(L/K\) preserve \(T\) and act transitively on its primes over \(\mathfrak p'\). To see this, if a prime \(Q\) over \(\mathfrak p'\) is outside the orbit of \(Q_2\), write \(P_1,\ldots,P_l\) for that orbit. Incomparability gives \(a_i\in Q\setminus P_i\) and \(b_{ij}\in P_j\setminus P_i\) for \(j\ne i\). Set \(a=\sum_i a_i\prod_{j\ne i}b_{ij}\). Modulo \(P_i\), only the \(i\)-th term survives and that term is nonzero, so \(a\in Q\) avoids every orbit prime. Each conjugate of \(a\) is also outside those primes. The field norm of \(a\) is a product of these conjugates with repetitions, including \(a\). This holds for finite normal inseparable extensions as well. Each distinct root of the minimal polynomial occurs with the same inseparable multiplicity, and the determinant norm further repeats by the degree over \(K(a)\). Thus the norm is integral over \(R\) and in \(K\), hence in \(R\). It belongs to \(Q\), so it is in \(\mathfrak p'\). The product expression shows it is outside \(Q_2\), a contradiction. An automorphism carrying \(Q_2\) to \(Q_2'\) carries \(Q_1\) to a prime contained in \(Q_2'\) over \(\mathfrak p\). Its contraction to \(A\) works.

For arbitrary \(F/K\) algebraic, enlarge \(A\) to the integral closure \(T\) of \(R\) in \(F\) and lift \(\mathfrak q'\) to \(Q'\) in \(T\). For an intermediate field \(K\subseteq E\subseteq F\), put \(T_E=T\cap E\), the integral closure of \(R\) in \(E\). It is integrally closed with fraction field \(E\). If \(z\in E\) is algebraic over \(K\), clearing denominators in its monic equation makes \(rz\) integral over \(R\) for some \(0\ne r\in R\), and transitivity gives the integrally closed assertion. Consider pairs \((E,P_E)\) with \(P_E\subseteq Q'\cap T_E\) contracting to \(\mathfrak p\), ordered by inclusion of fields and compatible prime contractions. The pair \((K,\mathfrak p)\) exists. A chain has an upper bound by taking the unions of fields and primes. The union prime is proper and primality is checked by taking any two elements in a common member of the chain. A maximal pair exists. If \(E\ne F\), adjoin one algebraic element to \(E\), obtaining a finite field extension \(E'/E\). The finite case applied to \(T_E\subseteq T_{E'}\), with upper prime \(Q'\cap T_{E'}\), extends \(P_E\) compatibly. This contradicts maximality. Thus the maximal prime is in \(T\), and contracting it to \(A\) proves the general case. This is the finite-normal-extension argument of \cite[Theorem 13.9 and Proposition 13.10]{eisenbud}.

For the flat case localize \(R\to A\) to the flat local map \(R_{\mathfrak p'}\to A_{\mathfrak q'}\). Such a map is faithfully flat. Indeed, a nonzero module contains a nonzero cyclic module \(R_{\mathfrak p'}/J\), where \(J\) is proper. After tensoring it is still nonzero since \(JA_{\mathfrak q'}\) is contained in the maximal ideal, and flatness preserves its injection into the given module. Consequently the tensor of \(A_{\mathfrak q'}\) with the residue field at \(\mathfrak pR_{\mathfrak p'}\) is nonzero. A prime of this fiber contracts to the designated base prime, by quotienting and localizing as in the spectrum correspondence. Its contraction to \(A\) is contained in \(\mathfrak q'\). Each of the one-step lifting results lifts a finite chain by iteration.
\end{proof}
\chapter{Noetherian modules and primary decomposition}
\label{p103:ch05}

A ring is Noetherian if its ideals satisfy the ascending chain condition. A module is Noetherian if its submodules satisfy the same condition. In either lattice the condition is equivalent to every subobject being finitely generated: a subobject with no finite generating set yields a strictly ascending sequence by successively adding elements, and a chain whose union is finitely generated stabilizes once the generators occur in it. We follow \cite[Sections 1.4 and 2.2 and Chapter 3]{eisenbud}.

\section{Finiteness and avoidance}

\begin{theorem}[Hilbert's basis theorem and finite modules]\label{p103:ALG-05-C02-R01}
If \(R\) is Noetherian, every polynomial ring in finitely many variables over \(R\) and every finite-type \(R\)-algebra is Noetherian. Every finite \(R\)-module and every submodule of such a module is both Noetherian and finitely presented.
\end{theorem}
\begin{proof}
An extension of two Noetherian modules is Noetherian. In an ascending chain, the intersections with the submodule and the images in the quotient stabilize. If two nested members have the same intersection and image, subtract a lift of the image of each element of the larger member to put its difference in the common intersection. Thus the nested members coincide. It follows by induction that finite free \(R\)-modules are Noetherian. A quotient or submodule of a Noetherian module is Noetherian. In particular, every finite module is a quotient of \(R^n\) with finitely generated kernel, and each of its submodules is finite and has a finite presentation by the same argument.

Given an ideal \(J\subseteq R[X]\), the leading coefficients of its nonzero members together with zero form an ideal \(L\subseteq R\). Align degrees by multiplying by powers of \(X\), and if the leading terms cancel their coefficient sum is zero. The scalar-multiplication check is the same, including when the scalar kills the leading coefficient. If \(L=0\), then \(J=0\). Otherwise choose finitely many \(f_i\in J\) with leading coefficients generating \(L\), and let \(D\) be the maximum of their degrees. Any \(f\in J\) of degree at least \(D\) can have its leading term canceled by a linear combination of suitable monomial multiples of the \(f_i\), giving a member of \(J\) of smaller degree. The members of \(J\) of degree less than \(D\) form a submodule of the finite free module \(R^D\), so finitely many generate them. These polynomials and the \(f_i\) generate \(J\) as an ideal. Induction gives the finite-variable version, including the zero-variable case. Every finite-type algebra is a quotient of such a polynomial ring and is Noetherian.
\end{proof}

\begin{lemma}[Finite prime avoidance]\label{p103:ALG-05-C02-S01}
Let \(J,I_1,\ldots,I_n\) be ideals with \(n\ge1\) and \(J\subseteq\bigcup_i I_i\). If at most two of the \(I_i\) are nonprime, or \(R\) contains an infinite field, then \(J\subseteq I_i\) for some \(i\).
\end{lemma}
\begin{proof}
For at most two nonprime ideals, remove any ideal contained in another on the list. For \(n=1\) there is nothing to show. For \(n=2\), if \(a,b\in J\) lie outside \(I_1,I_2\), respectively, then either one already lies outside the union or \(a+b\) does. For \(n\ge3\) arrange that \(I_n\) is prime. If \(J\) is contained in none of the ideals, induction gives \(x\in J\) outside \(I_1,\ldots,I_{n-1}\). If \(x\in I_n\), choose \(y\in J\setminus I_n\) and \(b_i\in I_i\setminus I_n\) for \(i<n\). The sum \(x+y\prod_{i<n}b_i\) avoids each \(I_i\) with \(i<n\) because the second term is in all those ideals. It avoids \(I_n\) by primality. This is a contradiction.

If \(R\) contains an infinite field \(k\), choose for each \(i\) a witness \(x_i\in J\setminus I_i\) under the supposition that no \(I_i\) contains \(J\). The finite-dimensional \(k\)-span \(W\) of these witnesses is contained in \(J\), and each \(W\cap I_i\) is a proper subspace. Choose a nonzero linear form on \(W\) vanishing on it. Their product is a nonzero polynomial, which cannot vanish on all of \(W\). In coordinates this follows by induction on the number of variables from the univariate root bound over an infinite field. A point where the product is nonzero avoids the union.
\end{proof}

\section{Associated primes and support}

For an \(R\)-module \(M\), \(\operatorname{Ass}_R(M)\) is the set of prime ideals occurring as annihilators of elements of \(M\). Equivalently \(\mathfrak p\) is in this set if \(R/\mathfrak p\) embeds in \(M\). Define \(\operatorname{Supp}_R(M)=\{\mathfrak p:M_{\mathfrak p}\ne0\}\). We will use the following elementary observation repeatedly. In a nonzero module over a Noetherian ring, choose an annihilator \(\operatorname{Ann}(x)\) maximal among the annihilators of nonzero elements. If \(ab\) kills \(x\) but \(b\) does not, then \(bx\ne0\), and \(\operatorname{Ann}(bx)\) contains both \(\operatorname{Ann}(x)\) and \(a\). Maximality implies that \(a\) kills \(x\). Thus the chosen annihilator is prime.

\begin{theorem}[Localization, exact sequences, and zero divisors]\label{p103:ALG-05-C02-S03}
Let \(R\) be Noetherian and \(M\) a finite \(R\)-module. Then
\[
 \operatorname{Ass}_{S^{-1}R}(S^{-1}M)
 =\{S^{-1}\mathfrak p:\mathfrak p\in\operatorname{Ass}_R(M),\ 
      \mathfrak p\cap S=\varnothing\}.
\]
The set of elements of \(R\) that kill some nonzero member of \(M\) is \(\bigcup_{\mathfrak p\in\operatorname{Ass}_R(M)}\mathfrak p\). For a short exact sequence of finite modules \(0\to M'\to M\to M''\to0\),
\[
 \operatorname{Ass}(M')\subseteq\operatorname{Ass}(M)
  \subseteq\operatorname{Ass}(M')\cup\operatorname{Ass}(M''),\qquad
 \operatorname{Supp}(M)=\operatorname{Supp}(M')\cup\operatorname{Supp}(M'').
\]
In a direct sum the associated-prime set is the union of the two sets.
\end{theorem}
\begin{proof}
If \(R/\mathfrak p\) embeds in \(M\) and \(\mathfrak p\) avoids \(S\), localization preserves the injection and \(S^{-1}(R/\mathfrak p)\) is the cyclic module with annihilator \(S^{-1}\mathfrak p\). Conversely suppose \(m/s\in S^{-1}M\) has prime annihilator \(P\). Write \(\mathfrak p=P\cap R\) and choose generators \(r_1,\ldots,r_k\) of \(\mathfrak p\). For each \(i\) some \(u_i\in S\) kills \(r_i m\). Then \(x=(\prod_i u_i)m\) is killed by \(\mathfrak p\). If \(r x=0\), the image of \(x\) in the localization is a unit times \(m/s\), so \(r\in\mathfrak p\). Hence \(\mathfrak p=\operatorname{Ann}(x)\), as required. The empty set of generators is harmless, with product \(1\).

If \(r\) kills a nonzero element, the nonzero submodule \(\ker(r:M\to M)\) has an element with prime annihilator by the maximal annihilator argument. That prime belongs to \(\operatorname{Ass}(M)\) and contains \(r\). Membership of \(r\) in an associated prime supplies a witness for the converse. The first exact-sequence inclusion follows from the injection \(M'\to M\). For the second, given an embedded cyclic module \(R/\mathfrak p\subseteq M\), if its intersection with \(M'\) is nonzero, any nonzero element of that intersection has annihilator \(\mathfrak p\) because \(R/\mathfrak p\) is a domain. Otherwise the cyclic module embeds in \(M''\). For a direct sum both summands embed in the sum. The support equality follows by localizing the exact sequence and observing that the middle term is zero precisely when both ends are zero.
\end{proof}

\begin{theorem}[Finiteness of associated primes and support]\label{p103:ALG-05-C02-R03}
Let \(R\) be Noetherian and \(M\) a finite \(R\)-module. Then \(\operatorname{Ass}(M)\) is finite, and it is nonempty when \(M\ne0\). The support is \(V(\operatorname{Ann}M)\), and every minimal member of the support is associated. If \(M=0\), both sets are empty.
\end{theorem}
\begin{proof}
Build a chain of submodules of \(M\) by embedding a cyclic \(R/\mathfrak p\) in each successive nonzero quotient, using the maximal annihilator argument. Its inverse image strictly enlarges the current submodule. Since \(M\) is Noetherian, the process ends in a finite filtration with prime cyclic factors. Theorem~\ref{p103:ALG-05-C02-S03} bounds \(\operatorname{Ass}(M)\) by the finite set of primes in this filtration. It is nonempty if \(M\ne0\), by the same maximal annihilator argument.

If \(\operatorname{Ann}M\) meets \(R\setminus\mathfrak p\), \(M_{\mathfrak p}=0\). Conversely if \(M_{\mathfrak p}=0\), select a denominator killing each member of a finite set of generators of \(M\). Their product annihilates \(M\) and avoids \(\mathfrak p\). Thus the support equals the displayed closed set. If \(\mathfrak p\) is minimal in the support, then \(M_{\mathfrak p}\ne0\) and its support over \(R_{\mathfrak p}\) consists just of the maximal ideal, by the localization correspondence and the support formula. Hence this maximal ideal is associated, and the localization formula in Theorem~\ref{p103:ALG-05-C02-S03} implies that \(\mathfrak p\) is associated to \(M\).
\end{proof}

\section{Primary submodules and their uniqueness}

A proper submodule \(N\subsetneq M\) is \emph{primary} if \(r m\in N\) and \(m\notin N\) imply \(r^nM\subseteq N\) for some \(n\ge1\). In this case \(\sqrt{\operatorname{Ann}(M/N)}\) is a prime and \(N\) is called primary for that prime. Indeed, in the nonzero quotient, if \((ab)^n\) annihilates the quotient and \(b^n\) does not, take \(m\) with \(b^nm\ne0\) and apply the primary condition to \(a^n\) and \(b^nm\). This puts \(a\) in the radical. The radical is precisely the set of zero divisors on \(M/N\). Every member of the radical is a zero divisor there, since a nilpotent injective endomorphism of a nonzero module is impossible. The primary condition proves the converse. Thus over a Noetherian ring the associated-prime set of \(M/N\) is the singleton consisting of this radical. Any associated prime contains the radical, since it is prime and contains the module annihilator, and is contained in the radical by the zero-divisor criterion. There is at least one associated prime by the maximal annihilator argument.

\begin{theorem}[Lasker--Noether decomposition for modules]\label{p103:ALG-05-C02-S02}
Let \(R\) be Noetherian, \(M\) a finite module, and \(N\subsetneq M\) a proper submodule. It has a finite decomposition
\[
 N=\bigcap_{i=1}^h N_i
\]
into primary submodules with distinct associated radicals \(\mathfrak p_i=\sqrt{\operatorname{Ann}(M/N_i)}\). The decomposition can be chosen irredundant, and for any such decomposition
\(\operatorname{Ass}(M/N)=\{\mathfrak p_1,\ldots,\mathfrak p_h\}\).
For every multiplicatively closed \(S\),
\[
 \{m\in M:m/1\in S^{-1}N\}
   =\bigcap_{\mathfrak p_i\cap S=\varnothing}N_i.
\]
In particular the component belonging to a minimal associated prime is unique.
\end{theorem}
\begin{proof}
Call a proper submodule irreducible if expressing it as an intersection of two submodules forces one to equal it. Every submodule is a finite intersection of irreducible ones, with the empty intersection allowed for \(M\). Otherwise the ascending chain condition gives a maximal counterexample. It cannot be irreducible, and writing it as an intersection of two strictly larger submodules would give finite decompositions of both. This is a contradiction.

An irreducible submodule is primary. Pass to the quotient so that the irreducible submodule is zero. If \(r\) kills a nonzero element, choose \(n\ge1\) at which the ascending chain \(\ker r^j\) has stabilized. Then \(r^nM\cap\ker r^n=0\). If \(r^nx\) is in the kernel, \(x\in\ker r^{2n}=\ker r^n\). The kernel is nonzero, so irreducibility forces \(r^nM=0\). Hence every zero divisor is nilpotent on the quotient.

An intersection of finitely many submodules primary for a fixed \(\mathfrak p\) is primary for \(\mathfrak p\). In the quotient by the intersection, multiplication by \(r\notin\mathfrak p\) is injective since the quotient embeds in the direct sum of the primary quotients, whereas \(r\in\mathfrak p\) acts nilpotently by taking a common exponent for the finitely many components. Therefore components with a common prime can be combined. Remove redundant components. Since \(N\ne M\), at least one proper component remains.

The diagonal embedding \(M/N\to\bigoplus_i M/N_i\) and Theorem~\ref{p103:ALG-05-C02-S03} put its associated primes among the \(\mathfrak p_i\). Conversely irredundancy provides an element in \(\bigcap_{j\ne i}N_j\) outside \(N_i\). Its class in \(M/N\) generates a nonzero cyclic submodule that embeds into \(M/N_i\): a multiple that belongs to \(N_i\) belongs to every component and hence to \(N\). The cyclic module has an associated prime, which must be \(\mathfrak p_i\), so this prime is associated to \(M/N\).

Localization preserves the finite intersection. If \(S\) meets \(\mathfrak p_i\), some power of an element of \(S\) kills \(M/N_i\), so the localization of that quotient vanishes. If it avoids \(\mathfrak p_i\), no member of \(S\) kills a nonzero class in \(M/N_i\). Hence the contraction to \(M\) of \(S^{-1}N_i\) is \(N_i\). This proves the formula, with empty intersection \(M\). For a minimal associated prime \(\mathfrak p\), take \(S=R\setminus\mathfrak p\). Exactly the component with prime \(\mathfrak p\) remains, and the formula identifies it intrinsically.
\end{proof}

\begin{corollary}[Primary decomposition of ideals]\label{p103:ALG-05-C02-R02}
Every proper ideal \(I\) of a Noetherian ring is a finite irredundant intersection of primary ideals with distinct radicals. The radicals are precisely \(\operatorname{Ass}_R(R/I)\). The component for each minimal associated prime is uniquely determined, whereas an embedded component need not be unique.
\end{corollary}
\begin{proof}
Take \(M=R\) in Theorem~\ref{p103:ALG-05-C02-S02}. The primary condition then has its ideal form \(ab\in Q,\ b\notin Q\Rightarrow a\in\sqrt Q\). For nonuniqueness, in \(k[X,Y]\) let \(I=(X^2,XY)\). Then
\[
 I=(X)\cap(X^2,Y)=(X)\cap(X^2,XY,Y^2).
\]
The first equality follows by checking monomials divisible by \(X\) that lie in \((X^2,Y)\). The second follows by the same check for \((X,Y)^2\). The last two ideals are distinct, and both have radical \((X,Y)\). Their quotients have nilpotent maximal ideal and every element outside it is a unit, so both are primary. Both intersections are irredundant, since \((X)\) and either other component are incomparable. Thus the component for the embedded maximal ideal varies.
\end{proof}
\chapter[Krull dimension and Hilbert functions]
{Krull dimension, affine algebras, and Hilbert functions}
\label{p103:ch06}

The length of a finite strictly increasing chain of primes is its number of strict inclusions. Krull dimension is the supremum of these lengths. The height of a prime is the dimension of the localization at that prime by Theorem~\ref{p103:ALG-05-C01-R02}. We allow infinite dimension and use the formal value \(-\infty\) for the empty spectrum. For a finite module \(M\) we write \(\dim M=\dim(R/\operatorname{Ann}M)\), equivalently the dimension of its support by Theorem~\ref{p103:ALG-05-C02-R03}.

\section{Height and systems of parameters}

\begin{theorem}[Krull's height theorem]\label{p103:ALG-05-C03-R01}
Let \(R\) be Noetherian and let \(\mathfrak p\) be minimal among the primes containing a proper ideal generated by \(r\) elements. Then \(\operatorname{ht}\mathfrak p\le r\). In particular a prime minimal over a principal ideal has height at most one.
\end{theorem}
\begin{proof}
We give the symbolic-power argument of \cite[Theorems 10.1--10.2]{eisenbud}. Localize at \(\mathfrak p\), so \((R,\mathfrak m)\) is local and the radical of the ideal generated by the given elements is \(\mathfrak m\). If \(r=0\), \(\mathfrak m\) is the only prime. For \(r=1\) write \(x\) for the generator. The ring \(R/(x)\) is Artinian. Its maximal ideal is its nilradical, finitely generated by nilpotents, hence nilpotent. Its successive powers have quotients finite-dimensional over the residue field, giving finite length.

Use the commutative local specialization of Theorem~\ref{p103:ALG-03-C03-R02} for Nakayama's lemma. Fix any prime \(\mathfrak q\subsetneq\mathfrak m\). Then \(x\notin\mathfrak q\). Define \(\mathfrak q^{(n)}=R\cap \mathfrak q^nR_{\mathfrak q}\), where the intersection denotes inverse image under the localization map. The descending images of \(\mathfrak q^{(n)}\) in \(R/(x)\) stabilize. Thus for large \(n\), every \(a\in\mathfrak q^{(n)}\) can be written \(a=xb+c\) with \(c\in\mathfrak q^{(n+1)}\). Since \(x\) is a unit at \(\mathfrak q\), this implies \(b\in\mathfrak q^{(n)}\). Hence
\(\mathfrak q^{(n)}=x\mathfrak q^{(n)}+\mathfrak q^{(n+1)}\).
Nakayama on the finite quotient yields equality of the two symbolic powers. Localizing at \(\mathfrak q\) gives equality of two consecutive ordinary powers of the maximal ideal there. Another application of Nakayama makes that power zero. Therefore \(R_{\mathfrak q}\) has just one prime, \(\mathfrak q\) has height zero, and every strict chain ending at \(\mathfrak m\) has length at most one.

For \(r>1\), write the generators as \(x_1,\ldots,x_r\). Any prime strictly below \(\mathfrak m\) is contained in a prime \(\mathfrak q\) immediately below \(\mathfrak m\), by the ascending chain condition. Some \(x_i\) is outside \(\mathfrak q\). Renumber it \(x_1\). Then \(\sqrt{(\mathfrak q,x_1)}=\mathfrak m\), so for each \(i\ge2\) write \(x_i^{e_i}=a_ix_1+y_i\) with \(y_i\in\mathfrak q\). Any prime containing \((x_1,y_2,\ldots,y_r)\) contains every \(x_i\), so is \(\mathfrak m\). Choose a prime \(\mathfrak q_0\subseteq\mathfrak q\) minimal over \((y_2,\ldots,y_r)\). Such a prime exists by Zorn's lemma, since the intersection of a descending chain of primes is prime. In \(R/\mathfrak q_0\) the maximal ideal is minimal over \((x_1)\). The principal case forbids the strict chain
\(\mathfrak q_0\subsetneq\mathfrak q\subsetneq\mathfrak m\).
Thus \(\mathfrak q\) is minimal over \((y_2,\ldots,y_r)\), and the induction hypothesis bounds its height by \(r-1\). Apply this to the next-to-last prime of any chain ending at \(\mathfrak m\), extending it to such a \(\mathfrak q\) if necessary. This bounds the length by \(r\). Relative versions of the argument apply in the specified quotient ring, where height is measured anew.
\end{proof}

\begin{theorem}[Systems of parameters]\label{p103:ALG-05-C03-S01}
Let \((R,\mathfrak m)\) be Noetherian local and \(M\ne0\) a finite module of dimension \(d\). For an ideal \(Q\subseteq\mathfrak m\), the quotient \(M/QM\) has finite length if and only if \(\sqrt{Q+\operatorname{Ann}M}=\mathfrak m\). The least number of elements of \(\mathfrak m\) needed to generate an ideal with that property is \(d\). In particular for \(x\in\mathfrak m\) one has
\(\dim(M/xM)\ge d-1\).
\end{theorem}
\begin{proof}
For each prime \(\mathfrak p\), the localization of \(M/QM\) is zero if \(\mathfrak p\) is outside \(\operatorname{Supp}M\) or does not contain \(Q\). Conversely if \(\mathfrak p\) is in the support and contains \(Q\), then \(M_{\mathfrak p}/Q_{\mathfrak p}M_{\mathfrak p}\ne0\) by local Nakayama. Thus the support of \(M/QM\) is \(V(Q+\operatorname{Ann}M)\). A finite module over \(R\) with support \(\{\mathfrak m\}\) is killed by a power of \(\mathfrak m\). The radical of its annihilator is \(\mathfrak m\) and its finitely many generators have powers in the annihilator. The filtration by powers of \(\mathfrak m\) then gives finite length. Conversely any finite composition series has factors \(R/\mathfrak m\), so support is \(\{\mathfrak m\}\) for a nonzero finite-length module. Here \(M/QM\) is nonzero by Nakayama. For the omitted ideal \(Q=R\), the quotient is zero and has length zero.

Replace \(R\) by \(R/\operatorname{Ann}M\) for the minimum calculation. The height theorem applied to a finite set of generators of the maximal ideal first shows \(d<\infty\), and applied to any ideal with radical the maximal ideal gives the lower bound \(r\ge d\). For the upper bound use induction on \(d\). If \(d=0\), no generators are needed. Otherwise the minimal primes are finitely many by Theorem~\ref{p103:ALG-05-C02-R03}, and each is properly contained in the maximal ideal. Prime avoidance gives \(x\in\mathfrak m\) outside all the minimal primes. Any chain in \(V(x)\) can be lengthened at its bottom by placing a minimal prime below its first prime, since the latter contains \(x\) and is not minimal. Consequently the quotient ring by \(x\) has dimension at most \(d-1\). Apply induction there and add \(x\) to the resulting generators. For arbitrary \(x\in\mathfrak m\), use a system of parameters for \(M/xM\) and prepend \(x\). This gives \(d\le 1+\dim(M/xM)\), including \(d=0\). These are the finite-colength and parameter arguments of \cite[Section 10.1]{eisenbud}.
\end{proof}

\section{Normalization and affine dimension}

\begin{lemma}[Making a polynomial monic]\label{p103:normalization-coordinate}
Given a nonconstant polynomial \(f\in k[X_1,\ldots,X_n]\) over a field and \(n\ge1\), a triangular change of variables and multiplication by a scalar make \(f\) monic in the last variable.
\end{lemma}
\begin{proof}
Choose an integer \(b\) larger than every individual exponent in the monomials of \(f\). Substitute \(X_i=Y_i+Y_n^{b^i}\) for \(1\le i<n\), and \(X_n=Y_n\). The greatest possible power of \(Y_n\) from a monomial with exponents \(a_i\) is \(a_n+\sum_{i<n}a_ib^i\), obtained uniquely by taking the pure \(Y_n\) terms. Distinct exponent vectors give distinct numbers. Thus the greatest power in the transformed polynomial has a single nonzero scalar coefficient and positive degree. Divide by it. The change is invertible by subtracting the same powers. For \(n=1\) only the rescaling is needed. No cardinality assumption on \(k\) enters.
\end{proof}

\begin{lemma}[Polynomial rings over a field]\label{p103:polynomial-dimension}
For any field \(k\), the polynomial ring \(k[X_1,\ldots,X_n]\) has dimension \(n\). If \(P\) is a nonzero prime, the quotient by \(P\) has dimension at most \(n-1\).
\end{lemma}
\begin{proof}
The coordinate prime chain gives the lower bound for the polynomial ring. Induct on \(n\) for the upper bound, starting with the field \(k\). For a nonzero prime \(P\), a nonzero \(f\in P\) is nonconstant. Lemma~\ref{p103:normalization-coordinate} makes it monic in the last variable, so the quotient by \(P\) is integral over the image of the polynomial ring in the other \(n-1\) variables. That image is a quotient of a ring of dimension \(n-1\), so the quotient by \(P\) has dimension at most \(n-1\), by Theorem~\ref{p103:ALG-05-C01-S04}. In any finite prime chain of positive length, if its first prime is nonzero apply this bound to the entire chain modulo that prime. If its first prime is zero, apply it to the part starting at the second prime. This gives the upper bound \(n\).
\end{proof}

\begin{theorem}[Noether normalization]\label{p103:ALG-05-C03-R02}
Let \(A\ne0\) be a finite-type algebra over an arbitrary field \(k\). There are algebraically independent \(y_1,\ldots,y_d\in A\) such that \(A\) is finite over \(k[y_1,\ldots,y_d]\). In fact \(d=\dim A\).
\end{theorem}
\begin{proof}
Choose a surjection from a polynomial ring in \(n\) variables and induct on \(n\). With zero kernel, the images of the variables themselves work, including \(n=0\). Otherwise a nonzero relation in its kernel is nonconstant since \(A\ne0\). By Lemma~\ref{p103:normalization-coordinate} take it monic in the last variable. Thus the image of the last variable is integral over the subalgebra generated by the others, and \(A\) is finite over that subalgebra. The subalgebra is nonzero and has fewer generators, so induction and transitivity of module finiteness prove existence. The dimension assertion follows from integral dimension invariance and Lemma~\ref{p103:polynomial-dimension}. The argument is the unrestricted-field normalization in \cite[Lemma 13.2 and Theorem 13.3]{eisenbud}.
\end{proof}

\begin{theorem}[Dimension of polynomial rings and affine domains]\label{p103:ALG-05-C03-R03}
For any field \(k\) and finite \(n\), \(\dim k[X_1,\ldots,X_n]=n\). For a finite-type \(k\)-domain \(A\),
\[
 \dim A=\operatorname{trdeg}_k\operatorname{Frac}A .
\]
For a nonzero general affine algebra its dimension is the maximum of the dimensions of its quotients by minimal primes.
\end{theorem}
\begin{proof}
The polynomial assertion is Lemma~\ref{p103:polynomial-dimension}. For an affine domain use Theorem~\ref{p103:ALG-05-C03-R02}. Its fraction field is algebraic over the fraction field of the polynomial subalgebra. Localizing \(A\) at the nonzero elements of that subalgebra gives a domain finite over a field and hence a field, necessarily \(\operatorname{Frac}A\). The independent variables give a transcendence basis. For a general affine algebra, each prime chain lies above some minimal prime and conversely each chain in a minimal-prime quotient lifts, so the stated maximum follows. The minimal primes are finite by Noetherianness and Theorem~\ref{p103:ALG-05-C02-R03}.
\end{proof}

\begin{theorem}[Weak and strong Nullstellensatz]\label{p103:ALG-05-C01-S05}
Let \(A\) be a finite-type algebra over any field \(k\). Every maximal ideal has residue field finite over \(k\), and every prime ideal is the intersection of the maximal ideals containing it. If \(k\) is algebraically closed, the maximal ideals in \(k[X_1,\ldots,X_n]\) are the kernels of evaluation at points of \(k^n\), and for every ideal \(J\),
\[
 I(V(J))=\sqrt J ,
\]
where \(V(J)\subseteq k^n\) is the common zero set and \(I(V(J))\) its vanishing ideal.
\end{theorem}
\begin{proof}
A field which is a finite-type \(k\)-algebra is finite over \(k\). Indeed, normalize it and observe that the polynomial subalgebra must itself be a field, since a field integral over a subdomain makes the subdomain a field. Thus the subalgebra has zero variables. Apply this to each maximal quotient.

Let \(P\) be prime and \(a\notin P\). The ring \((A/P)_{\bar a}\) is nonzero and of finite type over \(k\), with one more generator for the inverse of \(\bar a\). Take a maximal ideal there. Its contraction to \(A/P\) misses \(\bar a\). The quotient by the contraction is a \(k\)-subalgebra of the resulting finite extension field of \(k\). Any such subalgebra is a field, because a nonzero algebraic element has inverse in its polynomial subalgebra. Thus the contraction is maximal. Vary \(a\notin P\) to get the intersection assertion.

For algebraically closed \(k\), the maximal residue field is \(k\). The images of the variables are a point and the quotient map is its evaluation. Conversely every point evaluation is surjective with maximal kernel. The ideals of evaluation at points of \(V(J)\) are exactly the maximal ideals containing \(J\). Every prime containing \(J\) is an intersection of maximal ideals containing it, so Theorem~\ref{p103:ALG-05-C01-S02} identifies the intersection of all those maximal ideals with \(\sqrt J\). This includes the unit ideal, where \(V(J)\) is empty, and zero variables, where the ambient affine space is a singleton. The proof follows \cite[Theorem 1.6 and Corollary 13.12]{eisenbud}.
\end{proof}

\begin{theorem}[Finiteness of the integral closure of an affine domain]\label{p103:ALG-05-C01-S06}
Let \(A\) be a finite-type domain over any field \(k\) and \(L/\operatorname{Frac}A\) a finite field extension. The integral closure of \(A\) in \(L\) is finite as an \(A\)-module and is itself an affine \(k\)-domain.
\end{theorem}
\begin{proof}
By normalization \(A\) is finite over a polynomial ring \(B=k[t_1,\ldots,t_d]\). If \(K=\operatorname{Frac}B\), then \(L/K\) is finite. An element of \(L\) is integral over \(A\) exactly when it is integral over \(B\), by transitivity and inclusion. It suffices to prove module finiteness over \(B\), since the same generating set then generates over \(A\). The polynomial ring is Noetherian and integrally closed: Gauss's lemma makes it a unique factorization domain, and if a reduced fraction satisfies a monic equation then its denominator divides a power of its coprime numerator, so is a unit.

More generally let \(D\) be a Noetherian integrally closed domain, \(E=\operatorname{Frac}D\), and \(F/E\) a finite separable extension. Choose an \(E\)-basis \(v_i\) of \(F\) consisting of elements integral over \(D\), by clearing denominators in a monic equation after multiplying each initial basis member by a nonzero scalar from \(D\). The trace pairing is nondegenerate. For a primitive element with distinct conjugates the matrix of their values on the power basis is Vandermonde and invertible, and its transpose times itself is the trace matrix over a splitting field. If \(v_i^*\) is the dual basis, any \(z\in F\) integral over \(D\) satisfies
\[
 z=\sum_i \operatorname{Tr}_{F/E}(zv_i)\,v_i^* .
\]
The trace coefficients are integral over \(D\), as sums of integral conjugates, and belong to \(E\), so they are in \(D\). Thus the integral closure is a \(D\)-submodule of \(\sum_i Dv_i^*\) and is finite by Noetherianness. This proves the claim for \(B\) when \(L/K\) is separable.

In characteristic \(p>0\) let \(F\) be the maximal separable subextension of \(L/K\), and choose \(q=p^e\) so that \(L^q\subseteq F\). Such an \(e\) exists because \(L/F\) is finite purely inseparable. Take a common exponent for finitely many field generators. Take a primitive element \(\theta\) for \(F/K\) with separable minimal polynomial \(f\), and express the \(q\)-th powers of the chosen generators as \(K\)-linear combinations of powers of \(\theta\). Adjoin to \(k\) the \(q\)-th roots of the finitely many constants occurring in the rational functions that are coefficients of \(f\) and of those combinations. This gives a finite extension \(k'/k\) inside an algebraic closure. Put \(B'=k'[t_1^{1/q},\ldots,t_d^{1/q}]\) and \(K'=\operatorname{Frac}B'\). Here \(k'^q\subseteq k\), because \(k'\) was formed by adjoining \(q\)-th roots of constants from \(k\). The generators \(t_i^{1/q}\) are algebraically independent over \(k'\): raising any nonzero polynomial relation to the \(q\)-th power would give a nonzero polynomial relation among the \(t_i\) with coefficients in \(k\). The ring \(B'\) is finite over \(B\) by bounded powers of the roots and a finite \(k\)-basis of \(k'\).

Now \(\theta^{1/q}\) satisfies the polynomial obtained by taking the unique \(q\)-th roots of the coefficients of \(f\). They lie in \(K'\). This polynomial has distinct roots, since taking \(q\)-th powers sends its roots bijectively to the distinct roots of \(f\), preserving the count given by the degree. Thus \(L'=K'(\theta^{1/q})\) is finite separable over \(K'\). Taking \(q\)-th roots of the displayed expressions for the \(q\)-th powers of the generators of \(L\) shows \(L\subseteq L'\). The integral closure of \(B'\) in \(L'\) is finite over \(B'\) by the separable case, hence finite over \(B\). Every element of \(L\) integral over \(B\) lies in that finite \(B\)-module. These elements form a \(B\)-submodule and therefore are finite over \(B\), as desired. This is the trace and purely inseparable reduction of \cite[Corollary 13.13 and Proposition 13.14]{eisenbud}.
\end{proof}

\section{Graded Hilbert functions}

A nonnegatively graded algebra \(A=\bigoplus_{j\ge0}A_j\) is standard over \(A_0\) when it is generated as an \(A_0\)-algebra by finitely many elements of degree one. A finite graded module may be \(\mathbb Z\)-graded. It is bounded below because it has finitely many homogeneous generators. Write \(M(-1)_j=M_{j-1}\).

\begin{lemma}[The Hilbert series with an Artinian base]\label{p103:hilbert-serre-length}
If \(A_0=B\) is Artinian, \(A\) is standard over \(B\), and \(M\) is a finite graded \(A\)-module, then each \(M_j\) has finite length over \(B\). The series
\[
 H_M(t)=\sum_j \ell_B(M_j)t^j
\]
is a rational function with denominator a power of \(1-t\) and Laurent polynomial numerator. Its coefficients are eventually polynomial in \(j\). If the reduced denominator is \((1-t)^e\) with \(e>0\), the eventual degree is \(e-1\). If \(e=0\), the coefficients are eventually zero.
\end{lemma}
\begin{proof}
There is a homogeneous surjection from a polynomial ring in finitely many degree-one variables onto \(A\). Degree by degree the shifted free modules over that polynomial ring are finite over \(B\), and thus so is \(M_j\). Artinian commutative \(B\) has finite length over itself, so its finite modules have finite length. View \(M\) as a finite graded module over the polynomial ring and induct on the number of variables. With none it is supported in finitely many degrees. In the induction let \(x\) be the last variable, \(K=0:_M x\) and \(C=M/xM\). Both are finite graded modules by Noetherianness of the polynomial ring, and both are killed by \(x\), so the induction applies. From
\[
 0\longrightarrow K(-1)\longrightarrow M(-1)
   \xrightarrow{x} M\longrightarrow C\longrightarrow0
\]
obtain \((1-t)H_M=H_C-tH_K\). This yields the denominator bound. If \(P(t)\) is the Laurent numerator over the reduced denominator \((1-t)^e\), for large \(j\) its coefficients are finite linear combinations of \(\binom{j-a+e-1}{e-1}\) with coefficients from the terms \(t^a\) of \(P\). The leading coefficient is \(P(1)/(e-1)!\ne0\). If \(e=0\), the series itself is a Laurent polynomial.
\end{proof}

\begin{theorem}[Hilbert--Serre and dimension]\label{p103:ALG-05-C03-R04}
Let \(A\) be standard graded over a field \(k\), and let \(M\ne0\) be a finite graded \(A\)-module. If \(d=\dim M>0\), the reduced denominator of \(H_M\) has exponent \(d\), and \(j\mapsto\dim_k M_j\) agrees eventually with a polynomial of degree \(d-1\). If \(d=0\), the graded pieces are eventually zero. The zero module has zero series and polynomial.
\end{theorem}
\begin{proof}
Pass to the graded quotient \(C=A/\operatorname{Ann}M\), so \(M\) is faithful. The annihilator is homogeneous since one can test each homogeneous component of an element of \(A\) on homogeneous generators. If \(k\) is finite, extend it to the infinite field \(k(u)\). Each graded dimension is unchanged. Faithfulness is preserved. For homogeneous module generators \(m_i\), the injection \(C\to\bigoplus_i M\), \(c\mapsto(cm_i)_i\), remains injective after field extension. Dimension of \(C\) is unchanged as well. Use the ordinary Noether normalization \(k[y_1,\ldots,y_d]\hookrightarrow C\), whose injectivity and finiteness survive tensoring over \(k\), and apply integral dimension invariance over the larger field.

Over an infinite field the normalization of the standard graded algebra \(C\) can be taken in degree one. Start from a homogeneous surjection from a standard polynomial ring. If its kernel contains a nonzero homogeneous polynomial, choose a vector where that polynomial is nonzero, possible by the root argument in Lemma~\ref{p103:ALG-05-C02-S01}. A linear change of variables puts this vector on the last coordinate axis, so the relation becomes monic in that variable after scaling, with coefficients in the other variables. Repeat the induction in Theorem~\ref{p103:ALG-05-C03-R02}, retaining linear generators throughout. The resulting graded polynomial ring \(P\subseteq C\) has \(d\) variables, by the dimension equality. The module \(M\) is finite and faithful over \(P\).

If \(d=0\), \(M\) is finite-dimensional over \(k\), and only finitely many degrees are nonzero. For \(d>0\), let \(m_i\) be homogeneous \(P\)-module generators of degrees \(a_i\). In degree \(j\), the surjection from shifted copies of \(P\) gives
\(\dim_k M_j\le\sum_i\dim_k P_{j-a_i}\).
The injection \(P\to\bigoplus_i M\) by multiplication on these generators gives
\(\dim_k P_j\le\sum_i\dim_k M_{j+a_i}\).
For large \(j\), \(\dim_k P_j=\binom{j+d-1}{d-1}\). Comparing both inequalities with the eventual polynomial from Lemma~\ref{p103:hilbert-serre-length} forces its degree to be exactly \(d-1\). The relation with the reduced denominator in that lemma then gives exponent \(d\). Compare \cite[Section 1.9 and the graded dimension discussion in Section 13.4]{eisenbud}.
\end{proof}
\chapter[Completion and regular local rings]
{Filtrations, completion, and regular local rings}
\label{p103:ch07}

\section{Nakayama's lemma and Artin--Rees}

For a commutative ring \(R\), its Jacobson radical \(\operatorname{Jac}R\) is the intersection of its maximal ideals. If \(a\in\operatorname{Jac}R\), then \(1-a\) is a unit, for otherwise a maximal ideal would contain both \(a\) and \(1-a\). In a local ring \((R,\mathfrak m)\), the nonunits are precisely \(\mathfrak m\). Conversely, if the nonunits form a proper ideal, every proper ideal is contained in it and it is the unique maximal ideal.

\begin{theorem}[Nakayama over a commutative ring]\label{p103:ALG-05-C04-R01}
Let \(M\) be a finite \(R\)-module and \(I\subseteq\operatorname{Jac}R\). If \(M=N+IM\) for a submodule \(N\), then \(M=N\). In particular \(IM=M\) implies \(M=0\), and generators of \(M/IM\) lift to generators of \(M\).
\end{theorem}
\begin{proof}
Passing to \(M/N\) reduces the assertion to \(IM=M\). Choose generators \(m_i\) of \(M\) and a matrix \(U\) over \(I\) with \(\mathbf m=U\mathbf m\). The adjugate of \(1-U\) gives \(\det(1-U)\mathbf m=0\). The determinant is \(1-a\) for some \(a\in I\), a unit, so \(M=0\). With no generators the module is zero already. If \(N\) is the submodule spanned by proposed lifts, \(M=N+IM\) expresses exactly that their images generate. See \cite[Corollary 4.8]{eisenbud} and the general Nakayama lemma in Theorem~\ref{p103:ALG-03-C03-R02}.
\end{proof}

Let \(I\) be an ideal of a Noetherian ring \(R\), let \(M\) be a finite module, and let \(t\) mark degree. The Rees algebra of \(I\) and Rees module of \(M\) are
\[
 \mathcal R_I(R)=\bigoplus_{n\ge0} I^nt^n\subseteq R[t],
 \qquad \mathcal R_I(M)=\bigoplus_{n\ge0} I^nMt^n\subseteq M[t].
\]
The algebra is generated over \(R\) by \(a_it\) for a finite set of generators \(a_i\) of \(I\). It is Noetherian by the Hilbert basis theorem. The module is finite over it, generated by any finite generating set of \(M\) in degree zero. Taking the quotient by \(I\) in degree zero gives the associated graded ring and module, denoted \(\operatorname{gr}_I R\) and \(\operatorname{gr}_I M\), with pieces \(I^n/I^{n+1}\) and \(I^nM/I^{n+1}M\). In particular the ring is standard graded over \(R/I\) and the module is finite over it.

\begin{theorem}[Artin--Rees]\label{p103:ALG-05-C04-R02}
If \(N\subseteq M\), there is \(c\ge0\) such that for every \(n\ge c\),
\[
 I^nM\cap N=I^{n-c}(I^cM\cap N).
\]
\end{theorem}
\begin{proof}
The direct sum \(\bigoplus_{n\ge0}(I^nM\cap N)t^n\) is a graded submodule of \(\mathcal R_I(M)\). The latter is Noetherian, so choose finitely many generators of the submodule and replace each by its finitely many homogeneous components. Let \(c\) bound their degrees, taking \(c=0\) for the zero submodule. If a generator has degree \(e\le c\), its contribution to degree \(n\ge c\) is a product of an element of \(I^{n-e}\) and a member of \(I^eM\cap N\). Factor \(I^{n-e}=I^{n-c}I^{c-e}\). The \(I^{c-e}\) multiple belongs to \(I^cM\cap N\), which gives one inclusion. The other holds because multiplication by \(I^{n-c}\) sends \(I^cM\cap N\) into \(I^nM\cap N\). This is \cite[Lemma 5.1]{eisenbud}.
\end{proof}

\begin{theorem}[Krull intersection]\label{p103:ALG-05-C04-R04}
Let \(N=\bigcap_{n\ge0}I^nM\) for a finite \(R\)-module over a Noetherian ring. Then \(IN=N\), and some \(a\in I\) satisfies \((1-a)N=0\). In particular \(N=0\) if \(I\subseteq\operatorname{Jac}R\). If \(R\) is a domain and \(I\) is proper, \(\bigcap_n I^n=(0)\).
\end{theorem}
\begin{proof}
Apply Artin--Rees to the submodule \(N\) with \(n=c+1\). Since \(N\subseteq I^nM\) for all \(n\), the result says \(N=IN\). This is a finite submodule of the Noetherian module \(M\). The determinant argument in Nakayama gives \((1-a)N=0\) with \(a\in I\) before using any Jacobson assumption. If \(I\subseteq\operatorname{Jac}R\) the factor is a unit. For \(M=R\) in a domain it is nonzero when \(I\ne R\), and a nonzero scalar kills no nonzero element. For arbitrary finite modules the last argument requires a torsion condition. For \(R=\mathbb Z\), \(I=(2)\) and \(M=\mathbb Z/3\), all powers of \(I\) act surjectively. See \cite[Corollary 5.4]{eisenbud}.
\end{proof}

\section{Adic completion}

For an \(R\)-module \(M\), write \(\widehat M=\varprojlim_{n\ge1} M/I^nM\) for the module of compatible residue classes. The ring \(\widehat R\) acts on \(\widehat M\) by multiplying the classes at each index. The natural map from \(M\) sends an element to its residues. The \(I\)-adic completion of an arbitrary module need not be a separated copy of the original module, since the kernel is the intersection of the powers. We use \cite[Sections 7.1, 7.2 and 7.5]{eisenbud}.

\begin{theorem}[Exactness and the tensor comparison]\label{p103:ALG-05-C04-R03}
Let \(R\) be Noetherian and \(I\) an ideal. Completion preserves short exact sequences of finite \(R\)-modules. For every such \(M\), the natural map
\[
 \widehat R\otimes_R M\ \longrightarrow\ \widehat M,\qquad
 r\otimes m\longmapsto r\,\widehat m
\]
is an isomorphism, natural in \(M\).
\end{theorem}
\begin{proof}
Represent a short exact sequence as \(0\to A\to B\to C\to0\) with \(A\subseteq B\). The exact residue sequences are
\[
 0\to A/(A\cap I^nB)\to B/I^nB\to C/I^nC\to0.
\]
Artin--Rees makes the filtration \(A\cap I^nB\) cofinal with \(I^nA\). One inclusion is \(I^nA\subseteq A\cap I^nB\), and for \(n\ge c\) the latter is contained in \(I^{n-c}A\). Explicitly, reducing at any desired index after passing to a sufficiently deep term of the other filtration gives inverse maps on the two inverse limits. Independence of the deeper index follows from compatibility.

The inverse limit of the displayed exact sequences is exact. The middle kernel consists of compatible classes in the levelwise left terms. For surjectivity, start with a lift of the first class in \(C/IC\). Given a lift at level \(n\), take any lift at level \(n+1\). Its reduction differs from the chosen lift by a class in \(A/(A\cap I^nB)\). Lift that class to \(A/(A\cap I^{n+1}B)\), whose reduction map is surjective, and subtract it. Induction gives compatible lifts. This proves exactness after using cofinality.

By Theorem~\ref{p103:ALG-05-C02-R01}, \(M\) has a finite presentation \(R^a\to R^b\to M\to0\). Completions of the finite free modules are \(\widehat R^a,\widehat R^b\), as compatibility is coordinatewise. Applying exactness to the kernel and image sequences of the presentation gives a right exact presentation of \(\widehat M\) by these completed free modules. Tensoring the original presentation with \(\widehat R\) gives the same presentation, since tensor product is right exact. The comparison on each free module is the usual basis isomorphism, so it gives the required isomorphism on cokernels. Its stated formula on pure tensors proves naturality.
\end{proof}

\begin{theorem}[The completed ring]\label{p103:ALG-05-C04-S01}
For Noetherian \(R\) and any ideal \(I\), \(\widehat R\) is Noetherian and flat over \(R\). If \(J=I\widehat R\), there are compatible identifications
\[
 \widehat R/J^n\cong R/I^n\quad(n\ge1),\qquad
 \operatorname{gr}_{J}\widehat R\cong\operatorname{gr}_I R.
\]
It is faithfully flat over \(R\) exactly when \(I\subseteq\operatorname{Jac}R\). If \(R\) is local and \(I\) is its maximal ideal, \(\widehat R\) is local with maximal ideal \(J\) and the same residue field.
\end{theorem}
\begin{proof}
Apply Theorem~\ref{p103:ALG-05-C04-R03} to \(I^n\to R\to R/I^n\). The completion of the last term is itself, since sufficiently high powers act by zero. The image of \(\widehat{I^n}\) in \(\widehat R\) is \(I^n\widehat R=J^n\), using the tensor comparison. This identifies the quotient maps with the coordinate projections of the original inverse limit. Their successive kernels give the graded isomorphism. They also show that \(\widehat R\) is separated and complete for \(J\), in the sense that its natural map to the inverse limit of its \(J\)-power quotients is an isomorphism.

For Noetherianness, the graded ring \(\operatorname{gr}_{J}\widehat R\) is Noetherian by its identification with the standard graded algebra \(\operatorname{gr}_I R\) over the Noetherian ring \(R/I\). For any ideal \(L\subseteq\widehat R\) its initial forms make a homogeneous ideal in the graded ring. Choose finitely many homogeneous generators and lift them to \(l_i\in L\cap J^{d_i}\) of their respective degrees \(d_i\). If \(x\in L\cap J^n\), express its degree-\(n\) class using those generators and lift the coefficients in \(J^{n-d_i}\) for \(d_i\le n\). Subtracting the corresponding sum of multiples of \(l_i\) leaves an element of \(L\cap J^{n+1}\). Starting in degree zero and repeating, the successive corrections to the coefficient of each \(l_i\) tend to zero \(J\)-adically because \(n-d_i\) tends to infinity. Completeness gives limits of the coefficient sums and separatedness gives equality of \(x\) with the resulting finite combination of \(l_i\). The zero initial ideal gives \(L=0\) by separatedness. Thus every ideal is finitely generated. This argument does not assume the ideal \(L\) is already closed.

For flatness, every ideal \(K\) of \(R\) is finite. Exactness of completion and its tensor description identify \(\widehat R\otimes_R K\to\widehat R\otimes_R R\) with an injection. The ideal criterion for flatness (Theorem~\ref{p103:ALG-03-C02-R04}) proves the assertion. Every member of \(J\) is in the Jacobson radical of \(\widehat R\). For \(a\in J\) and \(b\in\widehat R\), the geometric series in \(ab\) converges \(J\)-adically and inverts \(1-ab\). If a maximal ideal \(\mathfrak m\) of \(R\) contains \(I\), then \(\widehat R/\mathfrak m\widehat R\cong R/\mathfrak m\) by reduction modulo \(J\) followed by reduction modulo \(\mathfrak m/I\). If \(\mathfrak m\) does not contain \(I\), write \(1=a+b\) with \(a\in I\), \(b\in\mathfrak m\). The image of \(b\) is a unit, so the same fiber is zero. Thus the criterion that a flat module be faithfully flat exactly when all its maximal fibers are nonzero (Theorem~\ref{p103:ALG-03-C02-S04}) gives precisely \(I\subseteq\operatorname{Jac}R\). When \(R\) is local and \(I=\mathfrak m\), the ideal \(J\) is in the Jacobson radical and its quotient is a field, so it is the unique maximal ideal. The proof includes \(I=0\) and \(I=R\) by the inverse-limit definitions.
\end{proof}

\section{Hilbert--Samuel and regular local rings}

\begin{theorem}[The Hilbert--Samuel dimension formula]\label{p103:ALG-05-C03-S02}
Let \((R,\mathfrak m)\) be Noetherian local, \(M\ne0\) a finite module of dimension \(d\), and \(Q\subseteq\mathfrak m\) an ideal with \(M/QM\) of finite length. Then
\[
 L_Q(n)=\ell_R(M/Q^{n+1}M)
\]
agrees for large \(n\) with a polynomial of degree \(d\). The layer length \(\ell_R(Q^nM/Q^{n+1}M)\) is eventually polynomial of degree \(d-1\) for \(d>0\), and is eventually zero for \(d=0\).
\end{theorem}
\begin{proof}
Pass to \(R/\operatorname{Ann}M\), which preserves dimensions and all these lengths. Then \(\sqrt Q=\mathfrak m\) by the parameter theorem, and \(R/Q\) has finite length. The ring \(\operatorname{gr}_Q R\) is standard over \(R/Q\), with \(\operatorname{gr}_Q M\) a finite graded module. Its layer lengths agree whether measured over \(R\) or \(R/Q\). Lemma~\ref{p103:hilbert-serre-length} gives eventual polynomial behavior of the layer lengths. Summing the layers gives eventual polynomial behavior of \(L_Q\), for the sum of the values of a polynomial is polynomial in the upper index. Its eventual polynomial is nonzero because \(M/QM\ne0\). If it has positive degree its leading coefficient is positive because the lengths are nonnegative for all large \(n\).

There is \(s\ge1\) such that \(\mathfrak m^s\subseteq Q\subseteq\mathfrak m\). Consequently
\[
 L_{\mathfrak m}(n)\le L_Q(n)\le L_{\mathfrak m}(s(n+1)-1).
\]
Comparing nonnegative eventually polynomial functions in both directions shows that their degrees are equal. This allows us to use any finite-colength proper ideal to bound the degree. For a parameter ideal generated by \(d>0\) elements, each \(Q^jM/Q^{j+1}M\) is generated as an image of a direct sum of \(\binom{j+d-1}{d-1}\) copies of \(M/QM\), via multiplication by the degree-\(j\) monomials. Thus \(L_Q(n)=O((n+1)^d)\). For \(d=0\) the entire nonzero module has finite length and is killed by a power of \(\mathfrak m\), so \(L_Q\) is eventually constant of degree zero. We have proved the upper bound in all cases.

For the reverse bound use induction on \(d\). If \(d>0\), \(L_{\mathfrak m}\) cannot be eventually constant. Equal consecutive lengths would give \(\mathfrak m^nM=\mathfrak m^{n+1}M\), so Nakayama on the finite module \(\mathfrak m^nM\) would make it zero. This would force dimension zero. Write \(e\ge1\) for the degree of \(L_{\mathfrak m}\). Choose \(x\in\mathfrak m\) outside the associated primes of \(M\) other than \(\mathfrak m\), by finiteness and prime avoidance. All minimal primes in the support are among these nonmaximal primes when \(d>0\). As in the proof of the parameter theorem, avoiding them gives \(\dim(M/xM)\le d-1\). The opposite inequality is part of that theorem, so equality holds. The module \(K=0:_M x\) has finite length. Its associated primes are among those of \(M\) and contain \(x\), so are confined to \(\mathfrak m\), and all minimal primes in its support are associated. This makes its support at most the closed point, and the finite-length argument of the parameter theorem applies, including \(K=0\).

For large \(n\), Artin--Rees applied to \(xM\subseteq M\) supplies a fixed \(c\ge1\) with
\[
 xM\cap\mathfrak m^{n+1}M\subseteq x\mathfrak m^{n+1-c}M .
\]
In the length of \(M/\mathfrak m^{n+1}M\), the contribution of the image of \(xM\) is therefore at least \(\ell_R(xM/x\mathfrak m^{n+1-c}M)\). Multiplication by \(x\) identifies this quotient with
\(M/(\mathfrak m^{n+1-c}M+K)\). The natural surjection from
\(M/\mathfrak m^{n+1-c}M\) to that module has kernel a quotient
of \(K\). The target length is therefore at least
\(L_{\mathfrak m}(n-c)-\ell_R(K)\), and subtracting gives
\[
 \ell_R\!\left((M/xM)/\mathfrak m^{n+1}(M/xM)\right)
 \le L_{\mathfrak m}(n)-L_{\mathfrak m}(n-c)+\ell_R(K).
\]
The right side is \(O((n+1)^{e-1})\). The left side has degree \(d-1\) by induction, with positive leading coefficient. Thus \(e\ge d\), which finishes the degree calculation. Taking the first difference of the cumulative length now proves the assertion about layers. This is the length-comparison approach of \cite[Chapter 12, especially Theorem 12.4]{eisenbud}.
\end{proof}

\begin{theorem}[Basic structure of a regular local ring]\label{p103:ALG-05-C03-S03}
Let \((R,\mathfrak m,k)\) be a Noetherian local ring with residue field \(k\) and dimension \(d\). The minimum size of a set of generators of \(\mathfrak m\) is \(\dim_k\mathfrak m/\mathfrak m^2\), at least \(d\). The ring is regular local if this number is \(d\). In this case
\[
 \operatorname{gr}_{\mathfrak m}R\cong k[X_1,\ldots,X_d],
\]
via any minimal set \(x_1,\ldots,x_d\) of generators of \(\mathfrak m\). The ring \(R\) is a domain, and such generators form a regular sequence, meaning that each successive \(x_i\) acts injectively on \(R/(x_1,\ldots,x_{i-1})\) and the final ideal is proper. For \(d=0\) the ring is a field.
\end{theorem}
\begin{proof}
Nakayama applied to the finite module \(\mathfrak m\) shows that lifts of a \(k\)-basis of \(\mathfrak m/\mathfrak m^2\) generate \(\mathfrak m\), and every generating set maps to a spanning set. The height theorem gives the lower bound \(d\). Assume equality. If \(d=0\), Nakayama gives \(\mathfrak m=0\), so \(R=k\). In positive dimension the degree-one images of a minimal generating set give a graded surjection from \(k[X_1,\ldots,X_d]\) to \(\operatorname{gr}_{\mathfrak m}R\). If the kernel is nonzero, every prime containing it is nonzero and the quotient has dimension at most \(d-1\) by the polynomial-dimension argument. By Theorem~\ref{p103:ALG-05-C03-R04}, the eventual degree of its graded dimensions is then at most \(d-2\), allowing the eventually zero case in dimension zero. This contradicts the degree \(d-1\) of the layer lengths of \(R\) in Theorem~\ref{p103:ALG-05-C03-S02}, since each layer is a \(k\)-vector space. Thus the graded surjection is an isomorphism.

By Krull intersection, a nonzero element has finite order in the \(\mathfrak m\)-power filtration and a nonzero initial form. The product of the initial forms of two nonzero elements is nonzero in the graded polynomial domain. This is the corresponding graded class of their product, proving the product is nonzero in \(R\). Finally \(x_1\ne0\) for \(d>0\) because its class is a basis member. It is therefore a nonzero divisor. The local ring \(R/(x_1)\) has dimension at least \(d-1\) by the parameter theorem and its maximal ideal is generated by the images of the other \(d-1\) elements. The height theorem gives the matching upper bound, so it too is regular local with these images as a minimal generating set. Induction gives the regular-sequence assertion, with the final ideal \(\mathfrak m\) proper. Compare \cite[Section 10.3]{eisenbud}.
\end{proof}

\section{Hensel lifting}

\begin{theorem}[Hensel's lemma]\label{p103:ALG-05-C04-S02}
Let \(R\) be separated and complete for the powers of an ideal \(I\). If \(f\in R[T]\), \(a\in R\), and \(d=f'(a)\) satisfy \(f(a)\in d^2I\), there is \(b\in a+dI\) with \(f(b)=0\). This \(b\) is unique within \(a+dI\) when multiplication by \(d\) is injective. In particular a simple root of the reduction of \(f\) over the residue field of a complete local ring lifts uniquely in its residue class.

If \(f\) is monic and its reduction modulo \(I\) factors as a product of two relatively prime monic polynomials, with relatively prime meaning that the two generate the unit ideal of \((R/I)[T]\), the factors lift uniquely to monic factors of \(f\) with their prescribed reductions and degrees.
\end{theorem}
\begin{proof}
Write \(f(a)=d^2 c\) with \(c\in I\), and expand
\[
 f(a+d t)=d^2\bigl(c+t+t^2h(t)\bigr)
\]
for some \(h\in R[T]\). The expansion uses the polynomial identity that the terms of order at least two in a displacement are divisible by its square. It does not divide by \(d\). Set \(t_0=0\) and \(t_{j+1}=-c-t_j^2h(t_j)\). All terms lie in \(I\). If \(s,t\in I\) and \(s-t\in I^n\), then
\[
 s^2h(s)-t^2h(t)
 =(s-t)(s+t)h(s)+t^2(h(s)-h(t))\in I^{n+1}.
\]
The last membership uses the polynomial difference identity. The iteration therefore converges \(I\)-adically to a fixed point \(t\in I\), because it is eventually compatible modulo each power. This gives \(b=a+dt\). A pair of fixed points differ successively by every power of \(I\), so are equal by separatedness. If \(d\) acts injectively, any root in \(a+dI\) gives a fixed point by cancellation of \(d^2\), establishing uniqueness. If \(d=0\), the assumptions already give \(f(a)=0\) and \(a+dI=\{a\}\). For a simple root in the residue field take any representative \(a\). The derivative is a unit in the local ring and \(f(a)\in I=d^2I\) for the maximal ideal \(I\), so there is a unique root in the chosen class modulo \(I\).

For the factorization, write the given residual factors as \(\bar g,\bar h\) of degrees \(r,s\). Over \(R/I\) the linear map
\[
 (\alpha,\beta)\longmapsto\alpha\bar h+\beta\bar g
 \quad(\deg\alpha<r,\ \deg\beta<s)
\]
is an isomorphism onto the polynomials of degree \(<r+s\). Given a polynomial in the target, reduce it modulo the monic \(\bar g\) and multiply by the inverse of \(\bar h\) there, supplied by the Bézout relation, to get the unique remainder \(\alpha\) of degree \(<r\). The difference is divisible by \(\bar g\), and monic division gives the unique \(\beta\) of degree \(<s\). This works also when a factor is \(1\), with no correction on that factor.

Start with arbitrary monic lifts of the fixed degrees. If \(f-gh\) has coefficients in \(I^n\), it has degree less than \(r+s\). Reduce its coefficients in \(I^n/I^{n+1}\), an \(R/I\)-module. The linear isomorphism just proved remains an isomorphism on coefficient vectors with values in this module, by tensoring an isomorphism of finite free modules. It gives unique correction classes for \(g,h\) to cancel the error modulo \(I^{n+1}\). Lift the corrections with coefficients in \(I^n\) and add them to the factors. Their product of corrections lies in \(I^{2n}\subseteq I^{n+1}\) for \(n\ge1\), and degrees and leading coefficients are unchanged. Passing to coefficientwise limits gives a factorization. Uniqueness follows by applying the same correction isomorphism to the differences of two solutions, inductively modulo every \(I^n\), and then using separatedness. The root and factor forms are consistent with \cite[Theorem 7.3 and Exercises 7.20--7.22]{eisenbud}.
\end{proof}
\chapter[Finite-group representations]{Finite-group representations and characters}
\label{p103:ch08}

We write \(A_G=k[G]\) for the group algebra over a field \(k\).
The notation \(\operatorname{Hom}_G(V,W)\) denotes the \(k\)-linear
intertwiners. Unless specified otherwise, representations are left
representations, and an irreducible representation is nonzero. The
linear algebra of representations starts over arbitrary fields.
Hermitian inner products and ordinary characters below use
\(\mathbb C\). The condition on the characteristic in Maschke's
theorem is stated independently. See \cite[Chapters 1--3 and 6]{serre}
and \cite[Chapters 17 and 21]{d07}.

\section{Group algebras, irreducibility, and averaging}

\begin{proposition}[Representations and group algebra modules]
\label{p103:ALG-07-C01-R01}
For any group \(G\) and field \(k\), a representation on a \(k\)-vector
space \(V\) is a group homomorphism
\(\rho\colon G\to\operatorname{GL}_k(V)\). It is equivalent to a left
\(k[G]\)-module structure on \(V\) extending the given \(k\)-structure.
Equivariant linear maps are exactly the \(k[G]\)-module maps.
Invariant subspaces, kernels and images of intertwiners, and quotient
representations have the corresponding module descriptions.
\end{proposition}

\begin{proof}
The elements of \(k[G]\) are formal sums \(\sum_{g} a_g g\) with finite
support. Multiplication is the bilinear extension of \((g,h)\mapsto gh\).
The action on \(V\) is \((\sum_g a_g g)v=\sum_g a_g\rho(g)v\).
Associativity follows by comparing the two finite double sums for the
action of a product, and the identity of \(G\) acts as the identity.
Conversely, a \(k[G]\)-action gives invertible operators at \(g\) with
inverse given by the action of \(g^{-1}\). The two constructions
recover one another since the \(g\)'s span the group algebra. For a
linear map, commuting with every \(g\) is equivalent by linearity to
commuting with the whole algebra. The module descriptions now follow
from closure under the basis elements \(g\). In particular, the
quotient action is \(g(v+U)=gv+U\), independent of the representative
because \(U\) is invariant.
\end{proof}

\begin{theorem}[Schur's lemma]
\label{p103:ALG-07-C01-R02}
A nonzero intertwiner between irreducible representations is an
isomorphism, and the intertwining endomorphisms of an irreducible
representation form a division ring. If \(V\) is irreducible,
finite-dimensional and \(k\) is algebraically closed, then each
intertwining endomorphism of \(V\) is a unique scalar multiple of the
identity.
\end{theorem}

\begin{proof}
The kernel of a nonzero map cannot be all of its simple source, and
the nonzero image must be all of its simple target. Thus the map
is bijective and its inverse is an intertwiner. On a finite-dimensional
space over an algebraically closed field an endomorphism \(T\) has an
eigenvalue \(\lambda\). The intertwiner \(T-\lambda 1\) has a nonzero
kernel, so it is zero. Evaluating on a nonzero vector gives uniqueness.
\end{proof}

\begin{theorem}[Maschke's theorem and the modular obstruction]
\label{p103:ALG-07-C01-R03}
Let \(G\) be finite and let \(n=|G|\). If \(n\cdot 1_k\ne 0\), every
invariant subspace of a \(k\)-representation has an invariant
complement. Every representation is then semisimple, with a finite
decomposition into irreducibles when it is finite-dimensional.
If \(n\cdot 1_k=0\), the augmentation \(k[G]\to k\) onto the trivial
representation does not split.
\end{theorem}

\begin{proof}
Extend a basis of an invariant subspace \(U\subseteq V\) to a basis of
\(V\), and choose a linear projection \(P\) onto \(U\). Form
\[
P_G=\frac{1}{n}\sum_{g\in G}\rho(g)P\rho(g^{-1}).
\]
Each summand has image in \(U\) and restricts to the identity there.
Hence \(P_G^2=P_G\) and its image is \(U\). Conjugating by \(\rho(h)\)
permutes the summands, so \(\ker P_G\) is invariant. In finite dimension,
repeatedly splitting a proper submodule yields a sum of irreducibles.
In arbitrary dimension the submodule complement characterization
of semisimplicity in \ref{p103:ALG-03-C03-R03} applies.
For the second assertion, in the left regular representation invariance
of a vector \(\sum_g a_g g\) forces all the \(a_g\)'s to be equal.
Its augmentation is therefore a multiple of \(n\cdot1_k=0\).
The image of \(1\) under an equivariant section would be invariant
with augmentation \(1\), a contradiction.
\end{proof}

\section{Traces, duals, and character pairings}

For a finite-dimensional representation, its character is
\(\chi_V(g)=\operatorname{tr}(\rho(g))\). This is constant on conjugacy
classes by invariance of trace under similarity.

\begin{proposition}[Characters of linear constructions]
\label{p103:ALG-07-C01-S01}
Over any field, the finite-dimensional dual \(V^*\) has action
\((g\lambda)(v)=\lambda(g^{-1}v)\), and
\[
\chi_{V^*}(g)=\chi_V(g^{-1}),\qquad
\chi_{V\oplus W}=\chi_V+\chi_W,\qquad
\chi_{V\otimes_k W}=\chi_V\chi_W.
\]
Here the tensor product has diagonal action. If \(2\) is invertible in
\(k\), the eigenspaces of the flip on \(V\otimes_k V\) for \(+1\) and
\(-1\) identify equivariantly with \(\operatorname{Sym}^2 V\) and
\(\bigwedge^2 V\), and their characters are
\[
\chi_{\operatorname{Sym}^2 V}(g)
   =\tfrac12\bigl(\chi_V(g)^2+\chi_V(g^2)\bigr),\qquad
\chi_{\bigwedge^2 V}(g)
   =\tfrac12\bigl(\chi_V(g)^2-\chi_V(g^2)\bigr).
\]
\end{proposition}

\begin{proof}
In dual bases the operator on \(V^*\) is the transpose of the operator
of \(g^{-1}\). The trace on a direct sum adds, and in tensor bases the
diagonal entries of \(A\otimes B\) are the pairwise products of the
diagonal entries of \(A,B\). This proves the three trace identities.

Let \(\tau(v\otimes w)=w\otimes v\). The projections
\((1+\tau)/2,(1-\tau)/2\) commute with the group action. The
symmetric-square quotient kills the image of \(1-\tau\).
The exterior-square quotient kills the span of all \(v\otimes v\),
which also contains \(v\otimes w+w\otimes v\) by expanding
\((v+w)\otimes(v+w)\). Thus its kernel is the \(+1\) eigenspace
when \(2\) is invertible. This proves the asserted identifications.
For \(A=\rho(g)\), the trace of \(A\otimes A\) on either eigenspace
is the full trace after multiplication by the corresponding
projection. In a tensor basis,
\(\operatorname{tr}(\tau(A\otimes A))=\sum_{i,j} A_{ij}A_{ji}
=\operatorname{tr}(A^2)\).
Since \(A^2=\rho(g^2)\), the formulas follow.
\end{proof}

\begin{theorem}[Unitarity and the ordinary character pairing]
\label{p103:ALG-07-C02-R01}
Let \(G\) be finite and let \(V,W\) be finite-dimensional complex
representations. Each representation admits a \(G\)-invariant
positive-definite Hermitian inner product, and
\(\chi_V(g^{-1})=\overline{\chi_V(g)}\). On complex functions on \(G\)
use the inner product linear in its first variable
\[
 \langle f,h\rangle_G=\frac1{|G|}\sum_{g\in G}f(g)\overline{h(g)}.
\]
Then
\[
 \langle\chi_V,\chi_W\rangle_G
 =\dim_{\mathbb C}\operatorname{Hom}_G(W,V)
 =\dim_{\mathbb C}\operatorname{Hom}_G(V,W).
\]
\end{theorem}

\begin{proof}
The average of \(h_0(gv,gw)\) over \(G\), where \(h_0\) is any
positive-definite Hermitian inner product, remains positive-definite
and is invariant by the substitution \(g\mapsto gh\). In an orthonormal
basis the matrix of \(g^{-1}\) is the conjugate transpose of that
of \(g\). On \(\operatorname{Hom}_{\mathbb C}(W,V)\) the action is
\(T\mapsto \rho_V(g)T\rho_W(g^{-1})\). Its invariants are the
intertwiners, and the group average is a projection onto them.
In the basis of rectangular matrix units, the trace of this action is
\(\operatorname{tr}(\rho_V(g))\operatorname{tr}(\rho_W(g^{-1}))\).
The trace of the projection is its rank, yielding the first equality.
Interchanging the inputs conjugates the pairing, whose value is a real
integer. This yields the second equality with the stated orientation.
\end{proof}

\section{The regular representation and Fourier inversion}

\begin{theorem}[Regular decomposition and multiplicity]
\label{p103:ALG-07-C02-R03}
Let \(G\) be finite. There are finitely many isomorphism classes
\(V_i\) of irreducible complex representations. With
\(d_i=\dim_{\mathbb C}V_i\) and \(A=\mathbb C[G]\), there is an
isomorphism of left \(A\)-modules
\[
 A\simeq \bigoplus_i V_i^{\oplus d_i},
 \qquad |G|=\sum_i d_i^2.
\]
The multiplicity of \(V_i\) in a finite-dimensional representation
\(V\) is \(\langle\chi_V,\chi_i\rangle_G\), where
\(\chi_i=\chi_{V_i}\). In particular, the ordinary character determines
\(V\) up to isomorphism.
\end{theorem}

\begin{proof}
The map \(a\mapsto av\) from \(A\) onto an irreducible representation
is surjective for any \(v\ne0\) because its image is a nonzero
subrepresentation. It splits by Maschke's theorem. Every irreducible
therefore occurs in the finite decomposition of \(A\). This also
proves that such irreducibles are finite-dimensional. Evaluating at
\(1\in A\) identifies \(\operatorname{Hom}_A(A,V_i)\) with \(V_i\).
Schur's lemma shows that each copy of \(V_i\) contributes one to the
dimension of this Hom space and all other copies contribute zero.
The multiplicity in \(A\) is \(d_i\), and dimensions give the
sum of squares. For general finite-dimensional \(V\), the character
pairing with \(\chi_i\) counts the summands of type \(V_i\), again by
Schur's lemma. Semisimplicity then gives the classification.
\end{proof}

\begin{theorem}[Matrix coefficient orthogonality and group Fourier transform]
\label{p103:ALG-07-C02-S02}
Choose an orthonormal basis for each unitary representative \(V_i\) and
write \(u^{(i)}_{ab}(g)=\rho_i(g)_{ab}\). Then
\[
 \left\langle u^{(i)}_{ab},u^{(j)}_{cd}\right\rangle_G
     =\frac{1}{d_i}\delta_{ij}\delta_{ac}\delta_{bd}.
\]
The \(\sqrt{d_i}\,u^{(i)}_{ab}\) are an orthonormal basis of
\(\mathbb C^G\). The action map
\[
 \mathbb C[G]\longrightarrow\prod_i\operatorname{End}_{\mathbb C}(V_i),
 \qquad a\longmapsto \widehat a=(\rho_i(a))_i
\]
is an algebra isomorphism. Its inverse has coefficients
\[
 a_g=\frac1{|G|}\sum_i d_i
       \operatorname{tr}\bigl(\rho_i(g^{-1})\widehat a_i\bigr)
       \qquad\text{for }a=\sum_g a_g g.
\]
If \(|G|\) is invertible in a field \(k\) and the group algebra splits
over \(k\) (all division-ring blocks of its semisimple decomposition
are \(k\)), the action map gives the analogous matrix algebra
decomposition over \(k\).
\end{theorem}

\begin{proof}
For a linear map \(B\colon V_j\to V_i\), average the maps
\(\rho_i(g)B\rho_j(g^{-1})\). The result intertwines. It is zero if
\(i\ne j\), and if \(i=j\) it is
\(\operatorname{tr}(B)1/d_i\) by Schur's lemma and invariance of trace.
Take \(B\) to be the matrix unit sending the \(d\)-th basis vector to
the \(b\)-th. The \((a,c)\)-entry of the average is the desired
inner product, giving the formula. The number of orthogonal
nonzero functions is \(\sum_i d_i^2=|G|\), which proves completeness.

If \(a\) acts by zero in each \(V_i\), it acts by zero on \(A\)
by its regular decomposition and hence \(a=a\cdot1=0\).
The action map is injective and the two dimensions are equal.
For \(g\in G\), left multiplication by \(g^{-1}a\) has constant
diagonal entry \(a_g\) in the group basis. The coefficient of \(x\)
in \((g^{-1}a)x\) is \(a_g\). Its trace is \(|G|a_g\). Computing
the same trace in the regular decomposition gives the inverse
formula, including the factor \(d_i\).

Over a split field, Maschke's theorem gives complete reducibility.
The Wedderburn decomposition in \ref{p103:ALG-03-C03-R04}
identifies the Schur endomorphism rings with \(k\). The same
Hom computation gives multiplicities \(d_i\) in \(k[G]\), so the same
injectivity and dimension argument proves the algebra decomposition.
The normalized Hermitian coefficient calculation above uses the
complex field.
\end{proof}

\begin{theorem}[Rows and columns of the character table]
\label{p103:ALG-07-C02-R02}
For a finite group \(G\), the irreducible complex characters form an
orthonormal basis of the complex class functions. Their number
equals the number of conjugacy classes. For \(g,h\in G\), column
orthogonality is
\[
 \sum_i\chi_i(g)\overline{\chi_i(h)}
    = \begin{cases}
        |C_G(g)|,& \text{if }g,h\text{ are conjugate},\\
        0,&\text{otherwise},
      \end{cases}
\]
where \(C_G(g)=\{x\in G:xg=gx\}\).
\end{theorem}

\begin{proof}
Row orthogonality follows from the character pairing and Schur's
lemma. Under the Fourier algebra isomorphism, the center has one
scalar coordinate for each irreducible, since the center of a full
matrix algebra consists of scalar matrices. In the group basis,
centrality says exactly that the coefficients are constant on
conjugacy classes. The class sums are a basis of the center, so its
dimension is the number of classes. Thus the orthonormal characters
are a basis of the class functions. Let \(D\) be the conjugacy class
of \(h\) and expand its indicator in this basis. The coefficient
of \(\chi_i\) is \(|D|\overline{\chi_i(h)}/|G|\).
Evaluating at \(g\) and multiplying by \(|G|/|D|=|C_G(h)|\)
gives the stated formula.
\end{proof}

\begin{theorem}[Isotypic projections and central idempotents]
\label{p103:ALG-07-C02-S01}
The sum of all submodules of a complex \(G\)-representation isomorphic
to \(V_i\) is the \(i\)-th isotypic summand. It is independent of the
chosen decomposition, and the projection onto it is the action of
\[
 e_i=\frac{d_i}{|G|}\sum_{g\in G}\chi_i(g^{-1})g\in\mathbb C[G].
\]
The \(e_i\) are precisely the primitive central idempotents. They are
pairwise orthogonal and sum to \(1\), and the class sums are another
basis of the center.
\end{theorem}

\begin{proof}
The coefficients of \(e_i\) are conjugacy invariant, so \(e_i\) is
central. For the tuple having identity in position \(i\) and zero
elsewhere, the inverse Fourier formula gives exactly the displayed
coefficients. These tuples are pairwise orthogonal idempotents summing
to the unit, and are precisely the primitive idempotents in the
center \(\prod_i\mathbb C\). On any direct sum of irreducibles,
\(e_i\) acts as the identity on summands of its type and as zero on
the others. This proves that its image is the sum of all irreducible
submodules of the stated type, regardless of the chosen decomposition.
The class-sum assertion follows from centrality in the group basis.
\end{proof}

\begin{proposition}[Abelian groups and products of groups]
\label{p103:ALG-07-C01-S02}
For a finite abelian group over a splitting field of characteristic
prime to its order, every irreducible representation is
one-dimensional. In particular the conclusion holds if the field
contains all roots of \(X^m-1\), where \(m\) is the exponent of the group
and the characteristic is prime to the order. For finite groups
\(G,K\), every irreducible complex representation of \(G\times K\)
is \(V_i\boxtimes W_j\) for unique irreducible types of \(G\) and \(K\).
The action on this outer tensor product is
\((g,k)\mapsto\rho_i(g)\otimes\sigma_j(k)\). When \(G=K\),
restricting to the diagonal recovers the ordinary tensor product.
\end{proposition}

\begin{proof}
The algebra of an abelian group is commutative, so each matrix
block in the split decomposition is of size \(1\). Under the given
root condition every group operator is diagonalizable because its
minimal polynomial divides the square-free split polynomial \(X^m-1\).
Commuting operators preserve each other's eigenspaces, so successively
choosing a nonzero eigenspace for each of the finitely many operators
produces a common eigenvector. Its span must be the irreducible module.
For the product, the character of \(V_i\boxtimes W_j\) is
\((g,k)\mapsto\chi_i(g)\psi_j(k)\). The character pairing factors, so
its square norm is \(1\), and the pairing between different pairs is
zero. In a semisimple representation the square norm is the sum of
squares of multiplicities, so the representations found are
irreducible and pairwise inequivalent. Their squared dimensions sum to
\((\sum_i d_i^2)(\sum_j \dim(W_j)^2)=|G||K|\).
The regular sum of squares for \(G\times K\) leaves no further type.
\end{proof}

\begin{example}[A character table]
The trivial and sign representations of \(S_3\) have rows
\((1,1,1)\) and \((1,-1,1)\) on the classes of the identity,
a transposition, and a three-cycle. The subspace
\(x_1+x_2+x_3=0\) of the permutation representation has row
\((2,0,-1)\) because the full trace counts fixed points.
Its squared norm is \((4+3\cdot0+2\cdot1)/6=1\). Thus it is
irreducible, and the sum \(1^2+1^2+2^2=6\) proves completeness.
\end{example}

\section{Arithmetic of ordinary characters}

\begin{theorem}[Integrality and divisibility of degrees]
\label{p103:ALG-07-C02-S03}
Every ordinary character value of a finite group is an algebraic
integer. For an irreducible \(V\) of degree \(d\), the scalar by
which a class sum acts on \(V\) is an algebraic integer. Moreover,
\[
 d\mid |G| \qquad\text{and}\qquad d\mid [G:Z(G)] .
\]
\end{theorem}

\begin{proof}
Each \(\rho(g)\) has finite order. Its eigenvalues are roots of unity,
so their sum is an algebraic integer. We use here the elementary
fact that sums and products of algebraic integers are algebraic
integers. Indeed, if \(\alpha,\beta\) satisfy monic integer polynomials,
\(\mathbb Z[\alpha,\beta]\) is generated as an abelian group by
bounded powers of \(\alpha,\beta\). It is a finitely generated
torsion-free abelian group, hence free. The characteristic polynomial
of multiplication by any of its elements annihilates that element
by evaluating the Cayley--Hamilton identity at \(1\).

For a conjugacy class \(D\), let \(z_D=\sum_{g\in D}g\).
The center of \(\mathbb Z[G]\) is free over \(\mathbb Z\) on the
\(z_D\)'s. The product of two class sums is central and integral
coefficientwise. Its coefficient at \(x\) counts the pairs of
factors with product \(x\), a count invariant under conjugation.
Thus multiplication by \(z_D\) is an integer matrix on this lattice.
Cayley--Hamilton, applied to \(1\), gives a monic integer relation
for \(z_D\). By Schur's lemma, the action of a central element on
\(V\) is scalar. The central scalar of \(z_D\) is
\(\omega_D=|D|\chi_V(g_D)/d\) for any \(g_D\in D\), by taking the
trace. Applying the representation to the monic relation shows
that \(\omega_D\) is an algebraic integer. Now row orthogonality gives
\[
 \frac{|G|}{d}
   =\sum_D \frac{|D|\chi_V(g_D)}{d}\,\chi_V(g_D^{-1})
   =\sum_D\omega_D\,\chi_V(g_D^{-1}).
\]
A rational algebraic integer is an integer. For a reduced fraction
\(a/b\) in a monic integral equation, clearing denominators shows
\(b\mid a^m\) for the degree \(m\) of that equation. Hence \(d\mid |G|\).

Let \(Z=Z(G)\), let \(q=[G:Z]\), and let
\(\lambda\colon Z\to\mathbb C^\times\) be the scalar by which \(Z\)
acts on \(V\). For each \(m\ge1\), the outer tensor power
\(V^{\boxtimes m}\) is irreducible for \(G^m\) by the product result.
The central subgroup
\[
 N_m=\{(z_1,\ldots,z_m)\in Z^m:z_1\cdots z_m=1\}
\]
has size \(|Z|^{m-1}\) and acts trivially since the product of the
\(\lambda(z_i)\)'s is \(1\). The descended representation of \(G^m/N_m\)
remains irreducible and has degree \(d^m\). The divisibility already
proved gives \(d^m\mid |Z|q^m\). If a prime had larger exponent in
\(d\) than in \(q\), that inequality of exponents would fail for
large \(m\). Therefore \(d\mid q\).
\end{proof}

\section{Frobenius--Schur indicators and real forms}

\begin{theorem}[Frobenius--Schur and realization over \(\mathbb R\)]
\label{p103:ALG-07-C02-S04}
Let \(V\) be an irreducible finite-dimensional complex representation
of a finite group \(G\), and set
\[
 \nu(V)=\frac{1}{|G|}\sum_{g\in G}\chi_V(g^2).
\]
The value is \(0\) when \(V\not\simeq V^*\), \(+1\) when \(V\) has a
nonzero invariant symmetric bilinear form, and \(-1\) when it has a
nonzero invariant alternating bilinear form. In either nonzero case
the form is nondegenerate. The \(+1\) case is exactly the case of a
real form, meaning a real \(G\)-representation whose complexification
is isomorphic to \(V\). The \(-1\) case has real character and has no
real form.

More generally, a finite-dimensional complex representation has a
real character if and only if its multiplicities agree within each
pair of non-self-dual irreducibles \(V_i,V_i^*\). It has a real form
if and only if it satisfies that condition and each irreducible
of indicator \(-1\) has even multiplicity. Equivalently, a real
form exists exactly when there is a nondegenerate invariant
symmetric complex bilinear form.
\end{theorem}

\begin{proof}
Invariant bilinear forms correspond to equivariant maps \(V\to V^*\).
The dual of an irreducible is irreducible. The annihilator in \(V\)
of a nonzero proper invariant subspace of \(V^*\) would be a
nonzero proper invariant subspace of \(V\). By Schur's lemma the
space of forms is zero or one-dimensional, and any nonzero form
is nondegenerate. Transposing is an involution on this space, so
in the one-dimensional case a nonzero form has a unique symmetry
sign \(\varepsilon=+1\) or \(-1\).

Averaging characters computes invariant dimensions. The character
difference of the symmetric and exterior squares is
\(g\mapsto\chi_V(g^2)\). The invariant dimensions in their duals are
the dimensions of invariant symmetric and alternating bilinear
forms, respectively. For any complex representation the trivial
multiplicity is the same in its dual, for example by semisimple
decomposition. Thus \(\nu(V)\) is precisely the difference of
those dimensions and takes the asserted values.

For the assertion about real forms, choose an invariant Hermitian
inner product \(h\), linear in the first variable. Write
\(B(v,w)=h(v,Jw)\) for the nonzero invariant bilinear form.
Existence and uniqueness of \(Jw\) follow from nondegeneracy of
\(h\). The convention for \(h\) makes \(J\) conjugate-linear.
Invariance of \(h,B\) gives \(Jg=gJ\), and nondegeneracy of \(B\)
gives invertibility. Hence \(J^2=\lambda 1\) by Schur's lemma.
Commuting \(J^2\) with \(J\) gives
\(\lambda=\overline\lambda\), since \(J\) is conjugate-linear and
invertible. The symmetry sign satisfies
\[
 h(Jv,Jv)=B(Jv,v)=\varepsilon B(v,Jv)
   =\varepsilon h(v,J^2v)=\varepsilon\lambda h(v,v).
\]
Thus \(\varepsilon\lambda>0\). Scaling \(J\) by the positive
number \(|\lambda|^{-1/2}\) gives \(J^2=\varepsilon 1\).
When \(\varepsilon=+1\), the fixed space \(U\) of \(J\) is
\(G\)-stable and real-linear. Each \(v\) decomposes uniquely as
\(u+iw\) with \(u,w\in U\), since
\(u=(v+Jv)/2\) and \(w=(v-Jv)/(2i)\) are fixed. Thus \(U\) is a
real form. Conversely a real form admits an invariant positive
real inner product by averaging, and complex bilinear extension
gives an invariant nondegenerate symmetric form. A self-dual
representation has real character because the character of
the dual is the conjugate. This proves the claims for irreducibles.

For arbitrary \(V\), equality with the dual character is equivalent
to the stated equality of multiplicities, by character
classification. A \(G\)-invariant bilinear form pairs different
isotypic components trivially unless the two types are dual, by
applying Schur's lemma to the maps into duals. In a non-self-dual
pair the two isotypic components must have equal multiplicities to
admit a perfect pairing. In a self-dual component
\(V_i\otimes\mathbb C^{m_i}\), any invariant bilinear form is
\(B_i\otimes F_i\), since the invariant pairing between any two
copies is a scalar times the fixed form \(B_i\). It is nondegenerate
and symmetric exactly when \(F_i\) is nondegenerate and has the
same symmetry sign as \(B_i\). A nondegenerate alternating form
exists on \(\mathbb C^{m_i}\) precisely when \(m_i\) is even.
An odd-size skew-symmetric matrix has determinant zero, and
the two-dimensional standard skew form gives all even cases by
direct sum. This proves necessity of the conditions on a real form.

To construct the real form under those conditions, use the
conjugation \(J\) just constructed on each positive-type copy.
For a non-self-dual pair, unitarity identifies \(V_i^*\) with the
conjugate vector space \(\overline{V_i}\) equivariantly by
\(\bar v\mapsto h(-,v)\). Swap the two conjugate summands by
\((v,\bar w)\mapsto(w,\bar v)\). For a pair of negative-type copies,
use the involution \((v,w)\mapsto(Jw,-Jv)\), where now \(J^2=-1\).
All these maps are conjugate-linear equivariant involutions.
Their direct sum yields a real form by taking fixed vectors, as
above. Together with the bilinear-form analysis this proves the
general criterion. Compare \cite[Section 13.2, Theorem 31 and
Proposition 39]{serre}.
\end{proof}
\chapter[Induction and Mackey theory]{Induction, coinduction, and Mackey theory}
\label{p103:ch09}

The field is \(k\), and all group representations are left
representations. In the double coset \(KsH\), \(K\) acts on the left
and \(H\) on the right. Throughout, \(G/H\) denotes the set of
left cosets \(gH\), without any normality assumption.
The algebraic constructions of induction and coinduction make sense
for arbitrary groups. Finiteness is used explicitly when passing
between products and sums, and when taking ordinary characters.
See \cite[Sections 3.3, 7.1--7.4, 8.1 and 9.1--9.3]{serre}
and \cite[Chapter 21]{d07}.

\section{Induction, function models, and adjunctions}

\begin{theorem}[Induction and restriction]
\label{p103:ALG-07-C03-R01}
Let \(H\le G\) and let \(V\) be an \(H\)-representation. Restriction
\(\operatorname{Res}_H^G\) is given by the same action on the subgroup.
Induction is given by
\[
 \operatorname{Ind}_H^G V=k[G]\otimes_{k[H]}V
\]
with \(G\) acting by left multiplication on \(k[G]\).
If \(T\) is a set of representatives for \(G/H\), the direct sum
\(\bigoplus_{t\in T}V\) maps isomorphically onto the underlying
vector space by \(v\mapsto t\otimes v\) on the \(t\)-summand.
If \([G:H]\) and \(\dim_k V\) are finite, the dimension is
\([G:H]\dim_k V\). An equivariant function model consists of the
functions \(f\colon G\to V\) supported in finitely many cosets
\(xH\) and satisfying \(f(xh)=h^{-1}f(x)\) for \(x\in G,h\in H\).
The action is \((gf)(x)=f(g^{-1}x)\). This model is independent of
the representatives.
\end{theorem}

\begin{proof}
Restriction is functorial since an intertwiner for the larger
group remains one on the subgroup. The group algebra \(k[G]\)
is a \((k[G],k[H])\)-bimodule, and the action on the tensor product
respects balancing by associativity. The group basis is partitioned
by the cosets \(tH\), giving \(k[G]=\bigoplus_{t\in T}t k[H]\)
as right \(k[H]\)-modules.
Tensoring this identity, using the balanced tensor construction
in \ref{p103:ALG-03-C02-S01}, gives the stated direct sum.
A group element \(g\) sends \(t\otimes v\) to \(t'\otimes hv\)
if \(gt=t'h\). Under a replacement \(t\mapsto th\), the same
tensor \(t\otimes v\) is \(th\otimes h^{-1}v\), so the
coordinates transform accordingly.

Send \(t\otimes v\), with \(t\in G\), to the function supported
on \(tH\) that takes value \(h^{-1}v\) at \(th\).
The assignment for \(th_0\otimes v\) equals the assignment for
\(t\otimes h_0v\). At \(th\) the former value is
\((h_0^{-1}h)^{-1}v=h^{-1}h_0v\). Thus the map descends to the
balanced tensor product. It is equivariant by left translation.
Its inverse is \(f\mapsto\sum_{t\in T} t\otimes f(t)\), a finite
sum by the support condition. Replacing \(t\) by \(th\) replaces
\(f(t)\) by \(h^{-1}f(t)\), so the sum is unchanged.
\end{proof}

\begin{theorem}[Frobenius reciprocity]
\label{p103:ALG-07-C03-R02}
For \(H\le G\), any \(H\)-representation \(V\) and any
\(G\)-representation \(W\), there is a \(k\)-linear natural isomorphism
\[
 \operatorname{Hom}_G(\operatorname{Ind}_H^G V,W)
   \simeq \operatorname{Hom}_H(V,\operatorname{Res}_H^G W).
\]
Thus induction is left adjoint to restriction. For finite groups
and complex finite-dimensional representations, this gives
\[
 \langle\chi_{\operatorname{Ind}V},\chi_W\rangle_G
    =\langle\chi_V,\chi_{\operatorname{Res}W}\rangle_H.
\]
The same pairing equality extends to ordinary complex class functions
by linearity and conjugate-linearity. For irreducible \(V,W\), it
identifies the multiplicity of \(W\) in \(\operatorname{Ind}V\) with
the multiplicity of \(V\) in \(\operatorname{Res}W\).
\end{theorem}

\begin{proof}
An intertwiner \(F\) on the left is sent to \(v\mapsto F(1\otimes v)\).
This map is \(H\)-equivariant because
\(1\otimes hv=h\otimes v=h(1\otimes v)\). Conversely an
\(H\)-intertwiner \(\phi\) defines
\(F(g\otimes v)=g\phi(v)\). It is balanced since
\(gh\phi(v)=g\phi(hv)\), and \(G\)-equivariant by the formula.
Both composites are identities, and pre- or postcomposition
commutes with the construction, giving naturality.

The complex character pairing counts intertwiner dimensions in
either order by \ref{p103:ALG-07-C02-R01}. Applying it to this
isomorphism proves the equality with the displayed orientation.
Irreducible characters are bases of class functions, so the
pairing equality extends as stated, with induction on class
functions given by the formula in \ref{p103:ALG-07-C03-S02}.
The multiplicity assertion is then the character multiplicity formula.
\end{proof}

\begin{theorem}[Coinduction and finite-index comparison]
\label{p103:ALG-07-C03-S01}
For \(H\le G\), the coinduced representation of \(V\) is
\[
 \operatorname{Coind}_H^G V
  =\operatorname{Hom}_{k[H]}({}_{k[H]}k[G],V),
 \qquad (gF)(a)=F(ag).
\]
It is right adjoint to restriction, as expressed by the natural isomorphism
\[
 \operatorname{Hom}_G(W,\operatorname{Coind}_H^G V)
 \simeq \operatorname{Hom}_H(\operatorname{Res}_H^G W,V).
\]
Induction identifies naturally with the subrepresentation of
coinduction supported in finitely many cosets \(Hx\), and the
identification is onto when \([G:H]\) is finite.
\end{theorem}

\begin{proof}
Right multiplication on \(k[G]\) commutes with left multiplication
by \(k[H]\). The displayed action has \(g_1(g_2F)(a)=F(ag_1g_2)\)
and unit the identity. Restriction to group elements identifies
coinduction with functions \(F:G\to V\) satisfying
\(F(hx)=hF(x)\). The action is \((gF)(x)=F(xg)\).
A \(G\)-map \(\Phi:W\to\operatorname{Coind} V\) gives
\(\psi(w)=\Phi(w)(1)\). For \(h\in H\),
\(\psi(hw)=\Phi(w)(h)=h\psi(w)\). Conversely, from an \(H\)-map
\(\psi\) take \(\Phi(w)(x)=\psi(xw)\). It satisfies the function
covariance and the \(G\)-map identity. Evaluation gives one
inverse, and \(G\)-equivariance of \(\Phi\) gives the other
because \(\Phi(w)(x)=\Phi(xw)(1)\). The formulas commute with
maps of \(V,W\).

For \(f\) in the induction function model, let \(F(x)=f(x^{-1})\).
Then \(F(hx)=f(x^{-1}h^{-1})=hF(x)\), and
\(f(g^{-1}x^{-1})=F(xg)\) verifies equivariance. Inversion maps
the support coset \(tH\) to \(Ht^{-1}\), interchanging the
support conditions as claimed. Applying inversion again is the
inverse, and this transformation commutes with maps of \(V\).
\end{proof}

\section{Transitivity, double cosets, and character calculations}

\begin{theorem}[Induction in stages and Mackey decomposition]
\label{p103:ALG-07-C03-R03}
For \(H\le K\le G\), induction from \(H\) through \(K\) is
naturally isomorphic to induction directly from \(H\) to \(G\).
For arbitrary \(H,K\le G\), choose representatives \(s\) of
\(K\backslash G/H\), set \(L_s=K\cap sHs^{-1}\), and define
\({}^{s}V\) on the underlying space of \(V\) by
\(\ell\cdot v=\rho(s^{-1}\ell s)v\) for \(\ell\in L_s\).
Then there is a natural isomorphism of \(K\)-representations
\[
 \operatorname{Res}_K^G \operatorname{Ind}_H^G V
 \simeq
 \bigoplus_{s\in K\backslash G/H}\operatorname{Ind}_{L_s}^K({}^{s}V).
\]
The summands are independent, up to the given identification,
of the choices of representatives.
\end{theorem}

\begin{proof}
The multiplication map
\(g\otimes(k\otimes v)\mapsto gk\otimes v\) is balanced at both
tensor steps. Its inverse is \(g\otimes v\mapsto g\otimes(1\otimes v)\).
A relation \(gh\otimes v=g\otimes hv\) is preserved since
\(gh\otimes(1\otimes v)=g\otimes(h\otimes v)
=g\otimes(1\otimes hv)\). The other composite is the identity by
balancing over \(k[K]\), and the maps are \(G\)-equivariant.

For Mackey's formula, the group basis gives a direct sum of
\((k[K],k[H])\)-bimodules \(k[G]=\bigoplus_s k[KsH]\).
The tensor product over \(k[H]\) commutes with the direct sum.
On a fixed block the map from the proposed \(K\)-summand to
the block is
\[
 k[K]\otimes_{k[L_s]}{}^s V \longrightarrow
 k[KsH]\otimes_{k[H]}V,\qquad
 a\otimes v\longmapsto as\otimes v\quad(a\in K).
\]
It is balanced because \(a\ell s=as(s^{-1}\ell s)\). An inverse
sends \(as h\otimes v\) to \(a\otimes hv\), where \(a\in K,h\in H\).
If \(as h=a's h'\), then
\(\ell=a'^{-1}a=s h'h^{-1}s^{-1}\in L_s\). Consequently
\[
 a\otimes hv
   =a'\otimes(s^{-1}\ell s)hv=a'\otimes h'v .
\]
This proves independence of the expression \(as h\), and
the formula is also balanced on the right over \(k[H]\).
The two composites fix the indicated generators and both maps
are \(K\)-equivariant.

If \(s'=a_0s h_0\), then \(L_{s'}=a_0L_sa_0^{-1}\) and
the change of summand coordinates is
\(a\otimes v\mapsto aa_0\otimes h_0v\), from the \(s'\)-model to
the \(s\)-model. Under the two block maps both tensors have
image \(as'\otimes v\). Since both block maps are isomorphisms,
this proves well-definedness, invertibility, and representative
independence. All formulas commute with \(H\)-maps of \(V\).
\end{proof}

\begin{theorem}[Induced character and projection formula]
\label{p103:ALG-07-C03-S02}
For finite \(G\), \(H\le G\), and complex finite-dimensional \(V\),
\[
 \chi_{\operatorname{Ind}V}(g)=
 \frac1{|H|}
   \sum_{\substack{x\in G\\x^{-1}gx\in H}}\chi_V(x^{-1}gx).
\]
The value at \(1\) is \([G:H]\dim_{\mathbb C}V\), and the character
vanishes on elements not conjugate into \(H\).
For any field and any \(H\)-representation \(V\) and
\(G\)-representation \(W\), there is a natural projection isomorphism
\[
 \operatorname{Ind}_H^G(V\otimes_k\operatorname{Res}_H^G W)
      \simeq (\operatorname{Ind}_H^G V)\otimes_k W.
\]
\end{theorem}

\begin{proof}
In the coset direct sum for induction, the operator of \(g\)
permutes summands according to \(tH\mapsto gtH\).
A diagonal block occurs precisely when \(gtH=tH\), in which
case it is the operator of \(t^{-1}gt\) on \(V\). The trace
is the sum of these block traces. Replacing \(t\) by \(th\)
conjugates \(t^{-1}gt\) inside \(H\), so summing over all
elements of each fixed coset multiplies its contribution
by \(|H|\). This proves the character formula and the two
consequences. The same finite sum defines induction on class
functions and agrees with linear extension from characters.

For the isomorphism, map
\(x\otimes(v\otimes w)\) to \((x\otimes v)\otimes xw\).
Replacing \(x\) by \(xh\) on the source gives, on the target,
\((x\otimes hv)\otimes x(hw)\), which is the image of the
balanced expression \(x\otimes(hv\otimes hw)\).
The inverse sends \((x\otimes v)\otimes w\) to
\(x\otimes(v\otimes x^{-1}w)\). For \(xh\otimes v\),
balancing in the source of the first map sends the latter
expression to
\(x\otimes(hv\otimes x^{-1}w)\), as required.
The maps are inverse on group-basis pure tensors. Equivariance
follows from \((gx)w=g(xw)\) and the diagonal action,
and naturality follows from commutation with intertwiners.
\end{proof}

\section{Irreducibility and normal subgroups}

\begin{theorem}[Mackey's irreducibility criterion]
\label{p103:ALG-07-C03-S03}
Let \(G\) be finite, let \(H\le G\), and let \(V\ne0\) be a
finite-dimensional complex representation of \(H\).
Then \(\operatorname{Ind}_H^G V\) is irreducible if and only if
\(V\) is irreducible and, for each double coset \(HsH\ne H\),
the \(H\cap sHs^{-1}\)-representations
\(\operatorname{Res}_{H\cap sHs^{-1}}^H V\) and \({}^s V\)
have no common irreducible constituent.
If \(H\) is normal, the second condition is that the conjugate
representations \({}^s V\) for \(s\notin H\) are not isomorphic to \(V\).
\end{theorem}

\begin{proof}
Put \(X=\operatorname{Ind}_H^G V\). Frobenius reciprocity and
Mackey decomposition, with \(K=H\), give
\[
 \dim\operatorname{End}_G(X)=
 \sum_{s\in H\backslash G/H}
   \dim\operatorname{Hom}_H
       \bigl(V,\operatorname{Ind}_{L_s}^H({}^s V)\bigr).
\]
For the identity double coset the contribution is
\(\dim\operatorname{End}_H(V)\). Each other term equals
\(\dim\operatorname{Hom}_{L_s}
    (\operatorname{Res}_{L_s}^H V,{}^s V)\), using right
adjunction and the finite-index comparison, or equality
of opposite Hom dimensions and left reciprocity.
By complete reducibility and Schur's lemma this term vanishes
exactly when there is no common simple constituent.
The identity term is at least one, and it equals one exactly
when \(V\) is irreducible. Since \(X\ne0\) by the direct-sum model,
the same criterion \(\dim\operatorname{End}_G(X)=1\) detects
irreducibility of \(X\). For a normal \(H\) the intersections
are all \(H\), and for irreducible \(V\) each conjugate is also
irreducible. This gives the asserted simplification.
\end{proof}

\begin{theorem}[Clifford restriction and the inertia subgroup]
\label{p103:ALG-07-C03-S04}
Let \(A\trianglelefteq G\) be a normal subgroup of a finite group
and let \(V\) be an irreducible complex representation of \(G\).
The irreducible types in \(\operatorname{Res}_A^G V\) form one
\(G\)-conjugacy orbit and occur with equal multiplicity.
Choose one of their isotypic components \(U\) and set
\(I=\{g\in G:gU=U\}\), the inertia subgroup. Then \(A\le I\),
\(U\) is an irreducible \(I\)-representation, and the map
\[
 \operatorname{Ind}_I^G U\longrightarrow V,\qquad
 g\otimes u\longmapsto gu
\]
is an isomorphism. If \(I=G\), restriction to \(A\) is isotypic.
\end{theorem}

\begin{proof}
The restriction to \(A\) is semisimple. If a submodule \(W\)
has type \(\rho\) over \(A\), then the action on \(gW\) under
the identification \(w\mapsto gw\) is \(a\mapsto\rho(g^{-1}ag)\).
Normality ensures \(g^{-1}ag\in A\). Taking sums over all
submodules of a type shows that \(g\) sends an isotypic component
onto the corresponding conjugate component, using \(g^{-1}\)
for the reverse inclusion. A sum of components over a whole
orbit is \(G\)-stable. Irreducibility of \(V\) forces a single
orbit, and the transported direct sums have equal multiplicities.

The stabilizer \(I\) contains \(A\) because \(U\) is an \(A\)-submodule.
If \(U'\subset U\) were a nonzero proper \(I\)-submodule, its
translates by representatives of \(G/I\) would lie in distinct
isotypic components. Their direct sum is \(G\)-stable, since
the action sends each \(gU'\) onto the component \(g'U'\) for
the coset \(g'I\) of the product, and \(I\) preserves \(U'\).
The sum is nonzero and proper, a contradiction. Finally the
displayed map is balanced because \(gi\,u=g(iu)\), is
\(G\)-equivariant, and is an isomorphism on each summand
\(g\otimes U\to gU\) in the direct-sum model.
\end{proof}

\begin{example}
Let \(H=A_3=\langle c\rangle\) in \(S_3\), and let \(V\) be a
one-dimensional representation with \(c\) acting by a primitive
cube root of unity. A transposition conjugates \(c\) to \(c^{-1}\),
so \(V\) and its nontrivial conjugate are distinct.
The normal-subgroup criterion shows that \(\operatorname{Ind}V\)
is irreducible. Its character is \((2,0,-1)\). The transpositions
are not conjugate into \(H\), and the contribution on the
three-cycles is the sum of the two primitive cube roots.
\end{example}

\section{Virtual characters and rational Artin induction}

\begin{theorem}[The rational induction criterion]
\label{p103:ALG-07-C03-S05}
For a finite group \(G\), the ring \(R(G)\) of virtual complex
characters is the free abelian group on its irreducible
characters, with product induced by tensor product.
Let \(\mathcal H\) be a family of subgroups of \(G\).
Then the map
\[
 \bigoplus_{H\in\mathcal H}\bigl(\mathbb Q\otimes_{\mathbb Z}R(H)\bigr)
       \xrightarrow{\ \sum_H\operatorname{Ind}_H^G\ }
       \mathbb Q\otimes_{\mathbb Z}R(G)
\]
is surjective if and only if the conjugates of the groups
\(H\in\mathcal H\) cover \(G\). In particular, for the family of
cyclic subgroups, each ordinary character is a rational linear
combination of characters induced from cyclic subgroups.
Equivalently, some positive integer multiple of each virtual
character is an integer linear combination of such inductions.
\end{theorem}

\begin{proof}
By complete reducibility and character classification, the
direct-sum monoid of representations is the free commutative
monoid on irreducible types, and its group completion maps
injectively to class functions. This gives the asserted basis
for \(R(G)\). The character of a tensor product is the
pointwise product. Induction and restriction are additive on
direct sums and therefore give maps on the group completions.

After tensoring with \(\mathbb C\), the virtual character bases
identify the spaces with the full spaces of complex class
functions. By Frobenius reciprocity, the adjoint of the sum of
the induction maps is simultaneous restriction, using the sum
of the subgroup pairings on the direct sum. Finite-dimensional
linear algebra makes the sum of induction maps surjective
exactly when simultaneous restriction is injective. The
orthogonal complement of the image is the kernel of the adjoint.
A class function in that kernel vanishes on all the subgroups,
and hence on their conjugates. Conversely, if a conjugacy class
is omitted by the union, its nonzero indicator is in the kernel.
Thus injectivity is equivalent to the conjugate covering condition.

The matrix for the sum of inductions in the irreducible bases
has integer entries, namely multiplicities. Its rank over
\(\mathbb Q\) equals its rank over \(\mathbb C\), since both
are determined by the same nonzero minors. The surjectivity
criterion therefore descends to \(\mathbb Q\). Each element
lies in the cyclic subgroup it generates, giving the last case.
Clearing denominators gives the stated integral-multiple form.
\end{proof}
\chapter{Complexes, homotopies, and resolutions}
\label{p103:ch10}

The module proofs in this chapter work over an arbitrary unital ring
on either fixed side. For chain complexes the differential lowers
degree. For cochain complexes it raises degree. Terms outside a
specified range are zero. Our signs for shifts and tensor products
will also be used in Chapter~\ref{p103:ch11}.

\section{Homology and homotopy}

\begin{definition}[Complexes and homology]
\label{p103:ALG-10-C01-R01}
A chain complex \(C\) is a family of modules \(C_n\) with maps
\(d_n:C_n\to C_{n-1}\) satisfying \(d_{n-1}d_n=0\). Its cycles and
boundaries in degree \(n\) are \(Z_n(C)=\ker d_n\) and
\(B_n(C)=\operatorname{im}d_{n+1}\), and its homology is
\(H_n(C)=Z_n(C)/B_n(C)\). The corresponding definitions for
\(d^n:C^n\to C^{n+1}\) are \(Z^n(C)=\ker d^n\),
\(B^n(C)=\operatorname{im}d^{n-1}\), and \(H^n(C)=Z^n(C)/B^n(C)\).
A complex is acyclic if all these quotients vanish.
\end{definition}

A chain map \(f:C\to D\) consists of maps \(f_n\) commuting with the
differentials. This commutation carries cycles to cycles and
boundaries to boundaries, so \(H_n(f)[z]=[f_n(z)]\) is well defined.
This formula preserves the identity, composition, and addition.
The cochain construction uses the same formula in degree \(n\).
A map inducing isomorphisms in all degrees is a
\emph{quasi-isomorphism}. See \cite[\S1.1]{weibel} and
\cite[\S6.1]{rotman}.

\begin{theorem}[Homotopy invariance and contraction]
\label{p103:ALG-10-C01-R02}
If chain maps \(f,g:C\to D\) satisfy
\[
 f_n-g_n=d_{n+1}^D s_n+s_{n-1}d_n^C,\qquad s_n:C_n\to D_{n+1},
\]
they induce the same map on homology. A chain homotopy equivalence
therefore induces isomorphisms. The same statements hold for
cochain complexes, with \(s\) lowering degree. A complex is
contractible, meaning that its identity is homotopic to zero, if and
only if it is acyclic and the short exact sequences
\[
 0\longrightarrow Z_n(C)\longrightarrow C_n
 \xrightarrow{d_n} Z_{n-1}(C)\longrightarrow 0
\]
split in every degree.
\end{theorem}

\begin{proof}
On a cycle \(z\), the difference \(f_n(z)-g_n(z)=d_{n+1}^D s_n(z)\)
is a boundary. If \(gf\) and \(fg\) are homotopic to the respective
identities, the induced maps are inverse by functoriality.
For a contracting homotopy, every cycle \(z\) equals \(d s(z)\), and
\(s\) restricted to the cycles is a section of the displayed map
in the adjacent degree. Conversely, write \(C_n=Z_n\oplus Z_{n-1}\)
using the sections. In these coordinates \(d_n(z,w)=(w,0)\), and
\(s_n(z,w)=(0,z)\) gives \(d s+s d=1\).
Replacing \(C^n\) by \(C_{-n}\) gives the cochain statements.
\end{proof}

For example, place
\(0\to\mathbb Z/2\xrightarrow{1\mapsto 2}\mathbb Z/4
\to\mathbb Z/2\to0\) in consecutive chain degrees. It is acyclic
but fails to be contractible. A lift of \(1\) along the last
surjection has order four and hence cannot define a section.
For the general statements see \cite[\S1.4]{weibel} and
\cite[Theorem 6.14 and Proposition 6.15]{rotman}.

\section{The homology sequence and mapping cones}

\begin{theorem}[Long exact homology and cohomology sequences]
\label{p103:ALG-10-C01-R03}
A degreewise short exact sequence of complexes
\(0\to A\xrightarrow{i}B\xrightarrow{p}C\to0\) gives a natural long
exact homology sequence with connecting maps
\(\partial:H_n(C)\to H_{n-1}(A)\). In cohomology the connecting map
is \(H^n(C)\to H^{n+1}(A)\).
\end{theorem}

\begin{proof}
For a cycle \(c\in C_n\), choose \(b\in B_n\) with \(p b=c\).
Then \(p(db)=0\), so there is a unique \(a\in A_{n-1}\) with
\(i a=db\). Injectivity of \(i\) and \(d^2=0\) give \(da=0\).
Replacing \(b\) by \(b+i a'\) changes \(a\) by \(da'\).
Replacing \(c\) by \(c+dc'\) can be accompanied by replacing \(b\)
by \(b+db'\), for a lift \(b'\) of \(c'\), and leaves \(a\)
unchanged. Thus \(\partial[c]=[a]\) is well defined and additive
and linear, since sums and scalar multiples of lifts are lifts.

At \(H_n(B)\), suppose \(b\) is a cycle with \(p b=dc'\).
Lift \(c'\) to \(b'\), and then \(b-db'\) comes from a cycle in
\(A_n\). Conversely \(p i=0\). At \(H_n(C)\), if
\(\partial[c]=0\), write the associated \(a=da'\). Then \(b-i a'\)
is a cycle lifting \(c\). The converse follows by choosing a cycle
lift. At \(H_{n-1}(A)\), if \(i a=db\), the element \(p b\) is a
cycle and connects to \([a]\). Conversely a connecting class has
image a boundary by construction. These three chases repeat in
each degree. A commuting map of short exact sequences carries a lift
and its differential to the corresponding lift and differential in
the other sequence, giving naturality. Reindexing with
\(C_n=C^{-n}\) gives the cochain case and its stated degree.
\end{proof}

This is the complex application of the snake construction in
Theorem~\ref{p103:ALG-10-C01-R04}. The lift calculation specifies the
sign used here. See \cite[Theorem 1.3.1]{weibel} and
\cite[Theorems 6.10 and 6.13]{rotman}.

\begin{theorem}[Shift and cone criterion]
\label{p103:ALG-10-C01-S01}
The chain shift has \(C[1]_n=C_{n-1}\) and differential \(-d_C\).
For \(f:C\to D\) the mapping cone is
\[
 \operatorname{Cone}(f)_n=D_n\oplus C_{n-1},\qquad
 d(y,x)=(d_D y+f(x),-d_C x).
\]
There is a short exact sequence \(0\to D\to\operatorname{Cone}(f)
\to C[1]\to0\). The connecting map
\(H_n(C[1])\to H_{n-1}(D)\) identifies with \(H_{n-1}(f)\).
Consequently \(f\) is a quasi-isomorphism if and only if its cone
is acyclic. The cochain cone of \(u:Y\to X\) has
\(\operatorname{Cone}(u)^n=X^n\oplus Y^{n+1}\) with
\(d(x,y)=(d_Xx+u(y),-d_Yy)\) and satisfies the same criterion.
\end{theorem}

\begin{proof}
In the first component of \(d^2(y,x)\), the terms involving \(x\)
are \(d_Df(x)-f(d_Cx)=0\). The second component is zero by \(d_C^2=0\).
The injection and projection onto the indicated summands are chain
maps. Lift the class of a cycle \(x\in C_{n-1}\) to \((0,x)\).
Its differential is \((f(x),0)\), giving the connecting map.
Exactness at the two terms adjacent to \(H_n(\operatorname{Cone}(f))\)
then proves one implication, and if the cone is acyclic exactness
gives injectivity and surjectivity of \(H(f)\). The same computation
gives the cochain assertion using the shifted terms \(Y^{n+1}\).
\end{proof}

The cone of \(1_C\) is contractible, with homotopy
\((y,x)\mapsto(0,y)\). A quasi-isomorphism need not be a homotopy
equivalence. The augmentation
\((\mathbb Z\xrightarrow{2}\mathbb Z)\to(\mathbb Z/2)[0]\)
is a quasi-isomorphism, but every chain map in the reverse direction
is zero. Thus no reverse chain map can induce the inverse in
degree zero. Compare \cite[\S1.5]{weibel}.

\section{Total complexes and acyclic assembly}

\begin{theorem}[First-quadrant assembly]
\label{p103:ALG-10-C01-S02}
Suppose a bicomplex \(E_{p,q}\) has commuting horizontal and vertical
chain differentials \(d_h,d_v\), and is supported in the first
quadrant, allowing a fixed translation of that quadrant. The
direct-sum total complex is
\[
 \operatorname{Tot}(E)_n=\bigoplus_{p+q=n} E_{p,q},
 \qquad d|_{E_{p,q}}=d_h+(-1)^p d_v .
\]
If all columns are acyclic, its total complex is acyclic. The same
holds if all rows are acyclic. These assertions also hold for
commuting first-quadrant cochain differentials and the indicated
cochain totalization. If the given differentials anticommute,
their ordinary sum is the total differential instead.
\end{theorem}

\begin{proof}
The two cross terms cancel, since the horizontal differential
changes \(p\) by one. Suppose columns are acyclic, and let \(z\)
be a chain cycle on a total diagonal. At its largest horizontal
index \(p\), the equation \(dz=0\) implies \(d_vz_{p,q}=0\) because
there is no horizontal contribution from index \(p+1\).
Choose \(w\in E_{p,q+1}\) with
\((-1)^p d_vw=z_{p,q}\). Subtracting \(dw\) removes the component
at \(p\) and changes only the component at \(p-1\).
There is a lower bound for \(p\), so this procedure expresses the
cycle as a finite boundary sum. For cochains start at the smallest
\(p\) and lift vertically from degree \(q-1\). Subtracting the
total coboundary then changes only the next index \(p+1\).
On the fixed diagonal there is an upper bound for \(p\), and a
cycle at the lowest vertical degree is zero if it has no
predecessor. This proves termination in that case as well.
For rows, interchange the two indices and multiply the summand
\(E_{p,q}\) by \((-1)^{pq}\). This gives an isomorphism of total
complexes since an index shift by one changes that sign by
\((-1)^p\) or \((-1)^q\).
\end{proof}

The augmented form used below follows from the cone criterion.
For example, augment a chain bicomplex in the horizontal direction
to a column \(Y_q\) at \(p=-1\). Write \(X\) for the unaugmented
total and \(f:X\to Y\) for the map induced by augmentation. On the
augmented total in degree \(n-1\), the coordinates are
\(Y_n\oplus X_{n-1}\) and the differential is
\((y,x)\mapsto(-d_Yy+f(x),d_Xx)\). Shifting it by \([1]\)
and changing the sign on \(x\) identifies it with
\(\operatorname{Cone}(f)\). Thus exact augmented rows imply that
\(f\) is a quasi-isomorphism. In the cochain direction an
augmentation \(Y^q\to E_{0,q}\) from \(p=-1\) produces
\(Y^{n+1}\oplus X^n\) in total degree \(n\), with differential
\((y,x)\mapsto(-d_Yy,f(y)+d_Xx)\). Reordering the summands gives
\(\operatorname{Cone}(Y\to X)\). The column version follows by the
same signed interchange of indices. Compare
\cite[\S1.2 and Lemma 2.7.3]{weibel}.

\section{Resolutions and comparison}

\begin{theorem}[Free, projective, and injective resolutions]
\label{p103:ALG-10-C02-R01}
Every module \(A\) has a free and hence a projective resolution
\(\cdots\to P_1\to P_0\to A\to0\) and an injective resolution
\(0\to A\to I^0\to I^1\to\cdots\). In a resolution the augmented
sequence is exact. When taking the homology of the resolution as
a complex, the term \(A\) and the augmentation are deleted.
\end{theorem}

\begin{proof}
Send the basis of a free module on the underlying set of \(A\) to
the corresponding elements of \(A\). Repeat on its kernel and then
on each subsequent kernel. The differential is the new surjection
followed by the inclusion of the kernel. This identifies its image
with the previous kernel at each step.

For the dual construction use the enough-injectives argument of
Theorem~\ref{p103:ALG-03-C02-S03}. Recall its concrete form.
The group \(\mathbb Q/\mathbb Z\) is injective by divisibility and
Baer's criterion over \(\mathbb Z\). A nonzero element of an
abelian group is detected by a homomorphism to this group, first
on its cyclic subgroup, then by extending the homomorphism.
For left modules the coinduced module
\(J=\operatorname{Hom}_{\mathbb Z}(R,\mathbb Q/\mathbb Z)\) has
action \((r f)(s)=f(sr)\). Evaluation at \(1\) identifies
\(\operatorname{Hom}_R(A,J)\) with
\(\operatorname{Hom}_{\mathbb Z}(A,\mathbb Q/\mathbb Z)\).
The inverse sends \(\chi\) to \(a\mapsto(s\mapsto\chi(sa))\).
This gives extension along every monomorphism, so \(J\) is
injective. The product of copies of \(J\) indexed by those
homomorphisms is injective by extending each component, and the
product evaluation map embeds \(A\). Right modules use the
opposite ring. Now embed the cokernel of \(A\hookrightarrow I^0\)
into \(I^1\), and repeat on successive cokernels. The resulting
augmented cochain sequence is exact.
\end{proof}

See \cite[\S\S2.2--2.3]{weibel} and
\cite[\S6.1, Proposition 6.4]{rotman}.

\begin{theorem}[Comparison and homotopy uniqueness]
\label{p103:ALG-10-C02-R02}
A module map \(u:A\to B\) lifts to a chain map from any projective
resolution \(P\) of \(A\) to any resolution \(Q\) of \(B\).
Any two lifts are homotopic. Dually \(u\) extends to a map from
any right resolution of \(A\) to any injective right resolution
of \(B\), unique up to homotopy. Two projective resolutions of
one module are homotopy equivalent, as are two injective
resolutions. Applying an additive functor to the appropriate
side and then taking homology gives functorial maps independent
of these choices, up to canonical natural isomorphism.
\end{theorem}

\begin{proof}
Lift \(u\epsilon_P\) along \(Q_0\to B\) using projectivity of
\(P_0\). Treat the augmentation as \(d_0\) during the induction.
If the lift has been constructed through degree \(n-1\), the map
\(f_{n-1}d_n^P\) takes values in \(\ker d_{n-1}^Q\), by the
preceding chain equation. Exactness gives an epimorphism
\(Q_n\to\ker d_{n-1}^Q\). Projectivity lifts the map to \(f_n\).

For two lifts \(f,g\), set \(s_{-1}=0\). Suppose the homotopy
equations hold in lower degrees. The residual
\(r_n=f_n-g_n-s_{n-1}d_n^P\) lies in \(\ker d_n^Q\).
For \(n=0\) this follows from equality after augmentation.
For \(n>0\), substitute the equation for \(f_{n-1}-g_{n-1}\)
into \(d_n^Q r_n\). The only remaining term contains
\(d_{n-1}^P d_n^P=0\). Lift \(r_n\) along
\(Q_{n+1}\to\ker d_n^Q\) to obtain \(s_n\).
Reverse all arrows, replace kernels by cokernels, and use
injectivity instead of projectivity to obtain the cochain result.
This is the same induction in the opposite abelian category.

Lifts of an identity in both directions compose to lifts of the
identity, so their composites are homotopic to identity maps.
Additive functors preserve the homotopy equation, with degrees
reversed for a contravariant functor. The homotopy invariance
theorem then gives the comparison isomorphisms after taking
homology. A lift of a composite and the composite of lifts lift
the same map, which proves the functoriality and the compatibility
of the comparison isomorphisms.
\end{proof}

Compare \cite[Comparison Theorems 2.2.6 and 2.3.7]{weibel} and
\cite[Theorem 6.16]{rotman}.

\begin{theorem}[Horseshoe construction]
\label{p103:ALG-10-C02-S01}
Given \(0\to A'\to A\to A''\to0\) and projective resolutions
\(P'\) and \(P''\) of the outside terms, there is a projective
resolution \(P\) of \(A\) and a short exact sequence of resolutions
\(0\to P'\to P\to P''\to0\), split in every nonaugmented degree,
with \(P_n=P'_n\oplus P''_n\). There is a dual construction for
injective resolutions. A map between short exact sequences lifts
to a commuting diagram of the chosen horseshoe resolutions.
\end{theorem}

\begin{proof}
At stage \(n\) let \(0\to K'_n\xrightarrow{i}K_n\to K''_n\to0\)
be the exact sequence of current kernels, with \(K_0=A\) and
the analogous convention outside. The given resolutions
provide epimorphisms \(p':P'_n\to K'_n\) and \(p'':P''_n\to K''_n\).
Choose a lift \(\ell:P''_n\to K_n\) of \(p''\), and define
\(p:P_n\to K_n\) by \(p(x,y)=i p'(x)+\ell(y)\).
For \(k\in K_n\), first lift its image in \(K''_n\) to \(y\).
Then lift \(k-\ell(y)\) from \(K'_n\) to \(x\).
Thus \(p\) is surjective. Projection on the second coordinate
sends \(\ker p\) onto \(\ker p''\). If \(p'' y=0\), choose
\(x\) with \(i p' x=-\ell y\). The kernel of that projection is
\(\ker p'\) in the first summand. This gives the next short exact
sequence of kernels. The differential \(P_n\to P_{n-1}\) for
\(n>0\) is \(p\) followed by the inclusion \(K_n\subseteq P_{n-1}\).
Its image is the preceding kernel, giving exactness and \(d^2=0\).
The diagonal blocks are the prescribed differentials. Its
off-diagonal block is the first component of the chosen lift
\(\ell\) into \(K_n\subseteq P'_{n-1}\oplus P''_{n-1}\).
The lower-left block is zero.

For naturality, suppose the maps of syzygies through stage \(n\)
are known. Choose lifts \(\alpha:P'_n\to\widetilde P'_n\) and
\(\gamma:P''_n\to\widetilde P''_n\) of the outside maps.
If \(u:K_n\to\widetilde K_n\) is the middle map, then
\(u\ell-\widetilde\ell\gamma\) lands in
\(\widetilde K'_n\). Lift that difference to a map
\(t:P''_n\to\widetilde P'_n\). The middle lift is
\((x,y)\mapsto(\alpha x+t y,\gamma y)\).
It commutes with the three epimorphisms and hence induces the
next diagram of kernels. This starts at the module maps and
constructs a compatible diagram in all degrees. Different choices
induce the same maps on homology after an additive functor by
the comparison theorem. Reversing arrows gives the dual statement.
\end{proof}

This degreewise splitting is the reason an additive functor
preserves the short exact sequence of resolutions even when the
functor itself fails to be exact on all modules. The connecting
maps of the resulting complexes are therefore natural. See
\cite[Horseshoe Lemma 2.2.8]{weibel} and
\cite[Proposition 6.24 and Theorem 6.27]{rotman}.
\chapter{Ext, Tor, and coefficient theorems}
\label{p103:ch11}

All resolutions in this chapter are ordinary nonnegative resolutions.
For Ext the modules are both left \(R\)-modules, with an identical
right-module theory over the opposite ring. For Tor the first module
is right and the second is left. The groups are abelian groups in
general and naturally \(R\)-modules when \(R\) is commutative.
In the latter case all coefficient sequences below are \(R\)-linear.
The two-sided tensor conventions and the ordinary tensor--Hom
adjunction are those of Theorem~\ref{p103:ALG-03-C02-S01}.

\section{The bifunctors and their balanced descriptions}

\begin{definition}[Resolution definitions of Ext and Tor]
\label{p103:ALG-10-C03-R01}
Let \(n\geq0\). For left modules \(A,B\), choose a projective
resolution \(P\to A\) and put
\[
 \operatorname{Ext}_R^n(A,B)=
 H^n\!\bigl(\operatorname{Hom}_R(P_\bullet,B)\bigr),\qquad
 d^n(\phi)=\phi\,d_{n+1}^P .
\]
For a right module \(A\) and a left module \(B\), choose a projective
resolution \(P\to A\) of right modules and put
\[
 \operatorname{Tor}^R_n(A,B)=H_n(P_\bullet\otimes_R B),
 \qquad d_n=d_n^P\otimes 1 .
\]
The resolution complexes in these formulas have their augmentations
deleted.
\end{definition}

Here Ext is contravariant in its first variable and covariant in
the second. A map \(A'\to A\) is lifted to \(P'\to P\) and then
used by precomposition, whereas \(B\to B'\) acts by
postcomposition. Tor is covariant in both variables, using
tensoring of the lifted map or of the given map on \(B\).
Theorem~\ref{p103:ALG-10-C02-R02} shows that every lift gives the
same induced map and that compositions agree. In particular a
lift of the identity between two resolutions induces the canonical
natural comparison, with inverse induced by a lift in the other
direction. These arguments also prove additivity of the
bifunctors. See \cite[\S\S2.4--2.5 and 3.1--3.3]{weibel}.

\begin{theorem}[Balance]
\label{p103:ALG-10-C03-S01}
For the right-left pair in Tor, choose projective resolutions
\(P\to A\) and \(Q\to B\). Both maps
\[
 P_\bullet\otimes_R B
 \longleftarrow\operatorname{Tot}(P_\bullet\otimes_R Q_\bullet)
 \longrightarrow A\otimes_R Q_\bullet
\]
induced by augmentations are quasi-isomorphisms. For the same-side
pair in Ext, choose projective \(P\to A\) and injective \(B\to I\).
Both maps
\[
 \operatorname{Hom}_R(P_\bullet,B)
 \longrightarrow\operatorname{Tot}\operatorname{Hom}_R(P_\bullet,I^\bullet)
 \longleftarrow\operatorname{Hom}_R(A,I^\bullet)
\]
are quasi-isomorphisms. The resulting identifications are natural
in \(A\) and \(B\).
\end{theorem}

\begin{proof}
For Tor, in bidegree \((p,q)\) use \(P_p\otimes_R Q_q\),
with commuting differentials \(d_P\otimes 1\) and \(1\otimes d_Q\)
before inserting the sign \((-1)^p\). Each \(Q_q\) is flat, so
each row remains exact when the augmentation
\(P_\bullet\to A\) is included and tensored with \(Q_q\).
The horizontal augmented form of
Theorem~\ref{p103:ALG-10-C01-S02} proves the right-hand
quasi-isomorphism. Each \(P_p\) is flat, so the vertical
augmentation gives the left-hand one. Flatness of projectives
was established in Theorem~\ref{p103:ALG-03-C02-R04}.

For Ext the unsigned horizontal and vertical cochain maps are
precomposition with \(d_P\) and postcomposition with \(d_I\).
They commute, so use the same total sign rule. The augmented rows
are exact by injectivity of \(I^q\) and the augmented columns by
projectivity of \(P_p\). The cochain form of the assembly theorem
therefore proves both comparisons.
For module maps in either variable, lift to maps of the relevant
resolutions. The augmentation diagrams commute and hence so do
the induced isomorphisms in homology. Homotopy uniqueness gives
independence of the chosen lifts on each outer complex. These
outer comparisons are isomorphisms on homology, so the middle
comparison cannot introduce any further dependence. This proves
naturality and choice-independence.
\end{proof}

For a commutative ring the tensor interchange on a total complex
sends \(a_p\otimes b_q\) to \((-1)^{pq}b_q\otimes a_p\).
It is a chain isomorphism by the sign computation of
Theorem~\ref{p103:ALG-10-C01-S02}. Together with balance it gives
the usual Tor symmetry. The balance assertion itself requires
only the stated left-right convention. See
\cite[Theorems 2.7.2 and 2.7.6]{weibel} and
\cite[Theorem 6.32]{rotman}.

\begin{theorem}[Computation from flat resolutions]
\label{p103:ALG-10-C02-R03}
An ordinary nonnegative flat resolution \(F\to A\) of a right
module computes \(\operatorname{Tor}^R_n(A,B)\) as
\(H_n(F\otimes_R B)\). A flat resolution \(G\to B\) of a left
module computes it as \(H_n(A\otimes_R G)\).
The isomorphisms are natural with respect to maps of the
augmented resolutions. They include all \(n\geq0\).
\end{theorem}

\begin{proof}
Choose a projective resolution \(Q\to B\). In
\(F\otimes_R Q\), both augmented directions are exact for
the same reason as in the proof of balance: \(F_p\) and \(Q_q\)
are flat. The augmented total maps are quasi-isomorphisms.
This gives \(H_n(F\otimes_R B)\simeq H_n(A\otimes_R Q)\), and
the latter is the balanced Tor model. The construction for
\(G\to B\) interchanges the two variables. Naturality follows
by the same commuting augmentation diagrams. A projective
resolution of \(A\) admits a lift over the identity into \(F\)
by comparison, and this lift induces the identification just
constructed, since both comparisons with \(A\otimes_R Q\)
agree. Thus the formula also identifies results from different
flat resolutions even when a map between them is not chosen.
\end{proof}

This argument uses nonnegative augmented resolutions, which give
first-quadrant assembly. Arbitrary unbounded complexes of flat
modules require a different condition for such a computation.
Compare \cite[Lemma 3.2.8]{weibel} and
\cite[Theorem 7.5]{rotman}.

\section{Long exact sequences, vanishing, and dimension shifting}

\begin{theorem}[Degree zero, exactness, and vanishing]
\label{p103:ALG-10-C03-R02}
In degree zero the natural identifications are
\(\operatorname{Ext}^0_R(A,B)\simeq\operatorname{Hom}_R(A,B)\) and
\(\operatorname{Tor}^R_0(A,B)\simeq A\otimes_R B\).
A short exact sequence in either variable gives a natural long
exact sequence. Its arrows reverse the order of the three
coefficient modules in the first Ext variable. Positive Ext groups
vanish when the first variable is projective or the second
injective. Positive Tor groups vanish when either variable is flat
on its corresponding side. Moreover a right module \(A\) is flat
if and only if \(\operatorname{Tor}^R_1(A,-)\) vanishes on all left
modules, and dually for the other side.
\end{theorem}

\begin{proof}
A degree-zero cocycle in \(\operatorname{Hom}_R(P_\bullet,B)\)
kills \(\operatorname{im}d_1^P\), so it factors uniquely
through \(A\). There are no coboundaries in degree zero.
In the Tor calculation the cokernel of \(d_1^P\otimes1\)
is \(A\otimes_R B\) by right exactness. Both identifications
commute with coefficient maps by their quotient constructions.

For a short exact sequence \(0\to A'\to A\to A''\to0\),
apply the projective horseshoe. In Tor, tensor with the
fixed \(B\), preserving short exactness because the resulting
sequence is split in each degree before tensoring. For Ext in
the first variable apply \(\operatorname{Hom}_R(-,B)\).
Degreewise splitting again gives short exactness, now in the
order \(A'',A,A'\). For a short exact sequence in the second
Ext variable, use \(\operatorname{Hom}_R(P_\bullet,-)\).
Projectivity of every \(P_n\) makes this degreewise exact.
Equivalently, use an injective horseshoe in the balanced Ext
model. To compare these two choices, apply the balanced total
construction to the three injective resolutions and the same \(P\).
Its degreewise short exact sequence receives both model sequences
by the augmentations. Their quasi-isomorphisms commute with
connecting maps by naturality. For Tor in the second variable use the balanced
resolution of that variable and its projective horseshoe.
Theorem~\ref{p103:ALG-10-C01-R03} gives the four long exact
sequences. The comparison maps from the horseshoe construction,
or the fixed first resolution in the second Ext variable, give
naturality. Applying the same compatible comparison over the
identity shows that changing a horseshoe or its resolutions
preserves the connecting maps.

A projective or injective coefficient can be resolved by itself
in degree zero, in the appropriate balanced description. For Tor
the same argument uses the flat-resolution theorem. Conversely,
for each injection \(B'\hookrightarrow B\), apply the second
variable Tor sequence to its cokernel \(B''\). If
\(\operatorname{Tor}_1^R(A,B'')=0\), exactness makes
\(A\otimes_R B'\to A\otimes_R B\) injective. Together with
right exactness this is the flatness criterion, and the first
variable argument gives the dual assertion.
\end{proof}

In particular, the first-variable Ext sequence contains
\[
 \operatorname{Ext}^n_R(A'',B)\longrightarrow\operatorname{Ext}^n_R(A,B)
 \longrightarrow\operatorname{Ext}^n_R(A',B)
 \longrightarrow\operatorname{Ext}^{n+1}_R(A'',B).
\]
The second-variable sequence follows the original coefficient
order, also increasing degree at the connecting map. Both Tor
sequences follow the original coefficient order and lower degree
at the connecting map. See \cite[\S\S2.4--2.5 and 3.2]{weibel}
and \cite[\S\S6.2 and 7.1]{rotman}.

For a basic computation, let \(R\) be commutative and \(r\in R\)
a nonzerodivisor. The augmented two-term complex
\(0\to R\xrightarrow{r}R\to R/(r)\to0\) is a free resolution.
For every \(R\)-module \(B\) this gives
\[
 \operatorname{Tor}_1^R(R/(r),B)=\{b:rb=0\},\qquad
 \operatorname{Ext}^1_R(R/(r),B)=B/rB .
\]
The higher groups vanish and the degree-zero formulas respectively
give \(B/rB\) and \(\{b:rb=0\}\). If \(r\) is a zero divisor
the first arrow fails to be injective, so this particular
two-term resolution does not apply.

\begin{theorem}[Dimension shifting]
\label{p103:ALG-10-C02-R04}
For \(0\to K\to P\to A\to0\) with \(P\) projective, and \(n\geq1\),
\[
 \operatorname{Ext}^{n+1}_R(A,B)\simeq\operatorname{Ext}^n_R(K,B),
 \qquad
 \operatorname{Tor}_{n+1}^R(A,B)\simeq\operatorname{Tor}_n^R(K,B),
\]
where the module sides are chosen for the respective theories.
The Tor formula requires only flatness of \(P\) and has the
corresponding second-variable version. For \(0\to B\to I\to L\to0\)
with \(I\) injective,
\(\operatorname{Ext}^{n+1}_R(A,B)\simeq\operatorname{Ext}^n_R(A,L)\)
for \(n\geq1\).
\end{theorem}

\begin{proof}
The displayed identifications are the connecting maps or their
inverses in the respective long exact sequences. For example,
the groups on either side of
\(\operatorname{Ext}^n_R(K,B)\to\operatorname{Ext}^{n+1}_R(A,B)\)
in the sequence coming from \(P\) vanish by projectivity in
positive degree. Flatness gives the analogous Tor vanishings,
and injectivity gives those for the sequence through \(I\).
At the degree-zero edge the same sequences instead yield
\[
 \begin{aligned}
 \operatorname{Ext}^1_R(A,B)
   &\simeq \operatorname{coker}
   \bigl(\operatorname{Hom}_R(P,B)\to\operatorname{Hom}_R(K,B)\bigr),\\
 \operatorname{Tor}_1^R(A,B)
   &\simeq \ker\bigl(K\otimes_R B\to P\otimes_R B\bigr),\\
 \operatorname{Ext}^1_R(A,B)
   &\simeq \operatorname{coker}
   \bigl(\operatorname{Hom}_R(A,I)\to\operatorname{Hom}_R(A,L)\bigr).
 \end{aligned}
\]
The second-variable Tor edge has the corresponding tensor kernel.
\end{proof}

See \cite[Exercise 2.4.3 and \S2.5]{weibel} and
\cite[\S6.2]{rotman}.

\section{Extensions and their addition}

\begin{theorem}[Classification of module extensions]
\label{p103:ALG-10-C03-S02}
The equivalence classes of short exact sequences
\(0\to B\to E\to A\to0\), under maps fixing the outside terms,
are naturally identified with \(\operatorname{Ext}^1_R(A,B)\).
The class is the image of \(1_A\) under the second-variable
connecting map \(\operatorname{Hom}_R(A,A)\to
\operatorname{Ext}^1_R(A,B)\). The split extension represents
zero, and Baer addition corresponds to addition in Ext.
\end{theorem}

\begin{proof}
Fix \(0\to K\to P_0\xrightarrow{\epsilon}A\to0\) with \(P_0\)
projective. For the given extension, lift \(\epsilon\) to
\(l:P_0\to E\). Restricting \(l\) to \(K\) takes values in the
embedded \(B\), giving \(\phi:K\to B\). Another lift changes
\(\phi\) by the restriction of a map \(P_0\to B\). This is the
class in the first cokernel of
Theorem~\ref{p103:ALG-10-C02-R04}. In the resolution complex its
cocycle is \(\phi d_1^P\) with \(d_1^P\) viewed as a map to \(K\).
The cochain connecting formula applied to \(1_A\) lifts the
degree-zero representative \(\epsilon\) to \(l\) and takes its
differential, precisely this cocycle.

Conversely the pushout of \(K\hookrightarrow P_0\) along \(\phi\) is given by
\[
 E_\phi=(B\oplus P_0)/\{(\phi(k),-k):k\in K\}.
\]
The map from \(B\) is injective because a relation with zero
second component has \(k=0\). The map \([b,p]\mapsto\epsilon(p)\)
is onto, and its kernel consists of the classes \([b,k]\),
which equal \([b+\phi(k),0]\). Thus \(E_\phi\) is an extension,
and the lift \(p\mapsto[0,p]\) gives the original \(\phi\).
For an extension with lift \(l\), \([b,p]\mapsto b+l(p)\),
using the inclusion of \(B\), is a map fixing both ends.
It is an isomorphism by the short five lemma. If
\(\phi'=\phi+\psi|_K\), the map
\([b,p]\mapsto[b-\psi(p),p]\) identifies \(E_\phi\) with
\(E_{\phi'}\) and fixes the ends. This proves the bijection and
its independence of representatives. The pushout for \(\phi=0\)
is \(B\oplus A\), so precisely the split class corresponds to zero.

Pushing out along \(B\to B'\) composes \(\phi\) with that map.
Pulling back along \(A'\to A\) uses a lift \(P'_0\to P_0\) and
composes \(\phi\) with the induced kernel map. These are the
two Ext actions, proving naturality. For addition, take the
direct sum of two extensions, pull back along the diagonal
\(A\to A\oplus A\), and push out along addition \(B\oplus B\to B\).
The paired lifts give \((\phi,\phi')\) on \(K\) in the pullback,
and addition gives \(\phi+\phi'\) in the pushout. Hence the
Baer operation is exactly the additive group law.
\end{proof}

Compare \cite[Theorem 3.4.3 and Corollary 3.4.5]{weibel} and
\cite[Theorems 7.30--7.35]{rotman}.

\section{Universal coefficients}

\begin{lemma}[Submodules of free modules over a PID]
\label{p103:pid-submodule-arbitrary}
If \(R\) is a principal ideal domain, every submodule \(N\) of a
free \(R\)-module \(F\) is free, with no cardinality restriction on
a basis of \(F\).
\end{lemma}
\begin{proof}
Well-order a basis \((e_i)_{i\in I}\) of \(F\). For each \(i\),
let \(F_{\le i}\) be the span of basis elements of index at most
\(i\), and let \(J_i\) be the image of \(N\cap F_{\le i}\)
under the \(i\)-th coordinate map. This is an ideal of \(R\).
If \(J_i\ne0\), choose a nonzero generator \(a_i\) of \(J_i\)
and a vector \(v_i\in N\cap F_{\le i}\) with \(i\)-th coordinate
\(a_i\). No vector is needed for \(J_i=0\).

In a nontrivial finite relation among the \(v_i\), take the
greatest index \(i\) with nonzero coefficient \(r_i\).
The \(i\)-th coordinate of the sum is \(r_i a_i\), which is
nonzero in the domain \(R\). Thus the chosen vectors are
linearly independent.

For generation, let \(x\in N\) be nonzero and take the largest
index \(i\) of its finite support. Its \(i\)-th coordinate lies
in \(J_i\) and is nonzero, so it is a multiple \(r a_i\).
The vector \(x-rv_i\) is in \(N\), and either vanishes or has
strictly smaller largest support index. By well-founded induction
on \(i\) the remainder is a finite linear combination of the
chosen vectors. Thus so is \(x\). For \(N=0\) the chosen family is
empty and is a basis of \(N\).
\end{proof}

\begin{theorem}[Universal coefficient sequences]
\label{p103:ALG-10-C03-R03}
Let \(C\) be a nonnegative chain complex.
For homology, assume that \(C\) is a complex of right \(R\)-modules
and that each \(C_n\) and \(B_n(C)\) is flat. For every left module \(M\)
and \(n\geq0\) there is a natural short exact sequence
\[
 \begin{split}
 0&\longrightarrow H_n(C)\otimes_R M
 \longrightarrow H_n(C\otimes_R M)\\
 &\longrightarrow\operatorname{Tor}_1^R(H_{n-1}(C),M)
 \longrightarrow0 .
 \end{split}
\]
For cohomology, assume that \(C\) is a complex of left \(R\)-modules
and that each \(C_n\) and \(B_n(C)\) is projective.
For every left module \(M\) there is a natural short exact sequence
\[
 \begin{split}
 0&\longrightarrow\operatorname{Ext}_R^1(H_{n-1}(C),M)
 \longrightarrow H^n\!\bigl(\operatorname{Hom}_R(C,M)\bigr)\\
 &\longrightarrow\operatorname{Hom}_R(H_n(C),M)
 \longrightarrow0 .
 \end{split}
\]
Here \(H_{-1}(C)=0\). The cohomology sequence splits after
choosing sections of the boundary projections. The homology
sequence splits under the additional condition that the boundary
projections split, in particular if all boundaries are projective.
Thus for a free complex over a PID both sequences split
noncanonically.
\end{theorem}

\begin{proof}
Write \(Z_n=Z_n(C)\), \(B_n=B_n(C)\), \(B_{-1}=0\), and let
\(j_n:B_n\to Z_n\) be the inclusion. The two basic exact
sequences are
\[
 0\to Z_n\to C_n\xrightarrow{d_n}B_{n-1}\to0,
 \qquad
 0\to B_n\xrightarrow{j_n}Z_n\to H_n(C)\to0 .
\]
In the homology case \(Z_n\) is flat. Indeed the first short
exact sequence puts \(\operatorname{Tor}_1^R(Z_n,M)\) between
\(\operatorname{Tor}_2^R(B_{n-1},M)\) and
\(\operatorname{Tor}_1^R(C_n,M)\), both zero for every \(M\).
The flatness criterion in
Theorem~\ref{p103:ALG-10-C03-R02} applies.
Write \(Z\) for the complex with terms \(Z_n\) and zero
differential, and \(L\) for the complex with terms \(B_{n-1}\)
and zero differential. The first short exact sequences assemble to
\(0\to Z\to C\to L\to0\), which remains degreewise exact after
tensoring with \(M\) since \(\operatorname{Tor}_1^R(B_{n-1},M)=0\).
In its homology sequence the connecting map is
\[
 B_{n-1}\otimes_R M\longrightarrow Z_{n-1}\otimes_R M,
 \qquad j_{n-1}\otimes1 .
\]
Indeed \(b\otimes m\) lifts to \(c\otimes m\) with \(d_C c=b\),
and its differential is \(j_{n-1}(b)\otimes m\) in the cycles.
Such tensors generate the source. Exactness now places
\(H_n(C\otimes_R M)\) between
\(\operatorname{coker}(j_n\otimes1)\) and
\(\ker(j_{n-1}\otimes1)\) in a short exact sequence.
The second basic sequence and flat resolution computation
identify these terms with the claimed tensor and Tor groups.
The first arrow is \([z]\otimes m\mapsto[z\otimes m]\).
Naturality follows from the short exact sequences of cycles and
boundaries and the natural connecting maps.

In the cohomology case \(C_n\to B_{n-1}\) splits by projectivity,
so each \(Z_n\) is projective. Applying \(\operatorname{Hom}_R(-,M)\)
to \(0\to Z\to C\to L\to0\) gives a short exact sequence of
cochain complexes, now in the order \(L,C,Z\).
The connecting map from \(\operatorname{Hom}_R(Z_n,M)\) to
\(\operatorname{Hom}_R(B_n,M)\) is \(j_n^*\). To see this,
extend the functional to \(C_n\) using a section, and
precompose with \(d_{n+1}\). The result factors through
\(C_{n+1}\twoheadrightarrow B_n\), and the induced functional
is exactly the restriction to \(B_n\). The long exact sequence
now identifies the outer terms as the cokernel of \(j_{n-1}^*\)
and the kernel of \(j_n^*\), respectively. The projective
presentation \(B_n\to Z_n\to H_n(C)\) identifies these with Ext
and Hom in the indicated degrees. The right arrow sends a
cocycle to its restriction to \(Z_n\), descended to \(H_n(C)\).
All the identifications are natural.

For the splitting in homology, choose a decomposition
\(C_n=Z_n\oplus B_{n-1}\) from a section of \(C_n\to B_{n-1}\).
In these coordinates the differential sends
\((z,b)\) to \((j_{n-1}(b),0)\). In the tensor complex the
cycles in degree \(n\) are the direct sum of \(Z_n\otimes_R M\)
and \(\ker(j_{n-1}\otimes1)\). All the boundaries are in the
first summand and equal the image of \(j_n\otimes1\).
Taking the quotient exhibits a splitting of the displayed
sequence. In cohomology the same decomposition makes
a cocycle in degree \(n\) a pair of functionals on
\(Z_n,B_{n-1}\), of which the first kills \(B_n\).
Coboundaries are zero on \(Z_n\) and give exactly
\(\operatorname{im}(j_{n-1}^*)\) on \(B_{n-1}\).
Thus restriction on the first summand has a section by taking
zero on the second summand. Both splittings depend on the
chosen boundary sections. Over a PID,
Lemma~\ref{p103:pid-submodule-arbitrary} gives free boundary
and cycle modules whenever the chain modules are free, so the
required projectivity and flatness follow at arbitrary rank.
\end{proof}

The boundary hypotheses are essential to these general formulas.
For example, set \(R=k[\epsilon]/(\epsilon^2)\) over a field \(k\),
take \(C_n=R\) for \(n\geq0\), and take every positive-degree
differential to be multiplication by \(\epsilon\).
Then \(H_n(C)=0\) for \(n>0\) because the image and kernel of
multiplication by \(\epsilon\) coincide. With
\(M=R/(\epsilon)\), every differential of \(C\otimes_R M\)
is zero, so \(H_n(C\otimes_R M)=k\) in every nonnegative degree.
For \(n\geq2\) the proposed outer homology terms would both
vanish. The same example gives a zero cochain differential in
\(\operatorname{Hom}_R(C,M)\). Compare
\cite[Theorems 3.6.1, 3.6.2, and 3.6.5]{weibel} and
\cite[Theorems 7.55, 7.59, and Corollary 7.56]{rotman}.

\section{K\"unneth for algebraic complexes}

\begin{theorem}[K\"unneth short exact sequence]
\label{p103:ALG-10-C03-R04}
Let \(C,D\) be nonnegative chain complexes of right and left
\(R\)-modules, respectively. Suppose that all \(C_p\) and all
\(B_p(C)\) are flat. Give the tensor complex the direct-sum total
differential
\[
 d(c\otimes x)=d_Cc\otimes x+(-1)^p c\otimes d_D x,\qquad c\in C_p .
\]
There is a natural short exact sequence
\[
 \begin{split}
 0&\longrightarrow\bigoplus_{p+q=n}H_p(C)\otimes_R H_q(D)
 \longrightarrow H_n(C\otimes_R D)\\
 &\longrightarrow\bigoplus_{p+q=n-1}
   \operatorname{Tor}_1^R(H_p(C),H_q(D))\longrightarrow0 .
 \end{split}
\]
The last sum is zero when \(n=0\). Over a PID the sequence
splits noncanonically if one of the complexes is free.
In particular this includes the case when both are free.
\end{theorem}

\begin{proof}
Use \(Z,L\) from the preceding proof for \(C\).
Tensor the sequence \(0\to Z\to C\to L\to0\) with \(D\)
and take direct-sum total complexes. For each pair of degrees
it is exact, since the right module \(L_p=B_{p-1}(C)\)
is flat, and it remains exact after summing the finitely many
components on each diagonal. For a flat module \(U\),
the natural map \(U\otimes_R H_q(D)\to H_q(U\otimes_R D)\)
is an isomorphism. In fact tensoring preserves the cycle
kernel, the boundary image and the quotient that defines
homology. The sign on a fixed vertical column does not change
cycles or boundaries. The homology of the first and last
complexes is consequently the direct sum, in degree \(n\),
of \(Z_p(C)\otimes_R H_q(D)\) and
\(B_{p-1}(C)\otimes_R H_q(D)\), respectively.

Under these identifications the connecting homomorphism
maps the \((p,q)\) summand of the last term to the
\((p-1,q)\) summand of the first by \(j_{p-1}\otimes1\).
For an elementary class \(b\otimes[x]\) choose a lift
\(c\) of \(b\) and a cycle representative \(x\).
The vertical term of \(d(c\otimes x)\) is zero and the
horizontal term is \(j_{p-1}(b)\otimes x\), proving the
formula on generators. Its cokernel from the next total
degree is the stated direct sum of tensor products, and
its kernel in degree \(n\) is the stated direct sum of
Tor groups. This follows componentwise from the two-term
flat resolution \(0\to B_p(C)\to Z_p(C)\to H_p(C)\to0\).
The long exact sequence yields the desired short exact
sequence. The first map on elementary classes is
\([c]\otimes[x]\mapsto[c\otimes x]\). The constructions
commute with maps of both complexes, giving naturality.

To split over a PID, first suppose both \(C\) and \(D\) are
free. Lemma~\ref{p103:pid-submodule-arbitrary} gives freeness of
their boundary and cycle modules. Choose sections of all
projections onto boundaries.
Using \(C_p=Z_p(C)\oplus B_{p-1}(C)\) and the analogous
decomposition for \(D\), split each complex as a direct sum
of two-term subcomplexes. The summand indexed by \(p\)
for \(C\) is \(B_p(C)\to Z_p(C)\) in degrees \(p+1,p\),
where the differential is \(j_p\). Similarly for \(D\)
with index \(q\). These two-term complexes are shifted
free resolutions of the respective homology modules.
For a fixed pair of summands, its tensor total homology
is the tensor product in degree \(p+q\), Tor in degree
\(p+q+1\), and zero otherwise. Indeed the signed total
comparison in the balance theorem applies to the two
unshifted resolutions. Upon shifting, the fixed extra
sign \((-1)^p\) on the second differential can be
removed by multiplying its upper term by \((-1)^p\).
No higher Tor is present since either resolution has
length at most one. Taking the direct sum of all pairs,
the map from the first term of the short exact sequence
is the isomorphism onto the summands with \(p+q=n\),
by its cycle formula and the tensor cokernel in those
summands. The remaining summands, with \(p+q=n-1\),
therefore map isomorphically onto the last term and
give a section. This constructs the splitting.

Suppose now that \(C\) is free and \(D\) is arbitrary.
Construct a degreewise surjective quasi-isomorphism
\(f:Q\to D\) with \(Q\) nonnegative and free.
Take a free surjection \(Q_0\to D_0\). Having constructed
lower degrees, let \(Q_n\) be free on the underlying
set of the module
\[
 \{(z,y)\in Z_{n-1}(Q)\times D_n:
       f_{n-1}(z)=d_D y\},\qquad n>0,
\]
and let \(d_Q,f_n\) be the two projections after the
free surjection. This gives \(d_Q^2=0\) and the chain
map identity. Every cycle \(y\) lifts to a cycle because
\((0,y)\) is in the displayed module. Consequently
every \(y\in D_n\) lifts, by first lifting its boundary
to a cycle in \(Q_{n-1}\). A cycle \(z\) in the kernel
of \(f_{n-1}\) is the boundary of a lift of \((z,0)\)
which also lies in \(\ker f\). Thus the kernel complex
is acyclic, and the long exact sequence proves the
quasi-isomorphism.

Since every \(C_p\) is flat, \(C\otimes_R Q\to C\otimes_R D\)
is degreewise surjective with kernel the total complex
of \(C_p\otimes_R(\ker f)_q\). Every column of that
bicomplex is acyclic by flatness. The first-quadrant
assembly theorem makes the total kernel acyclic, so the
tensor map is a quasi-isomorphism as well. Naturality
identifies the entire coefficient sequence with the
one for \(C,Q\), whose splitting was constructed above.
If only \(D\) is free, replace \(C\) by a free complex
by the same construction. Tensoring the kernel complex
with the flat terms of \(D\) now gives acyclic rows,
so the same comparison and naturality prove splitting
in the stated orientation. The signed interchange also
gives the version of the theorem with the flatness
hypotheses imposed on the second factor instead.
\end{proof}

This argument is algebraic and applies to tensor complexes.
Using it for topological products calls for a comparison between
the product chain complex and the tensor of the two chain
complexes. See \cite[Theorem 3.6.3]{weibel}.
\chapter[Flat base change and a local criterion]
{Flat base change and a local criterion}
\label{p103:ch12}

In this chapter all rings are commutative. We first give a
comparison that is independent of any local finiteness
hypothesis for Tor, and then use the finite module and
completion techniques for a relative local criterion.

\section{Derived groups under flat extension of scalars}

\begin{theorem}[Flat base change and localization]
\label{p103:ALG-10-C03-S03}
Let \(R\to S\) be a flat map of rings, \(A,B\) be \(R\)-modules,
and \(C\) be an \(S\)-module. For \(n\ge0\) there are natural
isomorphisms
\[
 \begin{aligned}
 \operatorname{Tor}_n^S(S\otimes_R A,C)
  &\simeq \operatorname{Tor}_n^R(A,\operatorname{Res} C),\\
 S\otimes_R\operatorname{Tor}_n^R(A,B)
  &\simeq \operatorname{Tor}_n^S(S\otimes_R A,S\otimes_R B).
 \end{aligned}
\]
The first group on the right has its \(S\)-module action through \(C\).
If \(R\) is Noetherian and \(A\) is finitely generated, there is
also the natural isomorphism
\[
 S\otimes_R\operatorname{Ext}_R^n(A,B)
 \simeq \operatorname{Ext}_S^n(S\otimes_R A,S\otimes_R B).
\]
In particular, if \(T\subset R\) is a multiplicative set, the
latter formula holds with \(S=T^{-1}R\) and
\(S\otimes_R(-)=T^{-1}(-)\).
These Ext and Tor groups vanish if and only if their localizations
at all prime ideals vanish, or equivalently at all maximal ideals
(the indicated finiteness hypothesis applies to the Ext comparison).
\end{theorem}
\begin{proof}
Let \(P_\bullet\to A\) be a projective resolution. Flatness of
\(S\) gives an exact augmented complex \(S\otimes_R P_\bullet\to
S\otimes_R A\). Each term is projective because extension of
scalars carries a split summand of a free module to a split
summand of a free module. The first comparison of complexes
sends
\[
 (s\otimes p)\otimes c\longmapsto p\otimes sc,
 \qquad p\otimes c\longmapsto(1\otimes p)\otimes c
\]
in the two directions. The second identifies
\(S\otimes_R(P_\bullet\otimes_R B)\) with
\((S\otimes_R P_\bullet)\otimes_S(S\otimes_R B)\), by associativity
and multiplication of the two \(S\)-scalars. An exact functor
commutes with homology because it preserves kernels, images,
and their quotient. This gives the second comparison and
explains the \(S\)-actions.

For Ext, take a finite free cover of \(A\) and then a finite free
cover of each successive kernel. All kernels are finitely generated
by Noetherianity, so every \(P_i\) is finite free. The natural map
\[
 S\otimes_R\operatorname{Hom}_R(P_i,B)
 \longrightarrow
 \operatorname{Hom}_S(S\otimes_RP_i,S\otimes_RB)
\]
sends \(s\otimes f\) to the map \(t\otimes p\mapsto st\otimes f(p)\).
For a finite free \(P_i\), both sides are the same finite power
of \(S\otimes_R B\) and the map is the coordinate identity.
It is also an isomorphism for finite projectives by taking
summands. Naturality in \(P_i\) makes these maps a comparison of
cochain complexes. Flatness again allows cohomology to pass
through \(S\otimes_R-\). All comparisons commute with maps on
the other variable and with lifts of maps in \(A\). Homotopy
uniqueness of lifts then proves naturality independently of
the chosen resolution.

Localization is flat by the fraction exactness argument of
Chapter~\ref{p103:ch04}. For any nonzero element \(x\) of an
\(R\)-module choose a maximal ideal containing its proper
annihilator. Its image in the corresponding localization is
nonzero, since any denominator killing \(x\) would belong to
that maximal ideal. This proves the vanishing tests.
\end{proof}
For the base-change and localization formulas see
\cite[Propositions 3.2.9 and 3.3.10 and Corollary 3.3.11]{weibel}.
Finite generation of the resolution terms matters for Hom.
Over a general ring the corresponding expression for an
infinite free module involves products rather than finite powers.

\section{The relative local flatness criterion}

\begin{theorem}[Local flatness via the residue field]
\label{p103:ALG-05-C04-S03}
Let \((R,\mathfrak m)\to(S,\mathfrak n)\) be a local
homomorphism of Noetherian local rings and \(M\) a finitely
generated \(S\)-module. Then \(M\) is flat over \(R\) if and
only if
\(\operatorname{Tor}_1^R(R/\mathfrak m,M)=0\).
The \(S\)-module \(M\) need not be finitely generated over \(R\).
\end{theorem}
\begin{proof}
Flatness gives the Tor vanishing by the flat computation in
Chapter~\ref{p103:ch11}. Conversely assume the vanishing.
Every simple \(R\)-module is isomorphic to \(R/\mathfrak m\),
so the long exact Tor sequence and a composition series give
\(\operatorname{Tor}_1^R(F,M)=0\) for every finite-length
\(R\)-module \(F\). In particular this applies to
\(R/(I+\mathfrak m^j)\) for every ideal \(I\) and \(j\ge1\), since
\(R/\mathfrak m^j\) has a finite filtration with layers
\(\mathfrak m^i/\mathfrak m^{i+1}\) finite-dimensional over the
residue field.

To prove the ideal test for flatness, fix a finitely generated
ideal \(I\) and an element \(z\in I\otimes_R M\) mapping to zero
in \(R\otimes_R M=M\). For each \(j\) the exact sequence
\[
 0\longrightarrow I/(I\cap\mathfrak m^j)
 \longrightarrow R/\mathfrak m^j
 \longrightarrow R/(I+\mathfrak m^j)\longrightarrow0
\]
stays injective on the first map after tensoring, by the
finite-length vanishing just proved. The image of \(z\) in
\((I/(I\cap\mathfrak m^j))\otimes_R M\) is therefore zero.
Right exactness puts \(z\) in the image of
\((I\cap\mathfrak m^j)\otimes_R M\to I\otimes_R M\).

Apply Artin--Rees to \(I\subseteq R\) for the ideal \(\mathfrak m\).
For some fixed \(c\), \(I\cap\mathfrak m^j\subseteq
\mathfrak m^{j-c}I\) for all \(j\ge c\). Balance of the tensor
then gives \(z\in \mathfrak m^{j-c}(I\otimes_R M)\) for those
\(j\). As an \(S\)-module \(I\otimes_R M\) is finite. An
\(R\)-generating set of \(I\) and an \(S\)-generating set of
\(M\) give generators by simple tensors. The map is local,
so \(\mathfrak m S\subseteq\mathfrak n\), and \(z\) lies in
every \(\mathfrak n^k(I\otimes_R M)\). The Krull intersection
theorem over \(S\) forces \(z=0\). The finite ideal test
\cref{p103:ALG-03-C02-R04} now proves flatness.
\end{proof}
This proof uses only the ordinary \(\operatorname{Tor}\)
long exact sequence, the Artin--Rees estimate over \(R\),
and separatedness for finite modules over \(S\). Compare the
local criterion and its finite module formulation in
\cite[Theorem 6.8 and Corollary 6.9]{eisenbud}.

\backmatter
\begin{thebibliography}{99}
\bibitem{lam} T. Y. Lam, \href{https://doi.org/10.1007/978-1-4612-0525-8}{\emph{Lectures on Modules and Rings}}, Graduate Texts in Mathematics 189, Springer, 1999.
\bibitem{dummit} David S. Dummit and Richard M. Foote, \emph{Abstract Algebra}, 3rd ed., Wiley, 2004.
\bibitem{eisenbud} David Eisenbud, \href{https://doi.org/10.1007/978-1-4612-5350-1}{\emph{Commutative Algebra with a View Toward Algebraic Geometry}}, Graduate Texts in Mathematics 150, Springer, 1995.
\bibitem{serre} Jean-Pierre Serre, \href{https://doi.org/10.1007/978-1-4684-9458-7}{\emph{Linear Representations of Finite Groups}}, translated by Leonard L. Scott, Graduate Texts in Mathematics 42, Springer, 1977.
\bibitem{weibel} Charles A. Weibel, \emph{An Introduction to Homological Algebra}, Cambridge University Press, 1994, first paperback printing 1995. \href{https://sites.math.rutgers.edu/~weibel/Hbook-corrections.html}{Author's corrections}.
\bibitem{rotman} Joseph J. Rotman, \href{https://doi.org/10.1007/b98977}{\emph{An Introduction to Homological Algebra}}, 2nd ed., Springer, 2009.
\bibitem{d07} MathGenome, \emph{Unified Abstract Algebra}, D07, MathGenome v0.8.0.
\bibitem{d05} MathGenome, \emph{Unified Linear Algebra}, D05, MathGenome v0.8.0.
\end{thebibliography}
\end{document}
